% !TeX root = ../大学生数学竞赛概率论与数理统计试题选编与专题训练.tex
% 由主文档通过 \input 载入；请编译主文档。
\setcounter{chapter}{4} % 保留原书章号。
\chapter{课程衔接训练：解答}

\clearpage\section*{事件、计数与几何概率：详解}\addcontentsline{toc}{section}{事件、计数与几何概率：详解}

\subsection*{训练题~\ref{q:Y1.1} 事件分解与德摩根律}\label{a:Y1.1}{\small 题目见第~\pageref{q:Y1.1} 页；来源：闫一博公开解答，原题 1.1。}
\begin{solve}这里用乘积表示交、上横线表示补，只有互不相容事件的并才使用加号。对任意样本点 $\omega$，$\omega\in A\cup B$ 等价于“$\omega\in A$，或者 $\omega\notin A$ 且 $\omega\in B$”，所以 $A\cup B=A+\overline AB$。按是否属于 $B$ 划分 $A$，得到 $A=AB+A\overline B$。对于任意有限或可数指标集 $J$，\[\omega\notin\bigcup_{j\in J}A_j\iff\forall j\in J,\ \omega\notin A_j,\]\[\omega\notin\bigcap_{j\in J}A_j\iff\exists j\in J,\ \omega\notin A_j.\]分别把右侧翻译为补集的交与并，即得两条德摩根律。此论证同时适用于题中的有限与可数版本。\end{solve}

\subsection*{训练题~\ref{q:Y1.2} 用集合语言表示事件}\label{a:Y1.2}{\small 题目见第~\pageref{q:Y1.2} 页；来源：闫一博公开解答，原题 1.2。}
\begin{solve}
（a）$A\overline{B}\:{\overline{C}}$；（b）$\overline{A}BC$；（c）$\overline{A}\:\overline{B}\:\overline{C}$；（d）$A\overline{B}\:\overline{C}+\overline{A}B\overline{C}+\overline{A}\:\overline{B}C$；（e）$A\cup B\cup C$.
\end{solve}

\subsection*{训练题~\ref{q:Y1.3} 有放回抽球与样本空间}\label{a:Y1.3}{\small 题目见第~\pageref{q:Y1.3} 页；来源：闫一博公开解答，原题 1.3。}
\begin{solve}
（a）
\[\Omega  = \{ \omega :\omega  = ({a_1},{a_2},{a_3}),{a_i} \in \{\text{红、黄、蓝、白} \} ,i = 1,2,3\} \]

（b）记取到三个球颜色互不相同为事件$A$，则
\[P(A) = \frac{{C_4^3}\cdot 3!}{{{4^3}}} = \frac{{24}}{{64}} = \frac{3}{8}\]
\end{solve}

\subsection*{训练题~\ref{q:Y1.4} 将事件列化为互斥事件列}\label{a:Y1.4}{\small 题目见第~\pageref{q:Y1.4} 页；来源：闫一博公开解答，原题 1.4。}
\begin{solve}取 $B_1=A_1$，对 $j\ge2$ 定义 $B_j=A_j\setminus\bigcup_{i<j}A_i$。若 $i<j$，$B_j$ 不含 $A_i$ 中的任何点，而 $B_i\subset A_i$，所以 $B_i\cap B_j=\varnothing$。若某点属于前 $n$ 个 $A_i$ 的并，取它第一次进入的最小指标 $j$，它便属于 $B_j$；反过来 $B_j\subset A_j$，故两个有限并相等。对无限并同样使用最早进入指标，或对有限并取递增极限，得到可数并相等。\end{solve}

\subsection*{训练题~\ref{q:Y1.5} 两阶段抽样的对称性}\label{a:Y1.5}{\small 题目见第~\pageref{q:Y1.5} 页；来源：闫一博公开解答，原题 1.5。}
\begin{solve}
（a）记任取$2$件且无次品为事件$A$，则
\[P(A) = \frac{{C_{97}^2}}{{C_{100}^2}} \]
记任取$2$件至少有一件次品为事件$B$，则
\[
P(B)=1-P(A)=\frac{147}{2475}
\]
（b）可视作从$100$件产品中取出$10$件不拿出，再从$10$件产品中抽$2$件，故概率与（a）相同.
\end{solve}

\subsection*{训练题~\ref{q:Y1.7} 无放回试钥匙}\label{a:Y1.7}{\small 题目见第~\pageref{q:Y1.7} 页；来源：闫一博公开解答，原题 1.7。}
\begin{solve}
记第$i$次打开房门为事件$A_i$，则

（a）记第三次打开房门为事件$B$
\[P(B)=P({\overline A _1}{\overline A _2}{A_3}) = \frac{4}{5} \times \frac{3}{4} \times \frac{1}{3} = \frac{1}{5}\]

（b）记三次内打开房门为事件$C$
\[\begin{aligned}P(C) &= 1 - P({\overline A _1}{\overline A _2}{\overline A _3}) \\&= 1 - \frac{4}{5} \times \frac{3}{4} \times \frac{2}{3} \\&= \frac{3}{5}\end{aligned}\]

（c）若有两把钥匙可开门，记三次内打开房门为事件$D$，则
\[\begin{aligned}P(D) &= 1 - P({\overline A _1}{\overline A _2}{\overline A _3}) \\&= 1 - \frac{3}{5} \times \frac{2}{4} \times \frac{1}{3} \\&= \frac{9}{{10}}\end{aligned}\]
\end{solve}

\subsection*{训练题~\ref{q:Y1.8} 生日相同的概率}\label{a:Y1.8}{\small 题目见第~\pageref{q:Y1.8} 页；来源：闫一博公开解答，原题 1.8。}
\begin{solve}
记$n$个人生日互不相同为事件$A$，至少有两个人生日相同为事件$B$，则
\[P(A) = \frac{{C_{365}^n}\cdot n!}{{{{365}^n}}},P(B) = 1 - P(A)\]
 当 $n>365$ 时，互不相同的概率为 $0$，至少两人同生日的概率为 $1$。\end{solve}

\subsection*{训练题~\ref{q:Y1.9} 恰好得到一副手套}\label{a:Y1.9}{\small 题目见第~\pageref{q:Y1.9} 页；来源：闫一博公开解答，原题 1.9。}
\begin{solve}先选择唯一完整的一副，共 $6$ 种；剩下两只必须分别来自另外五副中的不同两副，有 $\binom52$ 种选副方法，每副选左或右各有两种。因此\[P=\frac{6\binom52\,2^2}{\binom{12}4}=\frac{240}{495}=\frac{16}{33}.\]若剩下两只来自同一副，就会形成第二副匹配手套，必须排除。\end{solve}

\subsection*{训练题~\ref{q:Y1.10} 两次抽样的交集}\label{a:Y1.10}{\small 题目见第~\pageref{q:Y1.10} 页；来源：闫一博公开解答，原题 1.10。}
\begin{solve}
记两组号码中恰有$c$个号码相同为事件$A$，则
\[P(A) = \frac{{C_m^cC_{n - m}^{k - c}}}{{C_n^k}}\]
 组合数公式仅在 $\max(0,m+k-n)\le c\le\min(m,k)$ 的整数范围内为正，范围外概率为零。\end{solve}

\subsection*{训练题~\ref{q:Y1.11} 单位正方形中的圆域}\label{a:Y1.11}{\small 题目见第~\pageref{q:Y1.11} 页；来源：闫一博公开解答，原题 1.11。}
\begin{solve}
在区间$[0,1]$随机地取两点相互独立，因此可视作在$[0,1]\times [0,1]$中取$(x,y)$，满足$x^2+y^2<1$，也即以原点为圆心$1$为半径的圆，记任取两点平方和小于$1$为事件$A$，则
\[
P(A)=\frac{\pi}{4}
\]
\end{solve}

\subsection*{训练题~\ref{q:Y1.12} 到达时差与等待}\label{a:Y1.12}{\small 题目见第~\pageref{q:Y1.12} 页；来源：闫一博公开解答，原题 1.12。}
\begin{solve}
记第一人到达的时间为$x$，第二人到达的时间为$y$，即求$P(|x-y|>0.5h)$，记一人至少等另一人半小时为事件$A$，则
\[P(A) = \frac{{2 \times 0.5 \times (0.5 \times 0.5)}}{1\times 1} = \frac{1}{4}\]
\end{solve}

\subsection*{训练题~\ref{q:Y1.13} 硬币落格与边界区域}\label{a:Y1.13}{\small 题目见第~\pageref{q:Y1.13} 页；来源：闫一博公开解答，原题 1.13。}
\begin{solve}圆心到每条格线的距离必须大于半径 $1/2$。当 $a\le1$ 时，允许区域面积为零；当 $a>1$ 时，允许区域是边长 $a-1$ 的正方形。因此所求概率为\[p(a)=\begin{cases}0,&0<a\le1,\\(1-1/a)^2,&a>1.\end{cases}\]对 $a>1$，不等式等价于 $1-1/a<0.1$，即 $a<10/9$。合并两种情况得到 $0<a<10/9$。边界接触的概率为零，不影响结果。\end{solve}

\subsection*{训练题~\ref{q:Y1.14} 两次服务是否相遇}\label{a:Y1.14}{\small 题目见第~\pageref{q:Y1.14} 页；来源：闫一博公开解答，原题 1.14。}
\begin{solve}设两车到达时刻为 $X,Y$，以一小时起点为零、分钟为单位。它们在 $[0,60]^2$ 上均匀分布。不相遇等价于 $Y>X+5$ 或 $X>Y+8$，即两块互不相交的直角三角形。面积分别为 $55^2/2$ 和 $52^2/2$，故概率为\[\frac{55^2+52^2}{2\cdot60^2}=\frac{5729}{7200}.\]等号边界概率为零；先到哪一车决定应加哪一个服务时长。\end{solve}

\subsection*{训练题~\ref{q:Y1.15} 立方体内的球域}\label{a:Y1.15}{\small 题目见第~\pageref{q:Y1.15} 页；来源：闫一博公开解答，原题 1.15。}
\begin{solve}
  在区间$[0,1]$随机地取三点相互独立，因此可视作在$[0,1]\times [0,1]\times [0,1]$中取$(X,Y,Z)$，满足$X^2+Y^2+Z^2\leqslant 1$和$X^2+Y^2+Z^2=1$，也即以原点为球心$1$为半径的球体和球面，则
\[
P(X^2+Y^2+Z^2\leqslant 1)=\frac{4\pi}{3}\times\frac{1}{8}=\frac{\pi}{6}
\]
根据几何概型的特点，区域边界的测度为零，故
\[
P(X^2+Y^2+Z^2=1)=0
\]  
\end{solve}

\subsection*{训练题~\ref{q:Y1.17} 概率的次可加性}\label{a:Y1.17}{\small 题目见第~\pageref{q:Y1.17} 页；来源：闫一博公开解答，原题 1.17。}
\begin{solve}
取$A_0=\varnothing$，${B_j} = {A_j} - \bigcup\limits_{i = 0}^{j - 1} {{A_i}} $，$B_j$互不相容，由于$B_j\subset A_j$，故$ \bigcup\limits_{j = 1}^\infty  {{B_j}}  \subset \bigcup\limits_{j = 1}^\infty  {{A_j}} $.下面验证$ \bigcup\limits_{j = 1}^\infty  {{A_j}}  \subset \bigcup\limits_{j = 1}^\infty  {{B_j}} $，任取$x\in \bigcup\limits_{j = 1}^\infty  {{A_j}}$，存在$i\geqslant 1$（取首个包含$x$的$A_j$对应的下标为$i$），使得$x\in A_i$，但$x\notin A_0,A_1,\cdots,A_{i-1}$（若$i=1$，则$x\notin A_0$），故$x\in B_i$，于是有$ \bigcup\limits_{j = 1}^\infty  {{A_j}}  \subset \bigcup\limits_{j = 1}^\infty  {{B_j}} $，这表明
\[\bigcup\limits_{j = 1}^\infty  {{A_j}}  = \bigcup\limits_{j = 1}^\infty  {{B_j}}\] 根据可列可加性和单调性
\[\begin{aligned}
P(\bigcup\limits_{j = 1}^\infty  {{A_j}} ) &= P(\bigcup\limits_{j = 1}^\infty  {{B_j})}  \\&= \sum\limits_{j = 1}^\infty  {P({B_j})}  \leqslant \sum\limits_{j = 1}^\infty  {P({A_j})} 
\end{aligned}\]
\end{solve}

\subsection*{训练题~\ref{q:Y1.18} 可数个概率一事件}\label{a:Y1.18}{\small 题目见第~\pageref{q:Y1.18} 页；来源：闫一博公开解答，原题 1.18。}
\begin{solve}
应用De Morgan定律
\[\bigcap\limits_{j = 1}^\infty  {{A_j}}  = \overline {\bigcup\limits_{j = 1}^\infty  {\overline {{A_j}} } } \]
故
\[\begin{aligned}P(\bigcap\limits_{j = 1}^\infty  {{A_j}} ) &= P(\overline {\bigcup\limits_{j = 1}^\infty  {\overline {{A_j}} } } ) \\&= 1 - P(\bigcup\limits_{j = 1}^\infty  {\overline {{A_j}} } )\end{aligned}\]
根据概率的次可加性
\[\begin{aligned}0 \leqslant P(\bigcup\limits_{j = 1}^\infty  {\overline {{A_j}} } ) \leqslant \sum\limits_{j = 1}^\infty  {P(\overline {{A_j}} )}  &= \sum\limits_{j = 1}^\infty  {(1 - P({A_j}))}  \\&= \sum\limits_{j = 1}^\infty  0\\&=0 \end{aligned}\]
综上可得
\[P(\bigcap\limits_{j = 1}^\infty  {{A_j}} ) = 1 - P(\bigcup\limits_{j = 1}^\infty  {\overline {{A_j}} } ) = 1\]
同理可证，$P(\bigcap\limits_{j = 1}^n {{A_j}} ) = 1$
\end{solve}

\subsection*{训练题~\ref{q:Y1.19} 四个互斥区域的分割}\label{a:Y1.19}{\small 题目见第~\pageref{q:Y1.19} 页；来源：闫一博公开解答，原题 1.19。}
\begin{solve}四个事件按是否属于 $A$、是否属于 $B$ 划分整个样本空间。它们两两互斥，且并为 $\Omega$。因此由有限可加性，\[\begin{aligned}&P(AB)+P(\overline AB)+P(A\overline B)+P(\overline A\,\overline B)\\&\qquad=P(\Omega)=1.\end{aligned}\]这不需要 $A,B$ 独立。\end{solve}

\subsection*{训练题~\ref{q:Y1.20} 递增事件列的极限}\label{a:Y1.20}{\small 题目见第~\pageref{q:Y1.20} 页；来源：闫一博公开解答，原题 1.20。}
\begin{solve}
$\{A_j\}$是单调增序列，故
\[\begin{aligned}
P(\bigcup\limits_{j = 1}^\infty  {{A_j}} ) &= \mathop {\lim }\limits_{n \to \infty } P({A_n}) \\&= \mathop {\lim }\limits_{n \to \infty } {(1 - \frac{1}{j})^j} \\&= \frac{1}{e}
\end{aligned}\]
\end{solve}

\clearpage\section*{条件概率、独立性与贝叶斯公式：详解}\addcontentsline{toc}{section}{条件概率、独立性与贝叶斯公式：详解}

\subsection*{训练题~\ref{q:Y2.1} 至少一次求职成功}\label{a:Y2.1}{\small 题目见第~\pageref{q:Y2.1} 页；来源：闫一博公开解答，原题 2.1。}
\begin{solve}
记甲被单位$A$录取为事件$A$，被单位$B$录取为事件$B$，根据加法公式，
\[P(A \cup B) = P(A) + P(B) - P(AB) = 0.4 + 0.6 - 0.24 = 0.76\]
\end{solve}

\subsection*{训练题~\ref{q:Y2.2} 自习室占用问题}\label{a:Y2.2}{\small 题目见第~\pageref{q:Y2.2} 页；来源：闫一博公开解答，原题 2.2。}
\begin{solve}
每个学生走进自习室相互独立

（a）记第$i$个自习室无学生为事件$A_i$，根据容斥公式
\[P({A_1} \cup {A_2} \cup {A_3}) = \sum\limits_{i = 1}^3 {P({A_i})}  - \sum\limits_{1 \leqslant i < j \leqslant 3}^3 {P({A_i}{A_j})}  + P(\bigcap\limits_{i = 1}^3 {{A_i}} )\]
其中
\[\begin{aligned}P({A_i}) &= \frac{{{2^9}}}{{{3^9}}},P({A_i}{A_j}) \\&= \frac{{{1^9}}}{{{3^9}}},P(\bigcap\limits_{i = 1}^3 {{A_i}} ) \\&= 0\end{aligned}\]
代入可得
\[\begin{aligned}P({A_1} \cup {A_2} \cup {A_3}) &= 3 \cdot \frac{{{2^9}}}{{{3^9}}} - 3 \cdot \frac{{{1^9}}}{{{3^9}}} + 0 \\&= \frac{{511}}{{6561}}\end{aligned}\]

（b）记每个自习室至少有一个学生为事件$B$，则
\[P(B) = 1 - P({A_1} \cup {A_2} \cup {A_3}) = \frac{{6050}}{{6561}}\]
\end{solve}

\subsection*{训练题~\ref{q:Y2.3} 无人拿到自己帽子}\label{a:Y2.3}{\small 题目见第~\pageref{q:Y2.3} 页；来源：闫一博公开解答，原题 2.3。}
\begin{solve}
记第$i$个人没有取到自己的帽子为事件$A_i$，则$\overline{A}_i$表示第$i$个人取到了自己的帽子，根据De Morgan定律
\[P(\bigcap\limits_{i = 1}^6 {{A_i}} ) = 1 - P(\bigcup\limits_{i = 1}^6 {{{\overline A }_i}} )\]
根据容斥公式
\[\begin{aligned}P(\bigcup\limits_{i = 1}^6 {{{\overline A }_i}} ) &= \sum\limits_{n = 1}^6 {{{( - 1)}^{n - 1}}C_6^n\frac{{(6 - n)!}}{{6!}}}  \\&= \sum\limits_{n = 1}^6 {{{( - 1)}^{n - 1}}\frac{1}{{n!}}} \end{aligned}\]
故
\[\begin{aligned}P(\bigcap\limits_{i = 1}^6 {{A_i}} ) &= 1 - \sum\limits_{n = 1}^6 {{{( - 1)}^{n - 1}}\frac{1}{{n!}}}  \\&= \sum\limits_{n = 0}^6 {{{( - 1)}^n}\frac{1}{{n!}}} \end{aligned}\]
\end{solve}

\subsection*{训练题~\ref{q:Y2.4} 有放回点名的覆盖概率}\label{a:Y2.4}{\small 题目见第~\pageref{q:Y2.4} 页；来源：闫一博公开解答，原题 2.4。}
\begin{solve}
（a）记一名同学一学期不被点到为事件$A$，由于每次点名独立，故
\[P(A) = {\left( {1 - \frac{1}{n}} \right)^m}\]

（b）记第$i$个学生未被点到为事件$A_i$，则至少有一个学生未被点到的概率为$P(\bigcup\limits_{i = 1}^n {{A_i}})$，由于\[P({A_i}) = \frac{{{{(n - 1)}^m}}}{{{n^m}}},P({A_i}{A_j}) = \frac{{{{(n - 2)}^m}}}{{{n^m}}}, \cdots \]根据容斥公式
\[P(\bigcup\limits_{i = 1}^n {{A_i}} ) = \sum\limits_{k = 1}^n {{{( - 1)}^{k - 1}}} C_n^k\frac{{{{(n - k)}^m}}}{{{n^m}}}\]
\end{solve}

\subsection*{训练题~\ref{q:Y2.5} 信封匹配与容斥}\label{a:Y2.5}{\small 题目见第~\pageref{q:Y2.5} 页；来源：闫一博公开解答，原题 2.5。}
\begin{solve}记 $A_i$ 为第 $i$ 封信匹配的事件。任意指定 $k$ 封信都匹配时，其余信件可任意排列，故概率为 $(n-k)!/n!$。容斥公式给出\[\begin{aligned}P\left(\bigcup_{i=1}^nA_i\right)&=\sum_{k=1}^n(-1)^{k-1}\binom nk\frac{(n-k)!}{n!}\\&=\sum_{k=1}^n\frac{(-1)^{k-1}}{k!}.\end{aligned}\]特别地，当 $n\to\infty$ 时，它趋于 $1-e^{-1}$。使用容斥的原因是多个匹配事件会同时发生，简单相加会重复计数。\end{solve}

\subsection*{训练题~\ref{q:Y2.6} 独立任务的共同成功}\label{a:Y2.6}{\small 题目见第~\pageref{q:Y2.6} 页；来源：闫一博公开解答，原题 2.6。}
\begin{solve}
记第$j$个人破译成功为事件$A_j$，根据De Morgan定律
\[P(\bigcup\limits_{j = 1}^6 {{A_j}} ) = 1 - P(\bigcap\limits_{j = 1}^6 {{{\overline A }_j}} )\]
由于$6$个人未成功破译密码相互独立，故
\[P(\bigcup\limits_{j = 1}^6 {{A_j}} ) = 1 - \prod\limits_{j = 1}^6 {(1 - {p_j})} \]
\end{solve}

\subsection*{训练题~\ref{q:Y2.7} 发现信号的试验次数}\label{a:Y2.7}{\small 题目见第~\pageref{q:Y2.7} 页；来源：闫一博公开解答，原题 2.7。}
\begin{solve}
记甲发出的第$j$个信号被接受为事件$A_j$，根据De Morgan定律
\[P(\bigcup\limits_{j = 1}^n {{A_j}} ) = 1 - P(\bigcap\limits_{j = 1}^n {{{\bar A}_j}} )\]
由于每个信号未被接受相互独立，故
\[P(\bigcup\limits_{j = 1}^n {{A_j}} ) = 1 - \prod\limits_{j = 1}^n {0.9}  = 1 - {0.9^n}\]
由题意知
\[1 - {0.9^n} \geqslant 0.95 \Rightarrow {0.9^n} < 0.05 \Rightarrow n > {\log _{0.9}}(0.05) \Rightarrow n \geqslant 29\]
\end{solve}

\subsection*{训练题~\ref{q:Y2.9} 保证成功概率的人员数}\label{a:Y2.9}{\small 题目见第~\pageref{q:Y2.9} 页；来源：闫一博公开解答，原题 2.9。}
\begin{solve}
记$12$小时内密码未被破译为事件$A$，记第$i$名专家未破译密码为事件$A_i$，第$j$名破译员未破译密码为事件$B_j$，现求破译员人数使得$P(A)<0.05$，即求$n$使得
\[P(A) = P({A_1} \cdots {A_5}{B_1} \cdots {B_n}) < 0.05\]
由于每人未破译密码相互独立，故
\[P(A) = {(0.8)^5}{(0.95)^n}\]
于是有
\[{(0.8)^5}{(0.95)^n} < 0.05 \Rightarrow n \geqslant 37\]
\end{solve}

\subsection*{训练题~\ref{q:Y2.10} 德梅雷的骰子问题}\label{a:Y2.10}{\small 题目见第~\pageref{q:Y2.10} 页；来源：闫一博公开解答，原题 2.10。}
\begin{proof}
记投掷$24$次两个骰子没有出现一对$6$为事件$A$，则
\[P(A) = {\left( {\frac{{35}}{{36}}} \right)^{24}} > \frac{1}{2}\]
这表明掷$24$次两个骰子至少掷出一对$6$的概率小于二分之一.
\end{proof}

\subsection*{训练题~\ref{q:Y2.11} 条件成绩概率}\label{a:Y2.11}{\small 题目见第~\pageref{q:Y2.11} 页；来源：闫一博公开解答，原题 2.11。}
\begin{solve}
记甲微积分得优为事件$A$，线性代数得优为事件$B$

（a）
\[P(AB) = P(A)P(B|A) = ac\]

（b）\[P(B|\bar A) = \frac{{P(\bar AB)}}{{P(\bar A)}} = \frac{{P(B) - P(AB)}}{{P(\bar A)}} = \frac{{b - ac}}{{1 - a}}\]

（c）
\[P(B) = P(A)P(B|A) + P(\bar A)P(B|\bar A) = b\]

（d）
\[P(A \cup B) = P(A) + P(B) - P(AB) = a + b - ac\]

（e）
\[P(\overline A  \cup \overline B ) = 1 - P(AB) = 1 - ac\]
\end{solve}

\subsection*{训练题~\ref{q:Y2.12} 连续受损的生存概率}\label{a:Y2.12}{\small 题目见第~\pageref{q:Y2.12} 页；来源：闫一博公开解答，原题 2.12。}
\begin{solve}令 $A_i$ 表示第 $i$ 次落地后仍未损坏。题目给出的是在之前完好的条件下的生存概率，故\[\begin{aligned}P(A_1A_2A_3)&=P(A_1)P(A_2\mid A_1)P(A_3\mid A_1A_2)\\&=0.5\times0.3\times0.1=0.015.\end{aligned}\]这里使用条件概率乘法公式，不需要把三次摔落后是否损坏当作独立事件。\end{solve}

\subsection*{训练题~\ref{q:Y2.13} 已知一枚骰子的点数}\label{a:Y2.13}{\small 题目见第~\pageref{q:Y2.13} 页；来源：闫一博公开解答，原题 2.13。}
\begin{solve}
记第一个骰子掷出点数$i$为事件$A_i$，第二个骰子掷出点数$i$为事件$B_i$，则$A_3\cup B_3$表示至少有一个骰子的点数是$3$

（a）记两个骰子点数之和为$7$为事件$A$，则
\[P(A|{A_3} \cup {B_3}) = \frac{{P({A_4}{B_3} \cup {A_3}{B_4})}}{{P({A_3} \cup {B_3})}} = \frac{2}{{11}}\]

（b）记两个骰子点数之和为$6$为事件$B$，则
\[P(B|{A_3} \cup {B_3}) = \frac{{P({A_3}{B_3})}}{{P({A_3} \cup {B_3})}} = \frac{1}{{11}}\]
应当猜$7$.
\end{solve}

\subsection*{训练题~\ref{q:Y2.15} 两孩家庭中的条件事件}\label{a:Y2.15}{\small 题目见第~\pageref{q:Y2.15} 页；来源：闫一博公开解答，原题 2.15。}
\begin{solve}
记第一个孩子是女孩为事件$A$，第二个孩子是女孩为事件$B$，则
\[P(AB|A \cup B) = \frac{{P(AB)}}{{P(A \cup B)}} = \frac{{{{0.49}^2}}}{{1 - {{0.51}^2}}} = 0.3245\]
\end{solve}

\subsection*{训练题~\ref{q:Y2.16} 通知渠道的全概率}\label{a:Y2.16}{\small 题目见第~\pageref{q:Y2.16} 页；来源：闫一博公开解答，原题 2.16。}
\begin{solve}
（a）记张经理收到通知为事件$A$，根据全概率公式
\[P(A) = 0.3 \times 0.8 + 0.2 \times 0.9 + 0.5 \times 0.4 = 0.62\]

（b）记张经理按时来见老板为事件$B$，根据全概率公式
\[\begin{aligned}P(B)&=P(A)P(B\mid A)+P(\overline A)P(B\mid\overline A)\\&=0.62\times0.95+0.38\times0=0.589\end{aligned}\]
因此张经理不能按时来见老板的概率为$P(\overline{B})=1-P(B)=0.411$.
\end{solve}

\subsection*{训练题~\ref{q:Y2.17} 隐私调查的随机化回答}\label{a:Y2.17}{\small 题目见第~\pageref{q:Y2.17} 页；来源：闫一博公开解答，原题 2.17。}
\begin{solve}
（a）记调查家庭回答“是”为事件$A$，抽到红球为事件$B$，则$P(B)=p_0$,$P(\overline{B})=q_0$,$P(A|B)=p$,$P(A|\overline{B})=1-p$，根据全概率公式
\[{p_1} = P(A) = P(B)P(A|B) + P(\overline B )P(A|\overline B ) = {p_0}p + {q_0}(1 - p)\]
于是有
\[p = \frac{{{p_1} - {q_0}}}{{{p_0} + {q_0}}}\]
（b）将数据带入
\[\begin{aligned}p &= \frac{{{p_1} - {q_0}}}{{{p_0} - {q_0}}} \\&= \frac{{\frac{{195}}{{320}} - \frac{5}{8}}}{{\frac{3}{8} - \frac{5}{8}}} \\&= \frac{5}{{80}} \\&= 6.25\% \end{aligned}\]
\end{solve}

\subsection*{训练题~\ref{q:Y2.19} 比分差达到二的比赛}\label{a:Y2.19}{\small 题目见第~\pageref{q:Y2.19} 页；来源：闫一博公开解答，原题 2.19。}
\begin{solve}记甲最终获胜的概率为 $w$。第一轮甲连胜的概率为 $p^2$，乙连胜的概率为 $(1-p)^2$，双方各胜一局的概率为 $2p(1-p)$。在最后一种情况下，后续试验重新开始，故\[w=p^2+2p(1-p)w.\]每轮结束比赛的概率 $p^2+(1-p)^2>0$，从而比赛以概率 $1$ 结束。解得\[w=\frac{p^2}{p^2+(1-p)^2}.\]这里的赛制是固定的两局一轮；若改成任意相邻两局连胜即结束，递推状态和答案都需要改变。\end{solve}

\subsection*{训练题~\ref{q:Y2.20} 加强抽样的波利亚球袋}\label{a:Y2.20}{\small 题目见第~\pageref{q:Y2.20} 页；来源：闫一博公开解答，原题 2.20。}
\begin{solve}连续三次取黑的概率按条件概率相乘：\[P(B_1B_2B_3)=\frac{b(b+1)(b+2)}{(b+r)(b+r+1)(b+r+2)}.\]对任意初始正整数 $b,r$，将同时归纳证明 $P(B_n)=b/(b+r)$。$n=1$ 成立。若所有初始构成都满足第 $n-1$ 步结论，条件于第一步后，袋子分别变为 $(b+1,r)$ 或 $(b,r+1)$，于是\[\begin{aligned}P(B_n)&=\frac{b}{b+r}\frac{b+1}{b+r+1}+\frac{r}{b+r}\frac{b}{b+r+1}\\&=\frac{b}{b+r}.\end{aligned}\]因此 $P(B_2),P(B_3)$ 也都等于最初黑球比例。各次取球的边缘概率相同，但取球结果并不独立。\end{solve}

\subsection*{训练题~\ref{q:Y2.21} 已知正点反推交通方式}\label{a:Y2.21}{\small 题目见第~\pageref{q:Y2.21} 页；来源：闫一博公开解答，原题 2.21。}
\begin{solve}
记甲乘火车为事件$A$，甲正点到达为事件$B$，根据全概率公式
\[P(B) = P(A)P(B|A) + P(\overline A )P(B|\overline A ) = 0.4 \times 0.9 + 0.6 \times 0.8 = 0.84\]
再根据贝叶斯公式，有
\[P(A|B) = \frac{{P(A)P(B|A)}}{{P(B)}} = \frac{{0.4 \times 0.9}}{{0.84}} = \frac{3}{7}\]
\end{solve}

\subsection*{训练题~\ref{q:Y2.22} 观察合格品后的状态}\label{a:Y2.22}{\small 题目见第~\pageref{q:Y2.22} 页；来源：闫一博公开解答，原题 2.22。}
\begin{solve}
记机床工作状态良好为事件$A$，新开机器后生产的产品是合格品为事件$B$，根据全概率公式和贝叶斯公式
\[P(A|B) = \frac{{P(A)P(B|A)}}{{P(B)}} = \frac{{0.95 \times 0.99}}{{0.95 \times 0.99 + 0.05 \times 0.5}} = 0.9741\]
\end{solve}

\subsection*{训练题~\ref{q:Y2.23} 低先验概率下的连续检测}\label{a:Y2.23}{\small 题目见第~\pageref{q:Y2.23} 页；来源：闫一博公开解答，原题 2.23。}
\begin{solve}记真实患病事件为 $A$，每次阳性为 $B_i$。以下只按题设的抽象条件概率计算。一次阳性后，\[P(A\mid B_1)=\frac{0.0001}{0.0001+0.9999(0.05)}\approx0.00199621.\]两次检测在给定真实状态后独立，所以真阳性的联合条件概率为 $1$，两次假阳性的联合条件概率为 $0.05^2$。因此\[P(A\mid B_1B_2)=\frac{0.0001}{0.0001+0.9999(0.05)^2}\approx0.03846524.\]直接使用原始先验与联合似然，可避免把第一次后验过早舍入后再传入第二次计算。条件独立与无条件独立不是同一假设。\end{solve}

\subsection*{训练题~\ref{q:Y2.25} 目击证据与先验比例}\label{a:Y2.25}{\small 题目见第~\pageref{q:Y2.25} 页；来源：闫一博公开解答，原题 2.25。}
\begin{solve}记车型为富康的事件为 $A$，证人报告富康为 $B$。根据题设，先验为 $0.15$，两种真实车型下报告的概率为 $0.9$ 与 $0.2$。于是\[P(A\mid B)=\frac{0.15(0.9)}{0.15(0.9)+0.85(0.2)}=\frac{27}{61}\approx0.442623.\]给定车型后，两名证人报告条件独立，故两次相同指证后的后验为\[\frac{0.15(0.9)^2}{0.15(0.9)^2+0.85(0.2)^2}=\frac{243}{311}\approx0.781350.\]这里的先验来自题设中两车型相同的肇事概率假设；若各车型发生事件的概率不同，应先修正事件总体中的车型比例。\end{solve}

\subsection*{训练题~\ref{q:Y2.26} 无穷多次成功的条件}\label{a:Y2.26}{\small 题目见第~\pageref{q:Y2.26} 页；来源：闫一博公开解答，原题 2.26。}
\begin{solve}
甲每次能否击中目标相互独立，根据博雷尔—坎泰利引理，记甲第$j$次击中目标为事件$A_j$，若
\[\sum\limits_{j = 1}^\infty  {P({A_j})}  = \sum\limits_{j = 1}^\infty  {{j^{ - \alpha }}}  = \infty \]
则$P(A_n\: \text{i.o.})=1$，故$\alpha\leqslant 1$时甲能够击中目标无穷多次.
\end{solve}

\subsection*{训练题~\ref{q:Y2.27} 最终不再失败的条件}\label{a:Y2.27}{\small 题目见第~\pageref{q:Y2.27} 页；来源：闫一博公开解答，原题 2.27。}
\begin{solve}
记第$j$次试验失败为事件$A_j$，则$P(A_j)=j^{-\alpha}$，根据博雷尔—坎泰利引理
\[\left\{ \begin{gathered}
  P({A_n}\,{\text{i}}.{\text{o}}{\text{.) = 0}} \hfill \quad \alpha>1 \\
  P({A_n}\,{\text{i}}.{\text{o}}{\text{.) = 1}} \hfill \quad \alpha\leqslant 1\\ 
\end{gathered}  \right.\]
这表明$\alpha>1$时，某次试验后所有试验不再失败的概率为$1$，$\alpha\leqslant 1$时，某次试验后所有试验不再失败的概率为$0$.
\end{solve}

\subsection*{训练题~\ref{q:Y2.28} 加红球后的黑球出现次数}\label{a:Y2.28}{\small 题目见第~\pageref{q:Y2.28} 页；来源：闫一博公开解答，原题 2.28。}
\begin{solve}
记第$j$次取球摸到黑球为事件$A_j$，根据题意
\[P({A_j}) = \frac{1}{{1 + (1 + 1 + 2 +  \cdots  + (j - 1))}} = \frac{1}{{2 + \frac{{j(j - 1)}}{2}}}\]
根据博雷尔—坎泰利引理，由于
\[\sum\limits_{j = 1}^\infty  {\frac{2}{{4 + j(j - 1)}}}  < \infty \]
故$P({A_n}\,{\text{i}}.{\text{o}}{\text{.) = 0}}$，能够取到黑球有限次.
\end{solve}

\subsection*{训练题~\ref{q:Y2.29} 黑红同时增长的抽球}\label{a:Y2.29}{\small 题目见第~\pageref{q:Y2.29} 页；来源：闫一博公开解答，原题 2.29。}
\begin{solve}第 $j$ 次抽取前，黑球有 $j$ 个，红球有 $1+j(j-1)/2$ 个，所以\[p_j=P(B_j)=\frac{2j}{j^2+j+2}.\]由于每次增加的球数与抽到的颜色无关，袋中构成是确定的；独立随机抽取机制使事件 $B_j$ 相互独立。又 $p_j\sim2/j$，故 $\sum_jp_j=\infty$。由第二博雷尔—坎泰利引理，$P(B_j\text{ 无穷多次发生})=1$。若增球数量依赖抽取结果，就不能未经检查直接使用这个独立性论证。\end{solve}

\subsection*{训练题~\ref{q:Y2.30} 全概率公式的证明}\label{a:Y2.30}{\small 题目见第~\pageref{q:Y2.30} 页；来源：闫一博公开解答，原题 2.30。}
\begin{proof}
根据题意有
\[\begin{aligned}B \subset \bigcup\limits_{j = 1}^\infty  {{A_j}}  \Rightarrow B &= (\bigcup\limits_{j = 1}^\infty  {{A_j}} )B \Rightarrow B \\&= \bigcup\limits_{j = 1}^\infty  {({A_j}} B)\end{aligned}\]
由于$A_1,A_2,\cdots$互不相容，故$A_jB$也互不相容，因此有
\[P(B) = P(\bigcup\limits_{j = 1}^\infty  {({A_j}} B)) = \sum\limits_{j = 1}^\infty  {P({A_j}B)} \]
再根据乘法公式
\[P(B) = \sum\limits_{j = 1}^\infty  {P({A_j}B)}  = \sum\limits_{j = 1}^\infty  {P({A_j})P(B|{A_j})} \]
\end{proof}

\subsection*{训练题~\ref{q:Y2.31} 贝叶斯公式的证明}\label{a:Y2.31}{\small 题目见第~\pageref{q:Y2.31} 页；来源：闫一博公开解答，原题 2.31。}
\begin{proof}
根据条件概率公式
  \[P({A_j}|B) = \frac{{P({A_j}B)}}{{P(B)}} = \frac{{P({A_j})P(B|{A_j})}}{{P(B)}}\]
由于$A_1,A_2,\cdots$互不相容，故$A_jB$也互不相容，再根据全概率公式
\[\begin{aligned}P({A_j}|B) &= \frac{{P({A_j})P(B|{A_j})}}{{P(B)}} \\&= \frac{{P({A_j})P(B|{A_j})}}{{\sum\limits_{i = 1}^\infty  {P({A_i})P(B|{A_i})} }}\end{aligned}\]

\end{proof}

\subsection*{训练题~\ref{q:Y2.32} 串联与并联可靠性}\label{a:Y2.32}{\small 题目见第~\pageref{q:Y2.32} 页；来源：闫一博公开解答，原题 2.32。}
\begin{solve}
记元件$A$在一年中烧坏为事件$A$，元件$B$在一年中烧坏为事件$B$

（a）并联时，两元件均烧坏则系统断电，两元件烧断是相互独立的
\[P(AB) = P(A)P(B) = {p_1}{p_2}\]
（b）串联时，两元件至少有一个烧坏则系统断电，根据加法公式
\[P(A \cup B) = P(A) + P(B) - P(AB) = {p_1} + {p_2} - {p_1}{p_2}\]
\end{solve}

\subsection*{训练题~\ref{q:Y2.33} 比赛局数与胜率}\label{a:Y2.33}{\small 题目见第~\pageref{q:Y2.33} 页；来源：闫一博公开解答，原题 2.33。}
\begin{solve}令 $p>1/2$ 为甲每局获胜概率，$q=1-p$。三局两胜的概率为 $P_3=p^2+2p^2q$；五局三胜的概率为 $P_5=p^3+3p^3q+6p^3q^2$。按决胜局前已经取得的胜负数计数，以上各项互不重叠。化简得\[P_5-P_3=3p^2q^2(2p-1)>0.\]所以在各局独立且胜率固定的模型下，实力较强的一方更偏好五局三胜；若每局胜率随体力或其他状态变化，这个比较不能直接沿用。\end{solve}

\subsection*{训练题~\ref{q:Y2.34} 集体迟到的提前量}\label{a:Y2.34}{\small 题目见第~\pageref{q:Y2.34} 页；来源：闫一博公开解答，原题 2.34。}
\begin{solve}
由于每个人是否迟到相互独立，记无人迟到$t$min以上为事件$A_t$，则
\[P({A_t}) = {(1 - \frac{1}{{100{{(t + 1)}^2}}})^{1000}} > 0.99\]
计算可得$t\geqslant 31$min.
\end{solve}

\clearpage\section*{随机变量及其分布：详解}\addcontentsline{toc}{section}{随机变量及其分布：详解}

\subsection*{训练题~\ref{q:Y3.1} 分布函数的无穷远极限}\label{a:Y3.1}{\small 题目见第~\pageref{q:Y3.1} 页；来源：闫一博公开解答，原题 3.1。}
\begin{solve}
\[\begin{aligned}1&=P( - \infty  < X < \infty ) \\&= \sum\limits_{n =  - \infty }^\infty  {P(n < X \leqslant n + 1)}  \\&= \sum\limits_{n =  - \infty }^\infty  {F(n + 1) - F(n)} \end{aligned}\]
于是有
\[\mathop {\lim }\limits_{n \to \infty } F(n) - \mathop {\lim }\limits_{m \to  - \infty } F(m) = 1\]
由于分布函数是单调不减的有界函数，故极限存在（上式有意义），且
\[\begin{aligned}\mathop {\lim }\limits_{n \to   \infty } F\left( n \right) &= \mathop {\lim }\limits_{x \to   \infty } F(x),\mathop {\lim }\limits_{m \to  - \infty } F\left( m \right) \\&= \mathop {\lim }\limits_{x \to  - \infty } F(x)\end{aligned}\]
由于$0\leqslant F(x)\leqslant 1$，由$F(x)$的单调性可知
\[\begin{aligned}\mathop {\lim }\limits_{x \to \infty } F(x) &= F(\infty ) \\&= 1,\mathop {\lim }\limits_{x \to  - \infty } F(x) \\&= F( - \infty ) \\&= 0\end{aligned}\]
\end{solve}

\subsection*{训练题~\ref{q:Y3.2} 指数型密度的归一化}\label{a:Y3.2}{\small 题目见第~\pageref{q:Y3.2} 页；来源：闫一博公开解答，原题 3.2。}
\begin{solve}
根据归一化条件
\[1 = \int_{ - \infty }^\infty  {g(x)dx}  = \int_0^\infty  {\lambda {e^{ - 3x}}dx}  = \frac{\lambda }{3}\]
故$\lambda=3$.
\end{solve}

\subsection*{训练题~\ref{q:Y3.3} 二次指数可积的条件}\label{a:Y3.3}{\small 题目见第~\pageref{q:Y3.3} 页；来源：闫一博公开解答，原题 3.3。}
\begin{proof}
反证假设$a\geqslant 0$，则有$\exp(ax^2)\geqslant 1$，于是有
\[\begin{aligned}\int_{ - \infty }^\infty  {g(x)dx}  &= \int_{ - \infty }^\infty  {\exp (a{x^2} + bx + c)dx}  \geqslant \int_{ - \infty }^\infty  {\exp (bx + c)dx}  \\&= \left. {\frac{1}{b}\exp (bx + c)} \right|_{ - \infty }^\infty \end{aligned}\]
由于$a,b,c$均为常数，故无论上式中$b$取何值，均不满足归一化条件，反正假设不真，$a<0$.
\end{proof}

\subsection*{训练题~\ref{q:Y3.4} 最大值与最小值恒等式}\label{a:Y3.4}{\small 题目见第~\pageref{q:Y3.4} 页；来源：闫一博公开解答，原题 3.4。}
\begin{solve}对每个具体样本点，若 $X\ge Y$，则 $|X-Y|=X-Y$，代入两式分别得到 $X$ 与 $Y$；若 $X<Y$，则 $|X-Y|=Y-X$，分别得到 $Y$ 与 $X$。因此\[\max(X,Y)=\frac{|X-Y|+X+Y}{2},\]\[\min(X,Y)=\frac{X+Y-|X-Y|}{2}.\]将这两个恒等式相加，绝对值项抵消，得到 $\max(X,Y)+\min(X,Y)=X+Y$。结论逐点成立，不需要任何分布、矩或独立性假设。\end{solve}

\subsection*{训练题~\ref{q:Y3.5} 分段密度的归一化}\label{a:Y3.5}{\small 题目见第~\pageref{q:Y3.5} 页；来源：闫一博公开解答，原题 3.5。}
\begin{solve}
由于$f(x)$是$X$的概率密度，根据归一化条件
\[1 = \int_{ - \infty }^\infty  {f(x)dx}  = \int_1^2 {\frac{{(x - 1)}}{a}dx}  + \int_3^5 {\frac{{(x - 3)}}{a}dx} \]
计算可得
\[1 = \frac{1}{{2a}} - \frac{5}{{2a}} + \frac{9}{{2a}} = \frac{5}{{2a}} \Rightarrow a = \frac{5}{2}\]
\end{solve}

\subsection*{训练题~\ref{q:Y3.6} 二项计数的尾概率}\label{a:Y3.6}{\small 题目见第~\pageref{q:Y3.6} 页；来源：闫一博公开解答，原题 3.6。}
\begin{solve}设 $X$ 为八条短信中相应类别的条数。独立性及每条概率为 $0.2$ 给出 $X\sim\operatorname{Bin}(8,0.2)$。因此\[P(X\ge3)=1-\sum_{k=0}^{2}\binom8k0.2^k0.8^{8-k}=0.20308224.\]取补事件时必须包含 $X=0$，不能从 $k=1$ 开始相减。\end{solve}

\subsection*{训练题~\ref{q:Y3.7} 两组二项计数的比较}\label{a:Y3.7}{\small 题目见第~\pageref{q:Y3.7} 页；来源：闫一博公开解答，原题 3.7。}
\begin{solve}
记甲击中目标的次数为随机变量$X$，乙击中目标的次数为随机变量$Y$，$X\sim B(3,0,6),Y\sim B(3,0,7)$

（a）甲乙击中目标相互独立，故
\[P(X = Y) = \sum\limits_{k = 0}^3 {P(X = k)P(Y = k)} \]
计算可得
\[P(X = Y) = 0.32076\]

（b）甲乙击中目标相互独立，故
\[P(X > Y) = \sum\limits_{0 \leqslant j < i \leqslant 3} {P(X = i)P(Y = j)} \]
计算可得
\[P(X > Y) = 0.243\]
\end{solve}

\subsection*{训练题~\ref{q:Y3.8} 首次与第多次成功}\label{a:Y3.8}{\small 题目见第~\pageref{q:Y3.8} 页；来源：闫一博公开解答，原题 3.8。}
\begin{solve}
（a）记第一个打印错误所在页数为随机变量$X$，则$X$服从参数为$0.05$的几何分布，故
\[P(X=k)=0.95^{k-1} 0.05\]
（b）记第$r$打印错误所在页数为随机变量$Y$，则$Y$服从参数为$(r,0.05)$的Pascal分布，故
\[P(Y = k) = C_{k - 1}^{r - 1}{0.95^{k - r}}{0.05^r}\]
（c）记发现第$r$个错误前，无打印错误的页数为随机变量$Z$，则$Z$服从参数为$(r,0.05)$的负二项分布，故 
\[P(Z = k) = C_{k + r - 1}^{r - 1}{0.95^k}{0.05^r}\]
（d）记$n$页书的打印错误数为随机变量$\xi$，则$\xi\sim B(n,0.05)$
\end{solve}

\subsection*{训练题~\ref{q:Y3.9} 泊松分裂与独立性}\label{a:Y3.9}{\small 题目见第~\pageref{q:Y3.9} 页；来源：闫一博公开解答，原题 3.9。}
\begin{solve}记总邮件数为 $N$，看到的数量为 $X$，过滤数量为 $Z=N-X$。给定 $N=n$，有 $X\sim\operatorname{Bin}(n,0.8)$，所以 $h_k=\binom nk0.8^k0.2^{n-k}$。对非负整数 $k,m$，\[\begin{aligned}P(X=k,Z=m)&=P(N=k+m)P(X=k\mid N=k+m)\\&=e^{-\lambda}\frac{\lambda^{k+m}}{k!m!}0.8^k0.2^m\\&=\left[e^{-0.8\lambda}\frac{(0.8\lambda)^k}{k!}\right]\left[e^{-0.2\lambda}\frac{(0.2\lambda)^m}{m!}\right].\end{aligned}\]联合质量函数分解为两个泊松质量函数，故 $X\sim\operatorname{Pois}(0.8\lambda)$、$Z\sim\operatorname{Pois}(0.2\lambda)$ 且独立。因此已知 $X=k$ 后，$Z$ 仍为率 $0.2\lambda$ 的泊松分布。\end{solve}

\subsection*{训练题~\ref{q:Y3.13} 周期信号的均匀等待}\label{a:Y3.13}{\small 题目见第~\pageref{q:Y3.13} 页；来源：闫一博公开解答，原题 3.13。}
\begin{solve}
记救援者在$t$时刻到达可收到呼叫信号的范围，求救者发出求救信号为随机变量$X$，则$X\sim U(t,t+2)$，故
\[0.95 \leqslant P(t < X \leqslant t + \Delta t) = \int_t^{t + \Delta t} {\frac{1}{2}dt}  = \frac{1}{2}\Delta t\]
这表明至少停留$1.9$min才能以$0.95$的概率收到呼叫.
\end{solve}

\subsection*{训练题~\ref{q:Y3.14} 恒定风险率与指数寿命}\label{a:Y3.14}{\small 题目见第~\pageref{q:Y3.14} 页；来源：闫一博公开解答，原题 3.14。}
\begin{solve}
记热敏电阻的寿命（失效时间）为随机变量$X$，根据题意
\[P(X \leqslant t + \Delta t|X > t) = \lambda \Delta t + o(\Delta t)\]
记$X$的分布函数为$F(t)$，相应的概率密度为$f(t)$，于是有
\[\begin{aligned}\lambda \Delta t + o(\Delta t) &= \frac{{P(t < X \leqslant t + \Delta t)}}{{P(X > t)}} \\&= \frac{{F(t + \Delta t) - F(t)}}{{1 - F(t)}}\end{aligned}\]
两端同除$\Delta t$，取极限可得
\[\begin{aligned}\mathop {\lim }\limits_{\Delta t \to 0} \lambda  + o(1) &= \frac{1}{{1 - F(t)}}\mathop {\lim }\limits_{\Delta t \to 0} \frac{{F(t + \Delta t) - F(t)}}{{\Delta t}} \Rightarrow \lambda  \\&= \frac{{f(t)}}{{1 - F(t)}}\end{aligned}\]
上式右端是一个变量分离方程，积分可得
\[ - \ln (1 - F(t)) = \lambda t + C \Rightarrow F(t) = 1 - C{e^{ - \lambda t}}\]
根据归一化条件可得
\[\int_0^\infty  {f(t)dt}  = 1 \Rightarrow C = 1\]
于是得到
\[F(t) = 1 - {e^{ - \lambda t}}\]
这表明$X\sim \varepsilon(\lambda)$.
\end{solve}

\subsection*{训练题~\ref{q:Y3.15} 正态尺寸的合格概率}\label{a:Y3.15}{\small 题目见第~\pageref{q:Y3.15} 页；来源：闫一博公开解答，原题 3.15。}
\begin{solve}
记机床加工的部件长度为随机变量$X$，则$X\sim N(8,36\times 10^{-8})$，于是有
\[\begin{aligned}P(8 - 0.0015 \leqslant X \leqslant 8 + 0.0015) &= P( - 2.5 \leqslant \frac{{X - 8}}{{0.0006}} \leqslant 2.5) \\&= 2\Phi (2.5) - 1 \\&= 0.9876\end{aligned}\]
\end{solve}

\subsection*{训练题~\ref{q:Y3.16} 泊松变量的平方根}\label{a:Y3.16}{\small 题目见第~\pageref{q:Y3.16} 页；来源：闫一博公开解答，原题 3.16。}
\begin{solve}
$P(Y\geqslant 0)=1$，由离散变量的一一对应变换
\[P(Y = \sqrt k ) = P(X = k) = \frac{{{\lambda ^k}}}{{k!}}{e^{ - \lambda }},k = 0,1,2, \cdots \]
\end{solve}

\subsection*{训练题~\ref{q:Y3.17} 均匀电流的平方变换}\label{a:Y3.17}{\small 题目见第~\pageref{q:Y3.17} 页；来源：闫一博公开解答，原题 3.17。}
\begin{solve}
根据题意$I\sim U(8,9)$，故有概率密度$f_I(i)=1.i\in (8,9)$，$P(W>0)=1$，于是有
\[P(W = w) = P(2{I^2} = w) = P(I = \sqrt {\frac{w}{2}} ) = \frac{1}{4}\sqrt {\frac{2}{w}} dw\]
这表明$W$有概率密度
\[{f_W}(w) = \frac{1}{4}\sqrt {\frac{2}{w}},w\in(128,162) \]
\end{solve}

\subsection*{训练题~\ref{q:Y3.18} 指数变量的仿射变换}\label{a:Y3.18}{\small 题目见第~\pageref{q:Y3.18} 页；来源：闫一博公开解答，原题 3.18。}
\begin{solve}
根据题意$X\sim \varepsilon(\lambda)$，故有概率密度$f_X(x)=\lambda e^{-\lambda x}$，由于$a>0$，故$P(Y>b)=1$，于是有
\[P(Y = y) = P(aX + b = y) = P(X = \frac{{y - b}}{a}) = \frac{\lambda }{a}{e^{ - \frac{{\lambda (y - b)}}{a}}}dy\]
这表明$Y$有概率密度
\[{f_Y}(y) = \frac{\lambda }{a}{e^{ - \frac{{\lambda (y - b)}}{a}}},y > b\]
\end{solve}

\subsection*{训练题~\ref{q:Y3.19} 半柯西变量的对数}\label{a:Y3.19}{\small 题目见第~\pageref{q:Y3.19} 页；来源：闫一博公开解答，原题 3.19。}
\begin{solve}
$P(Y\in \mathbb{R})=1$，于是有 
\[P(Y = y) = P(\ln X = y) = P(X = {e^y}) = \frac{{2{e^y}}}{{\pi (1 + {e^{2y}})}}dy\]
这表明$Y$有概率密度
\[{f_Y}(y) = \frac{{2{e^y}}}{{\pi (1 + {e^{2y}})}},y \in ( - \infty ,\infty )\]
\end{solve}

\subsection*{训练题~\ref{q:Y3.20} 倒数、绝对值与正切变换}\label{a:Y3.20}{\small 题目见第~\pageref{q:Y3.20} 页；来源：闫一博公开解答，原题 3.20。}
\begin{solve}使用分段单调变换的密度公式。连续分布在零点及正切奇点处的总概率为零，所以可在这些点任意定义变换。\par（a）对 $y\ne0$，逆变换为 $x=1/y$，导数绝对值为 $1/y^2$，故\[f_{1/X}(y)=y^{-2}f(1/y).\]（b）对 $y>0$，两个逆像为 $y,-y$，因此\[f_{|X|}(y)=f(y)+f(-y),\]负半轴密度为零。\par（c）正切函数在各区间 $(-\pi/2+k\pi,\pi/2+k\pi)$ 上单调，逆像为 $x=\arctan y+k\pi$。因此\[f_{\tan X}(y)=\frac1{1+y^2}\sum_{k\in\mathbb Z}f(\arctan y+k\pi).\]各分支贡献非负，可以用非负积分交换定理求和；所得密度公式几乎处处成立。\end{solve}

\subsection*{训练题~\ref{q:Y3.21} 二项众数与参数估计}\label{a:Y3.21}{\small 题目见第~\pageref{q:Y3.21} 页；来源：闫一博公开解答，原题 3.21。}
\begin{solve}二项质量函数相邻比值为\[\frac{p_{k+1}}{p_k}=\frac{n-k}{k+1}\frac{p}{1-p}.\]它大于一等价于 $k+1<(n+1)p$。本题 $(n+1)p=20(0.7)=14$ 是整数，故 $k=13$ 与 $k=14$ 都达到最大概率，不能只保留一个。观察到 $X=9$ 后，以 $p$ 为变量的似然核为 $p^9(1-p)^{10}$，对数导数为 $9/p-10/(1-p)$，驻点为 $p=9/19$，二阶导数为负，故它是极大似然估计。\end{solve}

\subsection*{训练题~\ref{q:Y3.22} 迟到人数的众数和估计}\label{a:Y3.22}{\small 题目见第~\pageref{q:Y3.22} 页；来源：闫一博公开解答，原题 3.22。}
\begin{solve}迟到人数 $X\sim\operatorname{Bin}(95,p)$。二项分布的相邻比值表明，当 $(n+1)p$ 不是整数时唯一众数为其整数部分；若是整数则有相邻两个众数。本题 $96(0.05)=4.8$，所以最可能迟到人数为 $4$。观察到七人迟到时，似然核为 $p^7(1-p)^{88}$，对数导数过零点为 $p=7/95\approx0.073684$，即极大似然估计。众数求的是固定参数下最可能的结果，估计求的是固定结果下最合适的参数。\end{solve}

\subsection*{训练题~\ref{q:Y3.23} 泊松众数与似然}\label{a:Y3.23}{\small 题目见第~\pageref{q:Y3.23} 页；来源：闫一博公开解答，原题 3.23。}
\begin{solve}泊松相邻概率比为 $p_{k+1}/p_k=\lambda/(k+1)$。当 $\lambda=23.8$ 时，唯一众数为 $23$。观察到 $X=21$ 后，对数似然为 $-\lambda+21\log\lambda-\log(21!)$，导数为 $-1+21/\lambda$，故极大似然估计为 $\widehat\lambda=21$。若泊松参数恰为正整数 $r$，则 $r-1,r$ 同为众数，非整数情形才具有唯一的整数部分众数。\end{solve}

\subsection*{训练题~\ref{q:Y3.24} 计数结果与参数的区别}\label{a:Y3.24}{\small 题目见第~\pageref{q:Y3.24} 页；来源：闫一博公开解答，原题 3.24。}
\begin{solve}若 $X\sim\operatorname{Pois}(\lambda)$，相邻概率之比为\[\frac{P(X=k+1)}{P(X=k)}=\frac{\lambda}{k+1}.\]当 $\lambda=8.9$ 时，概率先增加后减少，唯一众数为 $8$。另一方面，已观测到 $X=12$ 时，关于 $\lambda$ 的似然为 $L(\lambda)=e^{-\lambda}\lambda^{12}/12!$，其对数导数为 $-1+12/\lambda$，故极大似然估计为 $\widehat\lambda=12$。前问固定参数比较不同结果，后问固定结果比较不同参数，两个优化变量不同。\end{solve}

\subsection*{训练题~\ref{q:Y3.25} 骰子样本的最大最小值}\label{a:Y3.25}{\small 题目见第~\pageref{q:Y3.25} 页；来源：闫一博公开解答，原题 3.25。}
\begin{solve}假定各次骰子公平且独立。最小值至少为 $k$ 的概率为 $((7-k)/6)^n$，所以\[P(m=k)=\left(\frac{7-k}{6}\right)^n-\left(\frac{6-k}{6}\right)^n.\]最大值不超过 $k$ 的概率为 $(k/6)^n$，所以\[P(M=k)=\left(\frac{k}{6}\right)^n-\left(\frac{k-1}{6}\right)^n.\]要求最小值为二、最大值为五，等价于所有点数在 $\{2,3,4,5\}$ 且二与五都出现。容斥得\[P(m=2,M=5)=\left(\frac46\right)^n-2\left(\frac36\right)^n+\left(\frac26\right)^n.\]对 $n=1$ 该式为零，与两个极值不可能不同相符。\end{solve}

\subsection*{训练题~\ref{q:Y3.26} 圆周投影的反正弦密度}\label{a:Y3.26}{\small 题目见第~\pageref{q:Y3.26} 页；来源：闫一博公开解答，原题 3.26。}
\begin{solve}令 $\Theta\sim U(0,2\pi)$，则 $X=R\cos\Theta$。对 $-R<x<R$，满足 $\cos\Theta\le x/R$ 的角度集合长度为 $2\pi-2\arccos(x/R)$，所以\[F_X(x)=1-\frac{1}{\pi}\arccos\frac{x}{R}.\]当 $x\le-R$ 时分布函数为 $0$，当 $x\ge R$ 时为 $1$。在内部求导得\[f_X(x)=\frac{1}{\pi\sqrt{R^2-x^2}},\quad |x|<R,\]其余处密度为零。端点处密度可任意定义。连续随机变量的单点概率为零，因此不能把 $P(X=x)$ 当作密度。\end{solve}

\subsection*{训练题~\ref{q:Y3.28} 正弦变换的两条分支}\label{a:Y3.28}{\small 题目见第~\pageref{q:Y3.28} 页；来源：闫一博公开解答，原题 3.28。}
\begin{solve}归一化条件为 $c\int_0^\pi x\,dx/\pi^2=c/2=1$，所以 $c=2$。当 $0<y<1$ 时，两个逆像为 $x_1=\arcsin y$ 和 $x_2=\pi-\arcsin y$，导数绝对值都为 $1/\sqrt{1-y^2}$。于是\[f_Y(y)=\frac{f_X(x_1)+f_X(x_2)}{\sqrt{1-y^2}}=\frac{2}{\pi\sqrt{1-y^2}}.\]其余处密度为零。两条分支都要计入；只保留 $\arcsin y$ 会漏掉 $X>\pi/2$ 的贡献。\end{solve}

\subsection*{训练题~\ref{q:Y3.29} 逆分布函数生成正态数}\label{a:Y3.29}{\small 题目见第~\pageref{q:Y3.29} 页；来源：闫一博公开解答，原题 3.29。}
\begin{solve}因 $\Phi$ 严格递增、连续且从实轴映到 $(0,1)$，对任意实数 $y$，\[P(Y\le y)=P(\Phi^{-1}(X)\le y)=P(X\le\Phi(y))=\Phi(y).\]因此 $Y$ 的分布函数正是标准正态分布函数，故 $Y\sim N(0,1)$。这个论证直接使用分布函数，不需要把零概率单点事件当作密度。\end{solve}

\subsection*{训练题~\ref{q:Y3.31} 分段单调变换}\label{a:Y3.31}{\small 题目见第~\pageref{q:Y3.31} 页；来源：闫一博公开解答，原题 3.31。}
\begin{solve}变换在正半轴严格增加且在 $1$ 处连续，逆函数在 $0<y<1$ 为 $\sqrt y$，在 $y>1$ 为 $y$。因此\[f_Y(y)=\begin{cases}\lambda e^{-\lambda\sqrt y}/(2\sqrt y),&0<y<1,\\\lambda e^{-\lambda y},&y>1,\\0,&y<0.\end{cases}\]在 $y=1$ 处没有原子，密度值可任意指定。两段积分分别为 $P(X<1)=1-e^{-\lambda}$ 和 $P(X>1)=e^{-\lambda}$，总和为一。\end{solve}

\clearpage\section*{联合分布与变量变换：详解}\addcontentsline{toc}{section}{联合分布与变量变换：详解}

\subsection*{训练题~\ref{q:Y4.3} 联合概率表与边缘分布}\label{a:Y4.3}{\small 题目见第~\pageref{q:Y4.3} 页；来源：闫一博公开解答，原题 4.3。}
\begin{solve}固定一个变量的取值，对另一个变量求和即可得到边缘概率。具体地，$P(X=i)=\sum_{j=1}^5p_{ij}$，所以按行相加得到
\begin{center}
    
    \begin{tabular}{|c|c|c|c|c|c|} \hline 
         $X$&  $1$& $2$&$3$& $4$&$5$\\ \hline 
         $P$& 
     $0.18$& $0.33$& $0.17$& $0.12$&$0.20$\\ \hline\end{tabular}

\end{center}
同理，$P(Y=j)=\sum_{i=1}^5p_{ij}$，按列相加得到\begin{center}
    
    \begin{tabular}{|c|c|c|c|c|c|} \hline 
         $Y$&  $1$& $2$&$3$& $4$&$5$\\ \hline 
         $P$& 
     $0.28$& $0.28$& $0.21$& $0.11$&$0.12$\\ \hline\end{tabular}

\end{center}

两组边缘概率各自之和均为 $1$，可据此检查求和是否遗漏。\end{solve}

\subsection*{训练题~\ref{q:Y4.6} 同参数二项变量之和}\label{a:Y4.6}{\small 题目见第~\pageref{q:Y4.6} 页；来源：闫一博公开解答，原题 4.6。}
\begin{solve}概率生成函数在独立和下相乘：\[E s^{X+Y}=(1-p+ps)^n(1-p+ps)^m=(1-p+ps)^{n+m}.\]这正是 $\operatorname{Bin}(n+m,p)$ 的生成函数，所以\[P(X+Y=k)=\binom{n+m}{k}p^k(1-p)^{n+m-k}.\]也可以把 $X,Y$ 分别表示成 $n,m$ 个独立伯努利试验的成功数，合并即为 $n+m$ 次试验。两组成功概率必须相同；参数不同的二项和一般不再是二项分布。\end{solve}

\subsection*{训练题~\ref{q:Y4.8} 函数依赖的离散向量}\label{a:Y4.8}{\small 题目见第~\pageref{q:Y4.8} 页；来源：闫一博公开解答，原题 4.8。}
\begin{solve}
    \[\sum\limits_{i = 1}^8 {P(X = i,Y = \frac{1}{i})}  = 1 \Rightarrow c = \frac{1}{8}\]
    于是得到$X,Y$的概率分布为
    \[P(X = i) = \frac{1}{8},P(Y = \frac{1}{i}) = \frac{1}{8},i = 1,2, \cdots ,8\]
\end{solve}

\subsection*{训练题~\ref{q:Y4.9} 最小值的尾概率}\label{a:Y4.9}{\small 题目见第~\pageref{q:Y4.9} 页；来源：闫一博公开解答，原题 4.9。}
\begin{solve}设 $Z=\min(X,Y)$，则 $\{Z>a\}=\{X>a,Y>a\}$。独立同分布给出 $P(Z>a)=P(X>a)^2$。又因 $a<b$，事件 $\{Z>b\}$ 包含于 $\{Z>a\}$，故\[\begin{aligned}P(a<Z\le b)&=P(Z>a)-P(Z>b)\\&=P(X>a)^2-P(X>b)^2.\end{aligned}\]此证明同时适用于连续、离散和混合分布；端点的严格与非严格关系由事件差确定，不需要另加连续性假设。\end{solve}

\subsection*{训练题~\ref{q:Y4.15} 三角支撑上的边缘密度}\label{a:Y4.15}{\small 题目见第~\pageref{q:Y4.15} 页；来源：闫一博公开解答，原题 4.15。}
\begin{solve}
\[{f_X}(x) = \int_0^x {{e^{ - x}}dy}  = x{e^{ - x}}\]
    \[{f_Y}(y) = \int_y^\infty  {{e^{ - x}}dx}  = {e^{ - y}}\]
 以上分别在 $x>0$、$y>0$ 成立，其余处密度为零。\end{solve}

\subsection*{训练题~\ref{q:Y4.16} 随机家庭规模与男孩数}\label{a:Y4.16}{\small 题目见第~\pageref{q:Y4.16} 页；来源：闫一博公开解答，原题 4.16。}
\begin{solve}
用$X$表示一个家庭的男孩个数，$Y$表示一个家庭的小孩数.    
\[P(X = k) = P(Y = n)P(X = k|Y = n)\]
    （i）$k=0$时，根据全概率公式
    \[\begin{aligned}P(X &= 0) \\&= 1 - \alpha \frac{p}{q} + {\sum\limits_{n = 1}^\infty  {\alpha {p^n}\left( {\frac{1}{2}} \right)} ^n} \\&= 1 - \alpha \frac{p}{q} + \alpha \frac{{\frac{p}{2}}}{{1 - \frac{p}{2}}} \\&= 1 - \frac{{\alpha p}}{q} + \frac{{\alpha p}}{{2 - p}}\end{aligned}\]
    （ii）$k\geqslant 1$时，根据全概率公式
    \[\begin{aligned}P(X &= k) \\&= \alpha {\sum\limits_{n = k}^\infty  {C_n^k\left( {\frac{p}{2}} \right)} ^n} \\&= \alpha {\sum\limits_{m = 0}^\infty  {C_{m + k}^k\left( {\frac{p}{2}} \right)} ^{m + k}} \\&= \alpha {\left( {1 - \frac{p}{2}} \right)^{ - (k + 1)}}{\left( {\frac{p}{2}} \right)^k}\end{aligned}\]
\end{solve}

\subsection*{训练题~\ref{q:Y4.22} 连续变量的离散随机缩放}\label{a:Y4.22}{\small 题目见第~\pageref{q:Y4.22} 页；来源：闫一博公开解答，原题 4.22。}
\begin{solve}给定 $Y=a_i\ne0$ 时，$Z=a_iX$ 的密度为 $|a_i|^{-1}f(z/a_i)$。对任意可测集合 $D$，全概率公式给出\[P(Z\in D)=\sum_i p_i\int_D\frac1{|a_i|}f(z/a_i)\,dz.\]被积项非负，交换求和与积分合法，故\[h(z)=\sum_i\frac{p_i}{|a_i|}f(z/a_i).\]它的积分等于 $\sum_i p_i=1$，因此是概率密度。绝对值保证负缩放时密度仍非负。\end{solve}

\subsection*{训练题~\ref{q:Y4.25} 博克斯—穆勒正态生成}\label{a:Y4.25}{\small 题目见第~\pageref{q:Y4.25} 页；来源：闫一博公开解答，原题 4.25。}
\begin{solve}对 $r>0$，$P(R>r)=P(U<e^{-r^2/2})=e^{-r^2/2}$，所以 $f_R(r)=re^{-r^2/2}$。又 $\Theta=2\pi V\sim U(0,2\pi)$，且它与 $R$ 独立。极坐标变换的雅可比绝对值为 $r$，除原点与角度接缝这些零面积集合外可逐区域使用变换公式，得到\[f_{X,Y}(x,y)=\frac1{2\pi}e^{-(x^2+y^2)/2}.\]在整个平面上，该密度等于两个标准正态密度的乘积，所以 $X,Y$ 独立且均为标准正态。这里没有把 $\arctan(y/x)$ 当成全平面唯一角度；角度应按象限选取。\end{solve}

\subsection*{训练题~\ref{q:Y4.30} 指数变量的极值与比值}\label{a:Y4.30}{\small 题目见第~\pageref{q:Y4.30} 页；来源：闫一博公开解答，原题 4.30。}
\begin{solve}当 $u>0$，最小值的生存函数为 $e^{-(\lambda+\mu)u}$，所以\[f_U(u)=(\lambda+\mu)e^{-(\lambda+\mu)u}.\]最大值的分布函数为 $(1-e^{-\lambda v})(1-e^{-\mu v})$，求导得\[f_V(v)=\lambda e^{-\lambda v}+\mu e^{-\mu v}-(\lambda+\mu)e^{-(\lambda+\mu)v},\quad v>0.\]对于比值 $W=X/Y$，采用 $x=wy$，逆雅可比为 $y$，故\[\begin{aligned}f_W(w)&=\lambda\mu\int_0^\infty ye^{-(\lambda w+\mu)y}\,dy\\&=\frac{\lambda\mu}{(\lambda w+\mu)^2},\quad w>0.\end{aligned}\]各密度在负半轴为零。若写最小值与最大值的联合密度，其支撑必须是 $0<u<v$，不能写成整个正象限。\end{solve}

\subsection*{训练题~\ref{q:Y4.34} 平方和与平方比的变换}\label{a:Y4.34}{\small 题目见第~\pageref{q:Y4.34} 页；来源：闫一博公开解答，原题 4.34。}
\begin{solve}正象限内变换一一对应于 $u>0,v>0$，逆变换为\[x=\sqrt{\frac{uv}{1+v}},\qquad y=\sqrt{\frac{u}{1+v}}.\]计算逆变换的雅可比绝对值，得\[\left|\frac{\partial(x,y)}{\partial(u,v)}\right|=\frac{1}{4\sqrt v(1+v)}.\]由于原密度为 $e^{-x-y}$，故\[f_{U,V}(u,v)=\frac{\exp\{-\sqrt{u/(1+v)}(1+\sqrt v)\}}{4\sqrt v(1+v)},\quad u,v>0.\]在其余区域上密度为零。雅可比乘在密度上，而不是乘在单点事件概率上。\end{solve}

\subsection*{训练题~\ref{q:Y4.36} 线性变换的密度公式}\label{a:Y4.36}{\small 题目见第~\pageref{q:Y4.36} 页；来源：闫一博公开解答，原题 4.36。}
\begin{solve}可逆性保证每个 $(u,v)$ 对应唯一逆像 $(a_1u+b_1v,c_1u+d_1v)$，逆变换导数矩阵恒为 $A^{-1}$。因此雅可比绝对值为 $|\det A^{-1}|=1/|\det A|$。对任意可测集合 $D$，变量替换公式给出\[\begin{aligned}P((U,V)\in D)&=\iint_{A^{-1}D}f(x,y)\,dx\,dy\\&=\iint_D\frac{f(a_1u+b_1v,c_1u+d_1v)}{|\det A|}\,du\,dv.\end{aligned}\]这证明了密度公式。当 $A$ 不可逆时，像可能集中在直线上，一般不存在二维勒贝格密度，不能沿用此式。\end{solve}

\subsection*{训练题~\ref{q:A4.4} 常数与任意变量独立}\label{a:A4.4}{\small 题目见第~\pageref{q:A4.4} 页；来源：ayhe123 公开讲义，原题 4.4。}
\begin{proof}
    $\because$ 当 $b\geq a$ 时有 $P(a\leq b)=1,P(a\leq b,Y\leq y)=P(Y\leq y)$, 当 $b<a$ 时有 $P(a\leq b)=P(a\leq b,Y\leq y)=0$, $\therefore\forall b,y$ 都有
    \[P(a\leq b,Y\leq y)=P(Y\leq y)P(a\leq b).\]

    由分布函数的乘积分解判据得 $a$ 与 $Y$ 独立.
\end{proof}

\subsection*{训练题~\ref{q:A4.5} 联合分布函数的右连续性}\label{a:A4.5}{\small 题目见第~\pageref{q:A4.5} 页；来源：ayhe123 公开讲义，原题 4.5。}
\begin{solve}固定 $x$。若 $y_1\le y_2$，事件 $\{X\le x,Y\le y_1\}$ 包含于 $\{X\le x,Y\le y_2\}$，由概率的单调性，$F(x,y_1)\le F(x,y_2)$。对任意序列 $y_n\downarrow y$，定义 $A_n=\{X\le x,Y\le y_n\}$，则 $A_n\downarrow\{X\le x,Y\le y\}$。概率的从上连续性给出\[F(x,y_n)=P(A_n)\longrightarrow P\left(\bigcap_nA_n\right)=F(x,y).\]故关于第二个坐标右连续。这个性质不要求随机向量具有密度。\end{solve}

\subsection*{训练题~\ref{q:A4.11} 曲边区域上的联合密度}\label{a:A4.11}{\small 题目见第~\pageref{q:A4.11} 页；来源：ayhe123 公开讲义，原题 4.11。}
\begin{solve}
    由
    \[\begin{aligned}\int_{-\infty}^\infty\int_{-\infty}^\infty f(x,y)\mathrm{d}x\mathrm{d}y&=\iint\limits_{x^2<y<1}ax^2y\mathrm{d}x\mathrm{d}y\\&=1\end{aligned}\]

    得 $a=\dfrac{21}{4}$. 有
    \[\begin{aligned}f_X(x)&=\int_{-\infty}^\infty f(x,y)\mathrm{d}y\\&=\int_{x^2}^1\dfrac{21}{4}x^2y\mathrm{d}y\\&=\dfrac{21}{8}x^2(1-x^2),\quad -1<x<1,\end{aligned}\]
    \[\begin{aligned}f_Y(y)&=\int_{-\infty}^\infty f(x,y)\mathrm{d}x\\&=\int_{-\sqrt{y}}^{\sqrt{y}}\dfrac{21}{4}x^2y\mathrm{d}x\\&=\dfrac{7}{2}y^2\sqrt{y},\quad 0<y<1.\end{aligned}\]

    $\because f_X(x)f_Y(y)\neq f(x,y)$, $\therefore X,Y$ 不独立.
\end{solve}

\subsection*{训练题~\ref{q:A4.14} 指数寿命的比较}\label{a:A4.14}{\small 题目见第~\pageref{q:A4.14} 页；来源：ayhe123 公开讲义，原题 4.14。}
\begin{solve}条件于 $Y=y>0$，指数变量的生存概率为 $P(X>y)=e^{-\lambda y}$。独立性保证给定 $Y$ 不改变 $X$ 的分布，故\[\begin{aligned}P(X>Y)&=\int_0^\infty e^{-\lambda y}\mu e^{-\mu y}\,dy\\&=\frac\mu{\lambda+\mu}.\end{aligned}\]可用对称性检查：当两率相同时概率为 $1/2$；当 $\lambda$ 很大、$X$ 往往很短时，该概率趋于零。\end{solve}

\subsection*{训练题~\ref{q:A4.17} 不独立变量的平方独立}\label{a:A4.17}{\small 题目见第~\pageref{q:A4.17} 页；来源：ayhe123 公开讲义，原题 4.17。}
\begin{solve}对任一变量积分，得到两个边缘密度均为区间 $(-1,1)$ 上的 $1/2$。联合密度不几乎处处等于边缘密度之积 $1/4$，所以不独立。另一方面，当 $0<u,v<1$ 时，\[\begin{aligned}P(X^2\le u,Y^2\le v)&=\int_{-\sqrt u}^{\sqrt u}\int_{-\sqrt v}^{\sqrt v}\frac{1+xy}{4}\,dy\,dx\\&=\sqrt{uv}.\end{aligned}\]含 $xy$ 的积分因对称性为零；而边缘分布函数分别为 $\sqrt u,\sqrt v$。在区间外按分布函数的 $0$、$1$ 延拓仍成立，因此平方变量独立。\end{solve}

\subsection*{训练题~\ref{q:A4.18} 曲线下均匀点的投影}\label{a:A4.18}{\small 题目见第~\pageref{q:A4.18} 页；来源：ayhe123 公开讲义，原题 4.18。}
\begin{solve}
    设 $x_0\in\mathbb{R}$. 有
    \begin{align*}
        P(X\leq x_0) & =\iint_D\dfrac{1}{m(D)}I[X\leq x_0]\mathrm{d}x\mathrm{d}y \\
        & =\int_{-\infty}^{x_0}\mathrm{d}x\int_0^{g(x)}\dfrac{1}{m(D)}\mathrm{d}y \\
        & =\int_{-\infty}^{x_0}\dfrac{g(x)}{m(D)}\mathrm{d}x.
    \end{align*}

    $\therefore X$ 有概率密度
    \[f_X(x)=\dfrac{g(x)}{m(D)},\]

    其中 $m(D)=\int_{-\infty}^\infty g(x)\mathrm{d}x$.
\end{solve}

\subsection*{训练题~\ref{q:A4.19} 离散与连续变量之差}\label{a:A4.19}{\small 题目见第~\pageref{q:A4.19} 页；来源：ayhe123 公开讲义，原题 4.19。}
\begin{solve}记 $b_k=\binom nkp^k(1-p)^{n-k}$。给定 $X=k$ 后，$Z=Y-k$ 是平移后的指数变量，所以\[F_Z(z)=\sum_{k=0}^n b_k(1-e^{-\lambda(z+k)})\mathbf1_{\{z\ge-k\}},\]\[f_Z(z)=\sum_{k=0}^n b_k\lambda e^{-\lambda(z+k)}\mathbf1_{\{z>-k\}}.\]第二式几乎处处成立；各端点处密度值可任意指定。它是有限个平移指数密度的混合，不能在求导后丢掉对 $k$ 的求和。\end{solve}

\subsection*{训练题~\ref{q:A4.23} 求和变换后的边缘积分}\label{a:A4.23}{\small 题目见第~\pageref{q:A4.23} 页；来源：ayhe123 公开讲义，原题 4.23。}
\begin{solve}令 $Z=X+Y,W=X$，逆变换为 $x=w,y=z-w$，雅可比绝对值为一。原支撑 $x,y>0$ 化成 $z>0,0<w<z$，因此\[f_{Z,W}(z,w)=\frac12ze^{-z}\mathbf1_{\{z>0,0<w<z\}}.\]求边缘必须对 $w$ 积分：\[f_Z(z)=\int_0^z\frac12ze^{-z}\,dw=\frac12z^2e^{-z},\quad z>0.\]这正是形状为三、率为一的伽马密度。原解漏掉边缘积分中的区间长度 $z$，其结果不能归一化，本题据此校正。\end{solve}

\subsection*{训练题~\ref{q:A4.27} 独立正态的线性组合}\label{a:A4.27}{\small 题目见第~\pageref{q:A4.27} 页；来源：ayhe123 公开讲义，原题 4.27。}
\begin{solve}正态分布的特征函数为 $\exp(i\mu t-\sigma^2t^2/2)$。独立性使联合乘积的期望分解，故\[\varphi_Y(t)=\exp\left\{it\left(b+\sum_ja_j\mu_j\right)-\frac{t^2}{2}\sum_ja_j^2\sigma_j^2\right\}.\]因此 $Y$ 服从均值 $b+\sum_ja_j\mu_j$、方差 $\sum_ja_j^2\sigma_j^2$ 的正态分布。若该方差为零，则分布退化在其均值处。仅知道每个边缘分布是正态而不假定独立或联合正态，不足以推出线性组合正态。\end{solve}

\subsection*{训练题~\ref{q:A4.28} 二维正态的条件分布}\label{a:A4.28}{\small 题目见第~\pageref{q:A4.28} 页；来源：ayhe123 公开讲义，原题 4.28。}
\begin{solve}把标准化二次型配方：\[z_1^2-2\rho z_1z_2+z_2^2=(1-\rho^2)z_1^2+(z_2-\rho z_1)^2.\]将联合密度除以 $X_1$ 的 $N(\mu_1,\sigma_1^2)$ 边缘密度，第一项指数和相应归一化因子抵消，剩下关于 $x_2$ 的正态密度。因此\[X_2\mid X_1=x_1\sim N(m(x_1),(1-\rho^2)\sigma_2^2),\]\[m(x_1)=\mu_2+\rho\frac{\sigma_2}{\sigma_1}(x_1-\mu_1).\]条件期望对 $x_1$ 线性，条件方差不依赖 $x_1$。连续条件分布通过密度比定义，不是以正概率事件 $X_1=x_1$ 作初等条件概率除法。\end{solve}

\subsection*{训练题~\ref{q:A4.32} 尾分布幂次与竞争概率}\label{a:A4.32}{\small 题目见第~\pageref{q:A4.32} 页；来源：ayhe123 公开讲义，原题 4.32。}
\begin{solve}由题设 $F_Y(y)=1-[1-F_X(y)]^\beta$。独立性允许条件于 $X$，所以\[P(X\ge Y)=E[F_Y(X)]=1-E[1-F_X(X)]^\beta.\]由于 $X$ 的分布连续，概率积分变换给出 $F_X(X)\sim U(0,1)$，即使 $F_X$ 不处处严格增加也成立。因此\[P(X\ge Y)=\int_0^1[1-(1-u)^\beta]du=\frac{\beta}{1+\beta}.\]该方法避免声称“有密度便处处可导”；绝对连续分布函数一般只在几乎处处具有导数。\end{solve}

\clearpage\section*{数学期望、方差与协方差：详解}\addcontentsline{toc}{section}{数学期望、方差与协方差：详解}

\subsection*{训练题~\ref{q:A5.5} 广告信息之间的等待}\label{a:A5.5}{\small 题目见第~\pageref{q:A5.5} 页；来源：ayhe123 公开讲义，原题 5.5。}
\begin{solve}
    设 $X$ 为相邻的两次广告短信中不是广告短信的数目, 则 $P(X=k)=p^k(1-p)$. 有
    \begin{align*}
        EX & =\sum\limits_{k=0}^\infty kP(X=k) \\
        & =\sum\limits_{k=1}^\infty kp^k(1-p) \\
        & =p(1-p)\sum\limits_{k=1}^\infty kp^{k-1} \\
        & =p(1-p)\left(\sum\limits_{k=1}^\infty\dfrac{\partial}{\partial x}x^k\right)\Bigg|_{x=p} \\
        & =p(1-p)\left(\dfrac{\partial}{\partial x}\sum\limits_{k=0}^\infty x^k\right)\Bigg|_{x=p} \\
        & =p(1-p)\left(\dfrac{\partial}{\partial x}\dfrac{1}{1-x}\right)\Bigg|_{x=p} \\
        & =p(1-p)\cdot\dfrac{1}{(1-p)^2}=\dfrac{p}{1-p}.
    \end{align*}

    把 $p=2\%$ 代入得 $EX=49$.
\end{solve}

\subsection*{训练题~\ref{q:A5.8} 标准正态半径的期望}\label{a:A5.8}{\small 题目见第~\pageref{q:A5.8} 页；来源：ayhe123 公开讲义，原题 5.8。}
\begin{solve}二维联合密度为 $(2\pi)^{-1}\exp[-(x^2+y^2)/2]$。用极坐标 $x=r\cos\theta,y=r\sin\theta$，雅可比为 $r$，积分角度得\[f_R(r)=r e^{-r^2/2},\qquad r>0.\]这称为瑞利分布。由于 $R$ 非负，期望积分的范围必须是正半轴：\[ER=\int_0^\infty r^2e^{-r^2/2}\,dr=\sqrt{\frac\pi2}.\]最后一步可分部积分，将它化为 $\int_0^\infty e^{-r^2/2}\,dr$。此外 $ER^2=EX^2+EY^2=2$，所以 $\operatorname{Var}(R)=2-\pi/2$。\end{solve}

\subsection*{训练题~\ref{q:A5.9} 交换对称与比例期望}\label{a:A5.9}{\small 题目见第~\pageref{q:A5.9} 页；来源：ayhe123 公开讲义，原题 5.9。}
\begin{solve}
    设
    \[Y_i=\dfrac{X_i}{X_1+X_2+\cdots+X_n},\quad (n=1,2,\cdots,n).\]
    
    $\because X_1,X_2,\cdots,X_n$ 独立同分布, $\therefore Y_1,Y_2,\cdots,Y_n$ 同分布. $\because P(X_i>0)=1,\therefore P(Y_i>0)=1$. 有
    \[1=E\dfrac{X_1+X_2+\cdots+X_n}{X_1+X_2+\cdots+X_n}=EY_1+EY_2+\cdots+EY_n=nEY_1,\]
    \[E\dfrac{X_1+X_2+\cdots+X_k}{X_1+X_2+\cdots+X_n}=EY_1+EY_2+\cdots+EY_k=kEY_1,\]

    $\therefore$
    \[E\dfrac{X_1+X_2+\cdots+X_k}{X_1+X_2+\cdots+X_n}=\dfrac{k}{n}.\]
\end{solve}

\subsection*{训练题~\ref{q:A5.12} 离散极值的尾和公式}\label{a:A5.12}{\small 题目见第~\pageref{q:A5.12} 页；来源：ayhe123 公开讲义，原题 5.12。}
\begin{solve}对任意非负整数值变量 $Z$，恒等式 $Z=\sum_{k\ge1}\mathbf1_{\{Z\ge k\}}$ 与非负积分交换定理给出 $EZ=\sum_{k\ge1}P(Z\ge k)$，允许两边都为无穷。最小值不小于 $k$ 当且仅当所有 $X_i\ge k$，由独立性其概率为 $v_k^n$。最大值小于 $k$ 当且仅当所有 $X_i<k$，概率为 $u_k^n$，所以最大值不小于 $k$ 的概率为 $1-u_k^n$。分别代入尾和公式即得结论。\end{solve}

\subsection*{训练题~\ref{q:A5.14} 曲边密度的归一化勘误}\label{a:A5.14}{\small 题目见第~\pageref{q:A5.14} 页；来源：ayhe123 公开讲义，原题 5.14。}
\begin{solve}固定 $x>1$，积分区域为 $1/x<y<x$，所以\[\int_1^\infty\!\int_{1/x}^{x}\frac{3}{2x^3y^2}\,dy\,dx=\frac32\int_1^\infty(x^{-2}-x^{-4})\,dx=1.\]由非负函数的积分交换定理，\[\begin{aligned}EY&=\frac32\int_1^\infty x^{-3}\!\int_{1/x}^x y^{-1}\,dy\,dx\\&=3\int_1^\infty x^{-3}\log x\,dx=\frac34.\end{aligned}\]末式分部积分即可得到。原资料题面中的 $3/(2x^2y^3)$ 在该区域上不能归一化，故不能直接作为概率密度。\end{solve}

\subsection*{训练题~\ref{q:A5.17} 全部指数寿命结束的等待}\label{a:A5.17}{\small 题目见第~\pageref{q:A5.17} 页；来源：ayhe123 公开讲义，原题 5.17。}
\begin{solve}
    记 $Y=\max(X_1,\cdots,X_5)$, 则
    \begin{align*}
        P(Y>y) & =1-P(Y\leq y) \\
        & =1-P(X_1\leq y,X_2\leq y,\cdots,X_5\leq y) \\
        & =1-P(X_1\leq y)P(X_2\leq y)\cdots P(X_5\leq y) \\
        & =1-(1-e^{-\lambda y})^5 \\
        & =1-\left(1+\sum\limits_{i=1}^5C_5^i(-e^{-\lambda y})^i\right) \\
        & =\sum\limits_{i=1}^5C_5^i(-1)^{i-1}e^{-\lambda iy}.
    \end{align*}

    由相应的分布或期望公式可得
    \begin{align*}
        EY & =\int_0^\infty P(Y>y)\mathrm{d}y \\
        & =\int_0^\infty\sum\limits_{i=1}^5C_5^i(-1)^{i-1}e^{-\lambda iy}\mathrm{d}y \\
        & =\sum\limits_{i=1}^5C_5^i(-1)^{i-1}\int_0^\infty e^{-\lambda iy}\mathrm{d}y \\
        & =\sum\limits_{i=1}^5C_5^i\dfrac{1}{\lambda i}(-1)^{i-1}=\dfrac{137}{60\lambda}.
    \end{align*}
\end{solve}

\subsection*{训练题~\ref{q:A5.20} 样本均值与方差的无偏性}\label{a:A5.20}{\small 题目见第~\pageref{q:A5.20} 页；来源：ayhe123 公开讲义，原题 5.20。}
\begin{solve}记共同均值为 $\mu$、方差为 $\sigma^2$。期望的线性性给出 $E\widehat\mu=\mu$。平方和恒等式为\[\sum_{j=1}^n(X_j-\widehat\mu)^2=\sum_{j=1}^n(X_j-\mu)^2-n(\widehat\mu-\mu)^2.\]因独立性使 $\operatorname{Var}(\widehat\mu)=\sigma^2/n$，两边取期望得到 $(n-1)\sigma^2$。除以 $n-1$ 即有 $E\widehat\sigma^2=\sigma^2$。有限二阶矩是这里取方差与平方和期望的前提。\end{solve}

\subsection*{训练题~\ref{q:A5.22} 最优逆方差加权}\label{a:A5.22}{\small 题目见第~\pageref{q:A5.22} 页；来源：ayhe123 公开讲义，原题 5.22。}
\begin{solve}独立性给出方差 $\sum_iw_i^2\sigma_i^2$。由柯西不等式，\[1=\left(\sum_i(w_i\sigma_i)\sigma_i^{-1}\right)^2\le\left(\sum_iw_i^2\sigma_i^2\right)\left(\sum_i\sigma_i^{-2}\right).\]等号成立当且仅当 $w_i\sigma_i=c/\sigma_i$。结合权重和约束得到\[w_i=\frac{\sigma_i^{-2}}{\sum_j\sigma_j^{-2}},\qquad\min\operatorname{Var}=\frac{1}{\sum_j\sigma_j^{-2}}.\]方差较小的观测获得较大权重；等方差时退化为算术平均。\end{solve}

\subsection*{训练题~\ref{q:A5.25} 非负变量的尾积分}\label{a:A5.25}{\small 题目见第~\pageref{q:A5.25} 页；来源：ayhe123 公开讲义，原题 5.25。}
\begin{proof}
    (1) 有 $\forall j,P(X>x|X=x_j)=I[x<x_j]$. 由全概率公式得
    \[P(X>x)=\sum\limits_{j=1}^\infty P(X>x|X=x_j)P(X=x_j)=\sum\limits_{j=1}^\infty p_jI[x<x_j].\]

    (2) 由 (1) 得
    \[\int_0^\infty P(X>x)\mathrm{d}x=\int_0^\infty\sum\limits_{j=1}^\infty p_jI[x<x_j]\mathrm{d}x.\]

    $\because\sum\limits_{j=1}^\infty p_jI[x<x_j]\leq\sum\limits_{j=1}^\infty p_j=1$, 由 Weierstrass 判别法得 $\sum\limits_{j=1}^\infty p_jI[x<x_j]$ 在 $(0,\infty)$ 上一致收敛, $\therefore$ 求和号和积分号可以交换次序. 有
    \begin{align*}
        \int_0^\infty\sum\limits_{j=1}^\infty p_jI[x<x_j]\mathrm{d}x & =\sum\limits_{j=1}^\infty p_j\int_0^\infty I[x<x_j]\mathrm{d}x \\
        & =\sum\limits_{j=1}^\infty p_j\int_0^{x_j}\mathrm{d}x \\
        & =\sum\limits_{j=1}^\infty p_jx_j=EX.\qedhere
    \end{align*}
\end{proof}

\subsection*{训练题~\ref{q:A5.26} 和与差的相关系数}\label{a:A5.26}{\small 题目见第~\pageref{q:A5.26} 页；来源：ayhe123 公开讲义，原题 5.26。}
\begin{solve}由协方差的双线性与对称性，\[\operatorname{Cov}(X+Y,X-Y)=\operatorname{Var}(X)-\operatorname{Var}(Y).\]又因为 $\operatorname{Cov}(X,Y)=0$，\[\operatorname{Var}(X+Y)=\operatorname{Var}(X-Y)=\operatorname{Var}(X)+\operatorname{Var}(Y)>0.\]将协方差除以两个标准差之积即可得到题设公式。尤其当原变量方差相同时，和与差不相关；只有在额外联合正态等条件下，才能进一步断言它们独立。\end{solve}

\subsection*{训练题~\ref{q:A5.28} 正方形密度的协方差}\label{a:A5.28}{\small 题目见第~\pageref{q:A5.28} 页；来源：ayhe123 公开讲义，原题 5.28。}
\begin{solve}
    有
    \[f_X(x)=\int_{-\infty}^\infty f(x,y)\mathrm{d}y=\int_0^1(x+y)\mathrm{d}y=x+\dfrac{1}{2},\quad x\in(0,1),\]
    \[EX=\int_{-\infty}^\infty xf_X(x)\mathrm{d}x=\dfrac{7}{12}.\]

    对称地, 有 $EY=\dfrac{7}{12}$.

    \[\begin{aligned}E(XY)&=\iint_{\mathbb{R}^2}xyf(x,y)\mathrm{d}x\mathrm{d}y\\&=\int_0^1\mathrm{d}x\int_0^1xy(x+y)\mathrm{d}y\\&=\dfrac{1}{3}.\end{aligned}\]

    $\therefore$
    \[\operatorname{Cov}(X,Y)=E(XY)-EXEY=-\dfrac{1}{144}.\]
\end{solve}

\subsection*{训练题~\ref{q:A5.30} 正态尺寸差与装配}\label{a:A5.30}{\small 题目见第~\pageref{q:A5.30} 页；来源：ayhe123 公开讲义，原题 5.30。}
\begin{solve}活塞能装入等价于 $X<Y$。独立正态变量之差仍为正态，\[Y-X\sim N(0.10,0.02^2+0.02^2).\]故所求概率为\[P(Y-X>0)=\Phi\left(\frac{0.10}{0.02\sqrt2}\right)\approx0.9997965.\]直接对差变量标准化可避免二维积分。方差相加依赖两个尺寸独立；若存在同批次相关性，应在差的方差中减去两倍协方差。\end{solve}

\subsection*{训练题~\ref{q:A5.32} 不相关但不独立的反例}\label{a:A5.32}{\small 题目见第~\pageref{q:A5.32} 页；来源：ayhe123 公开讲义，原题 5.32。}
\begin{solve}由柯西不等式，$E|X|<\infty$；由密度对称性，$EX=0$，且奇函数 $x|x|$ 的积分为零。因此\[\operatorname{Cov}(X,|X|)=E[X|X|]-EX\,E|X|=0.\]若 $X$ 与 $|X|$ 独立，则对任意 $a$，事件 $B_a=\{|X|\le a\}$ 既由 $X$ 决定又由 $|X|$ 决定，故 $P(B_a)=P(B_a)^2$。于是每个这样的概率只能为 $0$ 或 $1$，迫使 $|X|$ 几乎处处为常数。这与 $X$ 具有密度相矛盾。因此二者不独立。\end{solve}

\subsection*{训练题~\ref{q:A5.33} 中心正态的高阶矩}\label{a:A5.33}{\small 题目见第~\pageref{q:A5.33} 页；来源：ayhe123 公开讲义，原题 5.33。}
\begin{solve}正态分布的所有绝对矩有限，故可用对称性处理奇次幂，得到 $EX^{2m+1}=0$。设正态密度为 $f$，则 $f^{\prime}(x)=-xf(x)/\sigma^2$。对偶次矩分部积分，边界项因高斯衰减为零：\[\begin{aligned}EX^{2m}&=-\sigma^2\int_{-\infty}^{\infty}x^{2m-1}f^{\prime}(x)dx\\&=(2m-1)\sigma^2EX^{2m-2}.\end{aligned}\]从 $EX^0=1$ 递推得 $EX^{2m}=\sigma^{2m}(2m-1)!!$，其中双阶乘为所有正奇数 $1\cdot3\cdots(2m-1)$ 的乘积。\end{solve}

\subsection*{训练题~\ref{q:A5.35} 圆周与圆内均匀模型}\label{a:A5.35}{\small 题目见第~\pageref{q:A5.35} 页；来源：ayhe123 公开讲义，原题 5.35。}
\begin{solve}取 $\Theta\sim U(0,2\pi)$，令 $X=R\cos\Theta,Y=R\sin\Theta$。三角函数在整周期上的积分给出\[EX=EY=0,\quad EX^2=EY^2=\frac{R^2}{2},\quad EXY=0.\]故两个方差均为 $R^2/2$，协方差为零。但它们不独立：事件 $|X|>\sqrt3R/2$ 与 $|Y|>R/2$ 都有正概率，而其交集由 $X^2+Y^2=R^2$ 可知为空。圆内按面积均匀取点是另一种模型，此时两个方差均为 $R^2/4$。原资料解答混用了这两种均匀分布，本题已校订。\end{solve}

\clearpage\section*{条件期望与生成函数：详解}\addcontentsline{toc}{section}{条件期望与生成函数：详解}

\subsection*{训练题~\ref{q:A6.1} 随机组数下的条件二项}\label{a:A6.1}{\small 题目见第~\pageref{q:A6.1} 页；来源：ayhe123 公开讲义，原题 6.1。}
\begin{solve}
    $N_j$ 表示成功概率为 $0.6$ 的 $8$ 次独立重复试验中成功的次数. $\because N_1,N_2,\cdots$ 独立, $\therefore T|\{M=k\}=N_1+\cdots+N_k$ 表示成功概率为 $0.6$ 的 $8k$ 次独立重复试验中成功的次数, 服从二项分布 $\mathcal{B}(8k,0.6)$.
\end{solve}

\subsection*{训练题~\ref{q:A6.3} 几何个指数等待之和}\label{a:A6.3}{\small 题目见第~\pageref{q:A6.3} 页；来源：ayhe123 公开讲义，原题 6.3。}
\begin{solve}
    记 $q=1-p$。由独立指数变量之和的伽马分布公式, $W|\{N=n\}$ 有概率密度
    \[f_W(w;n)=\dfrac{\lambda^n}{(n-1)!}e^{-\lambda w}w^{n-1},\quad w>0.\]

    $\therefore$
    \[E(W|N=n)=\int_0^\infty wf_W(w;n)\mathrm{d}w=\int_0^\infty\dfrac{\lambda^n}{(n-1)!}e^{-\lambda w}w^n\mathrm{d}w.\]

    令 $t=\lambda w$, 得
    \begin{align*}
        \int_0^\infty\dfrac{\lambda^n}{(n-1)!}e^{-\lambda w}w^n\mathrm{d}w & =\dfrac{1}{\lambda(n-1)!}\int_0^\infty e^{-t}t^n\mathrm{d}t \\
        & =\dfrac{1}{\lambda(n-1)!}\Gamma(n+1)=\dfrac{n}{\lambda}.
    \end{align*}

    由相应的分布或期望公式可得
    \begin{align*}
        EW & =\sum\limits_{n=1}^\infty E(W|N=n)P(N=n) \\
        & =\sum\limits_{n=1}^\infty\dfrac{n}{\lambda}q^{n-1}p \\
        & =\dfrac{p}{\lambda}\left(\dfrac{\mathrm{d}}{\mathrm{d}x}\sum\limits_{n=0}^\infty x^n\right)\Bigg|_{x=q} \\
        & =\dfrac{p}{\lambda}\left(\dfrac{1}{(1-x)^2}\right)\Bigg|_{x=q}=\dfrac{1}{\lambda p}.
    \end{align*}
\end{solve}

\subsection*{训练题~\ref{q:A6.4} 复合泊松等待的均值方差}\label{a:A6.4}{\small 题目见第~\pageref{q:A6.4} 页；来源：ayhe123 公开讲义，原题 6.4。}
\begin{solve}
    设 $X_1,X_2,\cdots$ 是在第 $i$ 个红灯处的等待时间, 则 $X|\{N=n\}=X_1+X_2+\cdots+X_n$, $X_i$ 独立同分布, $X_1\sim\mathcal{U}(0,2)$, $N$ 与 $\{X_i\}$ 独立. 有
    \[E(X|N=n)=EX_1+EX_2+\cdots+EX_n=nEX_1=n,\]
    \begin{align*}
        E(X^2|N=n) & =E(X_1+X_2+\cdots+X_n)^2 \\
        & =E\left(X_1^2+X_2^2+\cdots+X_n^2+\sum\limits_{i=1}^n\sum\limits_{\substack{j=1\\j\neq i}}^nX_iX_j\right) \\
        & =nEX_1^2+\sum\limits_{i=1}^n\sum\limits_{\substack{j=1\\j\neq i}}^nEX_iEX_j \\
        & =\dfrac{4n}{3}+n(n-1).
    \end{align*}

    由相应的分布或期望公式可得
    \[EX=\sum\limits_{n=0}^\infty P(N=n)E(X|N=n)=\lambda,\]
    \begin{align*}
        EX^2 & =\sum\limits_{n=0}^\infty P(N=n)E(X^2|N=n) \\
        & =\sum\limits_{n=0}^\infty\left(\dfrac{4n}{3}+n(n-1)\right)\dfrac{\lambda^n}{n!}e^{-\lambda}  \\
        & =\dfrac{4}{3}\sum\limits_{n=1}^\infty\dfrac{\lambda^n}{(n-1)!}e^{-\lambda}+\sum\limits_{n=2}^\infty\dfrac{\lambda^n}{(n-2)!}e^{-\lambda}  \\
        & =\dfrac{4}{3}\lambda+\lambda^2.
    \end{align*}

    $\therefore$
    \[\operatorname{Var}(X)=EX^2-(EX)^2=\dfrac{4}{3}\lambda.\]
\end{solve}

\subsection*{训练题~\ref{q:A6.9} 条件指数与条件伽马}\label{a:A6.9}{\small 题目见第~\pageref{q:A6.9} 页；来源：ayhe123 公开讲义，原题 6.9。}
\begin{solve}
    $f_{X|Y}(x|y),f_{Y|X}(y|x)$ 可以由 $f(x,y)$ 除以一个归一化因子得到. 有
    \[f_{X|Y}(x|y)=\dfrac{f(x,y)}{\int_0^\infty f(x,y)\mathrm{d}x}=\dfrac{ye^{-y(x+1)}}{e^{-y}}=ye^{-yx},\]
    \[f_{Y|X}(y|x)=\dfrac{f(x,y)}{\int_0^\infty f(x,y)\mathrm{d}y}=(x+1)^2ye^{-y(x+1)}.\]
\end{solve}

\subsection*{训练题~\ref{q:A6.14} 威布尔寿命的剩余等待}\label{a:A6.14}{\small 题目见第~\pageref{q:A6.14} 页；来源：ayhe123 公开讲义，原题 6.14。}
\begin{solve}
    由相应的分布或期望公式可得
    \begin{align*}
        E(Y|Y>t) & =\dfrac{E(YI[Y>t])}{P(Y>t)} \\
        & =\dfrac{\int_t^\infty yf(y)\mathrm{d}y}{\int_t^\infty f(y)\mathrm{d}y} \\
        & =\dfrac{ab\int_t^\infty y^{1/2}e^{-ay^{1/2}}\mathrm{d}y}{ab\int_t^\infty y^{-1/2}e^{-ay^{1/2}}\mathrm{d}y} \\
        & =\dfrac{\int_t^\infty y^{1/2}e^{-ay^{1/2}}\mathrm{d}y}{\int_t^\infty y^{-1/2}e^{-ay^{1/2}}\mathrm{d}y}.
    \end{align*}

    令 $u=ay^{1/2}$, 则 $y=(u/a)^2,\mathrm{d}y=\dfrac{2}{a^2}u\mathrm{d}u$. 有
    \begin{align*}
        \dfrac{\int_t^\infty y^{1/2}e^{-ay^{1/2}}\mathrm{d}y}{\int_t^\infty y^{-1/2}e^{-ay^{1/2}}\mathrm{d}y} & =\dfrac{\int_{a\sqrt{t}}^\infty(u/a)e^{-u}(2/a^2)u\mathrm{d}u}{\int_{a\sqrt{t}}^\infty(u/a)^{-1}e^{-u}(2/a^2)u\mathrm{d}u} \\
        & =\dfrac{\int_{a\sqrt{t}}^\infty u^2e^{-u}\mathrm{d}u}{a^2\int_{a\sqrt{t}}^\infty e^{-u}\mathrm{d}u} \\
        & =\dfrac{e^{-a\sqrt{t}}(2+2a\sqrt{t}+a^2t)}{a^2e^{-a\sqrt{t}}} \\
        & =\dfrac{2+2a\sqrt{t}+a^2t}{a^2}.
    \end{align*}

    $\therefore$
    \[E(Y-t|Y>t)=E(Y|Y>t)-E(t|Y>t)=\dfrac{2+2a\sqrt{t}+a^2t}{a^2}-t.\]
\end{solve}

\subsection*{训练题~\ref{q:A6.15} 泊松观察后的伽马后验}\label{a:A6.15}{\small 题目见第~\pageref{q:A6.15} 页；来源：ayhe123 公开讲义，原题 6.15。}
\begin{solve}先验密度与似然分别为\[\pi(y)=\frac{\beta^\alpha}{\Gamma(\alpha)}y^{\alpha-1}e^{-\beta y},\quad P(N=n\mid y)=\frac{y^ne^{-y}}{n!}.\]由贝叶斯公式，其乘积中依赖 $y$ 的部分为 $y^{\alpha+n-1}e^{-(\beta+1)y}$。将其归一化得到\[Y\mid N=n\sim\operatorname{Gamma}(\alpha+n,\beta+1).\]所以 $E(Y\mid N=n)=(\alpha+n)/(\beta+1)$。须始终区分率参数与尺度参数；这里的伽马均值是形状除以率。\end{solve}

\subsection*{训练题~\ref{q:A6.17} 七个接口的第三次失效}\label{a:A6.17}{\small 题目见第~\pageref{q:A6.17} 页；来源：ayhe123 公开讲义，原题 6.17。}
\begin{solve}系统在第三个接口失效时停止工作，所以 $Y=X_{(3)}$ 是七个寿命的第三小值。对 $y\ge0$，记 $q=1-e^{-\lambda y}$，则在时刻 $y$ 已失效数量服从 $\operatorname{Bin}(7,q)$。故\[P(Y>y)=\sum_{k=0}^{2}\binom7k(1-e^{-\lambda y})^k e^{-(7-k)\lambda y}.\]利用顺序统计量密度或对其求导，得到\[f_Y(y)=105\lambda(1-e^{-\lambda y})^2e^{-5\lambda y},\quad y>0.\]指数中的系数是 $5$；原解写为 $2$ 会导致积分不等于一。此处还可由独立指数间隔求 $EY=(1/7+1/6+1/5)/\lambda$。\end{solve}

\subsection*{训练题~\ref{q:A6.20} 指数顺序统计量的间隔}\label{a:A6.20}{\small 题目见第~\pageref{q:A6.20} 页；来源：ayhe123 公开讲义，原题 6.20。}
\begin{solve}原样本连续，所以并列值的概率为零。顺序统计量在 $0<x_{(1)}<\cdots<x_{(n)}$ 上的联合密度为\[n!\lambda^n\exp\left(-\lambda\sum_{i=1}^nx_{(i)}\right).\]令 $x_{(i)}=\sum_{j=1}^iy_j$，变换矩阵为下三角且对角元均为一，故雅可比为一；新支撑为所有 $y_j>0$。又\[\sum_{i=1}^nx_{(i)}=\sum_{j=1}^n(n-j+1)y_j.\]因此联合密度为\[\prod_{j=1}^n\left[(n-j+1)\lambda e^{-(n-j+1)\lambda y_j}\right].\]每个因子是一个完整指数密度，支撑为直积，故各间隔独立，率分别为 $n\lambda,(n-1)\lambda,\ldots,\lambda$。\end{solve}

\subsection*{训练题~\ref{q:A6.21} 仿射变换的概率母函数}\label{a:A6.21}{\small 题目见第~\pageref{q:A6.21} 页；来源：ayhe123 公开讲义，原题 6.21。}
\begin{solve}
    $Y$ 的取值为 $\{ja+b|j\in\mathbb{R}\}$. 有
    \begin{align*}
        g(s) & =\sum\limits_{j=0}^\infty s^jP(X=j) \\
        & =\sum\limits_{j=0}^\infty(s^{1/a})^{ja}P(Y=ja+b) \\
        & =(s^{1/a})^{-b}\sum\limits_{j=0}^\infty(s^{1/a})^{ja+b}P(Y=ja+b) \\
        & =s^{-b/a}h(s^{1/a}).
    \end{align*}

    $\therefore$
    \[h(s^{1/a})=g(s)s^{b/a}\Rightarrow h(s)=g(s^a)s^b.\]
\end{solve}

\subsection*{训练题~\ref{q:A6.24} 两组硬币成功数之差}\label{a:A6.24}{\small 题目见第~\pageref{q:A6.24} 页；来源：ayhe123 公开讲义，原题 6.24。}
\begin{solve}设甲乙的正面次数为 $H_1,H_2$，彼此独立且均为 $\operatorname{Bin}(n,1/2)$。由于 $n-H_2$ 仍服从相同二项分布且与 $H_1$ 独立，\[H_1-H_2+n=H_1+(n-H_2)\sim\operatorname{Bin}(2n,1/2).\]因此对整数 $-n\le k\le n$，\[P(H_1-H_2=k)=\frac{\binom{2n}{n+k}}{2^{2n}},\]范围外概率为零。将负幂生成函数转化为这个非负二项和，可以避免移位时错写幂次或求和下标。\end{solve}

\subsection*{训练题~\ref{q:A6.29} 均匀和的特征函数极限}\label{a:A6.29}{\small 题目见第~\pageref{q:A6.29} 页；来源：ayhe123 公开讲义，原题 6.29。}
\begin{solve}直接积分可得 $\varphi_X(t)=\sin t/t$，$t\ne0$，并在 $t=0$ 连续定义为 $1$。独立性及缩放性质给出\[\varphi_{Y_n}(t)=\left[\frac{\sin(t\sqrt{3/n})}{t\sqrt{3/n}}\right]^n.\]对每个固定 $t$，由 $\sin u/u=1-u^2/6+O(u^4)$，括号内为 $1-t^2/(2n)+O(n^{-2})$。所以\[\varphi_{Y_n}(t)\to e^{-t^2/2}.\]极限是标准正态特征函数且在零点连续，由莱维连续性定理，$Y_n\xrightarrow dN(0,1)$。余项展开是对 $n\to\infty$ 而 $t$ 固定进行的。\end{solve}

\subsection*{训练题~\ref{q:A6.30} 指数和的特征函数极限}\label{a:A6.30}{\small 题目见第~\pageref{q:A6.30} 页；来源：ayhe123 公开讲义，原题 6.30。}
\begin{solve}指数变量的特征函数为\[\varphi_X(t)=\int_0^\infty\lambda e^{-(\lambda-it)x}dx=\frac{\lambda}{\lambda-it}.\]独立性及中心化缩放给出\[\varphi_{Y_n}(t)=\left[\frac{e^{-it/\sqrt n}}{1-it/\sqrt n}\right]^n.\]对固定 $t$，在原点附近使用复对数的连续分支，\[\log\varphi_{Y_n}(t)=n[-it/\sqrt n-\log(1-it/\sqrt n)]\to-\frac{t^2}2.\]线性项相消，二次项给出上述极限，余项为 $O(n^{-1/2})$。莱维连续性定理得到 $Y_n\xrightarrow dN(0,1)$。\end{solve}

\subsection*{训练题~\ref{q:A6.34} 二项分布的泊松极限}\label{a:A6.34}{\small 题目见第~\pageref{q:A6.34} 页；来源：ayhe123 公开讲义，原题 6.34。}
\begin{solve}对固定实数 $t$，特征函数为\[\varphi_n(t)=[1+p_n(e^{it}-1)]^n.\]由 $np_n\to\lambda$ 可知 $p_n\to0$、$np_n^2\to0$。在 $1$ 的邻域展开对数，得到\[\log\varphi_n(t)=np_n(e^{it}-1)+O(np_n^2)\to\lambda(e^{it}-1).\]因此特征函数逐点趋于 $\exp\{\lambda(e^{it}-1)\}$，后者是泊松分布的特征函数且在原点连续。由莱维连续性定理得到结论。条件 $np_n\to\lambda$ 不等同于每个 $n$ 都满足 $p_n=\lambda/n$。\end{solve}

\subsection*{训练题~\ref{q:A6.35} 给定总和后的部分和}\label{a:A6.35}{\small 题目见第~\pageref{q:A6.35} 页；来源：ayhe123 公开讲义，原题 6.35。}
\begin{solve}联合分布在交换坐标后不变，而 $S_n$ 也不变，因此各个 $E(X_j\mid S_n)$ 几乎处处相等。条件期望的线性性给出\[\sum_{j=1}^nE(X_j\mid S_n)=E(S_n\mid S_n)=S_n.\]于是每一项等于 $S_n/n$，再对前 $k$ 项求和即得结论。可积性保证这些条件期望存在；一般不需要条件密度存在。\end{solve}

\clearpage\section*{大数定律与中心极限定理：详解}\addcontentsline{toc}{section}{大数定律与中心极限定理：详解}

\subsection*{训练题~\ref{q:A7.2} 截尾期望的概率上界}\label{a:A7.2}{\small 题目见第~\pageref{q:A7.2} 页；来源：ayhe123 公开讲义，原题 7.2。}
\begin{proof}
    $\because X$ 有概率密度 $f(x)$, $\therefore$
    \begin{align*}
        MP(X\geq M) & =M\int_M^\infty f(x)\mathrm{d}x \\
        & =\int_M^\infty Mf(x)\mathrm{d}x \\
        & \leq\int_M^\infty xf(x)\mathrm{d}x.\qedhere
    \end{align*}
\end{proof}

\subsection*{训练题~\ref{q:A7.3} 柯西不等式与尾概率}\label{a:A7.3}{\small 题目见第~\pageref{q:A7.3} 页；来源：ayhe123 公开讲义，原题 7.3。}
\begin{solve}记 $A=\{X>M\}$，则 $P(A)=1-F(M)$。由非负性与柯西不等式，\[MP(A)\le E[X\mathbf1_A]\le\sqrt{EX^2\,P(A)}.\]除以 $M>0$ 即得\[P(X>M)\le\frac{\sqrt{(1-F(M))EX^2}}M.\]若 $P(A)=0$，结论仍成立。这里不需要密度存在，也不需要把弱不等式误写成对所有情形都严格的不等式。\end{solve}

\subsection*{训练题~\ref{q:A7.4} 最小值的概率界}\label{a:A7.4}{\small 题目见第~\pageref{q:A7.4} 页；来源：ayhe123 公开讲义，原题 7.4。}
\begin{solve}令 $A=\{X_1\ge M\}$。因 $X_1\mathbf1_A\ge M\mathbf1_A$，有 $MP(A)\le E(X_1\mathbf1_A)$。再用柯西不等式得到\[MP(A)\le\sqrt{EX_1^2\,P(A)}.\]将这一不等式两侧除以 $M$ 并取 $n$ 次幂，同时使用独立性 $P(\min_iX_i\ge M)=P(A)^n$，即得结论。当 $P(A)=0$ 时两侧均为零，无须进行除以 $P(A)$ 的操作。\end{solve}

\subsection*{训练题~\ref{q:A7.5} 伽马矩与切比雪夫界}\label{a:A7.5}{\small 题目见第~\pageref{q:A7.5} 页；来源：ayhe123 公开讲义，原题 7.5。}
\begin{solve}该密度为形状 $m+1$、率一的伽马分布，其中 $m$ 为非负整数。由伽马积分，\[EX=m+1,\quad EX^2=(m+1)(m+2),\]\[\operatorname{Var}(X)=m+1.\]切比雪夫不等式给出\[P(|X-(m+1)|\ge m+1)\le\frac1{m+1}.\]取补事件，得到 $P(0<X<2(m+1))\ge m/(m+1)$。这是下界，最后一步不能写成概率恰好等于 $m/(m+1)$。\end{solve}

\subsection*{训练题~\ref{q:A7.10} 三角函数样本均值之比}\label{a:A7.10}{\small 题目见第~\pageref{q:A7.10} 页；来源：ayhe123 公开讲义，原题 7.10。}
\begin{solve}由均匀密度积分，\[E\sin X_1=\frac2\pi\int_0^{\pi/2}\sin x\,dx=\frac2\pi,\]\[E\cos X_1=\frac2\pi\int_0^{\pi/2}\cos x\,dx=\frac2\pi.\]两个函数均有界，所以强大数定律分别适用；在两个概率一事件的交集上，两组样本平均都收敛到 $2/\pi$。分母极限非零，除法在该点连续，故比值 $Y_n$ 几乎必然收敛到 $1$。\end{solve}

\subsection*{训练题~\ref{q:A7.14} 比例调查的样本量}\label{a:A7.14}{\small 题目见第~\pageref{q:A7.14} 页；来源：ayhe123 公开讲义，原题 7.14。}
\begin{solve}样本比例 $\widehat p$ 的方差为 $p(1-p)/n\le1/(4n)$。正态近似要求\[z_{0.995}\frac{1}{2\sqrt n}\le0.01.\]采用 $z_{0.995}=2.575829\ldots$，得到 $n\ge16588$；若将分位数上调为 $2.58$，得到 $n\ge16641$。这是正态近似下的规划值，不是所有 $p$ 的有限样本精确覆盖保证。若仅用切比雪夫不等式，\[P(|\widehat p-p|\ge0.01)\le\frac{1}{4n(0.01)^2},\]取 $n\ge250000$ 可给出至少 $99\%$ 的严格但较粗的保证。\end{solve}

\subsection*{训练题~\ref{q:A7.15} 两组样本均值差的近似}\label{a:A7.15}{\small 题目见第~\pageref{q:A7.15} 页；来源：ayhe123 公开讲义，原题 7.15。}
\begin{solve}由中心极限定理，各组标准化样本均值渐近正态；两组独立保证两项独立。于是\[\overline X-\overline Y\ \dot\sim\ N\left(\mu_1-\mu_2,\frac{\sigma_1^2}{n}+\frac{\sigma_2^2}{m}\right).\]更严格地，把差除以其真实标准差；独立性和两个分量的正态极限使该标准化差趋于 $N(0,1)$。若两总体本来就是正态分布，上式为任意 $n,m$ 下的精确分布。样本容量不同，第二项分母应为 $m$。\end{solve}

\subsection*{训练题~\ref{q:A7.20} 泊松概率和的二分极限}\label{a:A7.20}{\small 题目见第~\pageref{q:A7.20} 页；来源：ayhe123 公开讲义，原题 7.20。}
\begin{solve}令 $Z_1,Z_2,\ldots$ 独立同分布，且 $Z_j\sim\operatorname{Pois}(1)$，则 $S_n=\sum_{j=1}^nZ_j\sim\operatorname{Pois}(n)$。由于 $EZ_j=\operatorname{Var}(Z_j)=1$，中心极限定理给出\[\frac{S_n-n}{\sqrt n}\xrightarrow{d}N(0,1).\]第一个和等于 $P((S_n-n)/\sqrt n\le0)$，故极限为 $\Phi(0)=1/2$。第二个和的标准化阈值为 $-1/\sqrt n\to0$。对任意 $\varepsilon>0$，当 $n$ 足够大时可将该概率夹在阈值 $-\varepsilon$ 与 $0$ 的概率之间，再令 $\varepsilon\downarrow0$，仍得到 $1/2$。\end{solve}

\subsection*{训练题~\ref{q:A7.21} 缩放收敛与紧性}\label{a:A7.21}{\small 题目见第~\pageref{q:A7.21} 页；来源：ayhe123 公开讲义，原题 7.21。}
\begin{solve}令 $Z_n=\sqrt n\xi_n$。依分布收敛于一个有限实值分布意味着分布族渐近紧：对任意 $\delta>0$，可选 $M$ 使 $\limsup_nP(|Z_n|>M)<\delta$，因为标准正态两侧尾概率可任意小，且 $\pm M$ 为连续点。固定 $\varepsilon>0$，当 $n$ 足够大时 $\varepsilon\sqrt n>M$，故\[P(|\xi_n|>\varepsilon)=P(|Z_n|>\varepsilon\sqrt n)\le P(|Z_n|>M).\]取上极限后再令 $\delta\downarrow0$，得到依概率收敛到零。不能直接把随 $n$ 改变的阈值代入只对固定点成立的分布函数收敛式。\end{solve}

\subsection*{训练题~\ref{q:A7.22} 不同参数二项和的极限}\label{a:A7.22}{\small 题目见第~\pageref{q:A7.22} 页；来源：ayhe123 公开讲义，原题 7.22。}
\begin{solve}可在等分布的模型中将每个 $X_j$ 写成 $j$ 个参数为 $p_j$ 的独立伯努利变量之和，并令不同组也相互独立。记 $B_n^2=\operatorname{Var}(S_n)$，则\[B_n^2=\sum_{j=1}^n jp_j(1-p_j)\ge0.09\frac{n(n+1)}2\longrightarrow\infty.\]每个中心化伯努利变量的绝对值不超过 $1$。对任意 $\varepsilon>0$，最终有 $\varepsilon B_n>1$，所以林德伯格条件中的所有截尾二阶矩均为零。林德伯格中心极限定理遂给出所需结论。使用这一形式的定理时，须把二项变量拆成有界的独立小项，不能只因原变量独立便直接断言结论。\end{solve}

\subsection*{训练题~\ref{q:A7.25} 小扰动不改变分布极限}\label{a:A7.25}{\small 题目见第~\pageref{q:A7.25} 页；来源：ayhe123 公开讲义，原题 7.25。}
\begin{solve}固定 $x\in\mathbb R$ 和 $\varepsilon>0$。由事件包含关系，\[P(\xi_n\le x-\varepsilon)-P(|\eta_n|>\varepsilon)\le P(\xi_n+\eta_n\le x),\]\[P(\xi_n+\eta_n\le x)\le P(\xi_n\le x+\varepsilon)+P(|\eta_n|>\varepsilon).\]令 $n\to\infty$，外侧概率分别趋于 $\Phi(x-\varepsilon)$ 与 $\Phi(x+\varepsilon)$；再由 $\Phi$ 连续令 $\varepsilon\downarrow0$，得到中间分布函数趋于 $\Phi(x)$。因此依分布收敛成立，整个证明不需要 $\xi_n$ 与 $\eta_n$ 独立。\end{solve}

\clearpage\section*{样本与抽样分布：详解}\addcontentsline{toc}{section}{样本与抽样分布：详解}

\subsection*{训练题~\ref{q:M-F-C4} 统计量与未知参数}\label{a:M-F-C4}{\small 题目见第~\pageref{q:M-F-C4} 页；来源：MIT 期末试卷，概念题 4。}
\begin{solve}统计量是样本的可测函数，其定义不能包含未知参数。因此 $\overline X$ 和 $S$ 是统计量；$1/\lambda$ 是总体均值，包含未知参数；$\overline X-1/\lambda$ 虽然也是随机量，但不属于本模型的统计量。注意“分布依赖参数”和“计算公式包含未知参数”不同：$\overline X$ 的抽样分布依赖 $\lambda$，仍然是统计量。观测前用大写表示随机统计量，观测后用 $\overline x,s$ 表示已经计算出的数值。\end{solve}

\subsection*{训练题~\ref{q:B8.1} 样本均值的均值与方差}\label{a:B8.1}{\small 题目见第~\pageref{q:B8.1} 页；来源：编订补充。}
\begin{solve}由期望的线性性，不需要独立性即可得到 $E\overline X=\mu$。独立性使交叉协方差为零，所以\[\operatorname{Var}(\overline X)=\frac{1}{n^2}\sum_{j=1}^n\sigma^2=\frac{\sigma^2}{n},\qquad\operatorname{Cov}(X_i,\overline X)=\frac{\sigma^2}{n}.\]若各观测仍有相同均值、方差但存在相关性，则\[\operatorname{Var}(\overline X)=\frac{\sigma^2}{n}+\frac{2}{n^2}\sum_{i<j}\operatorname{Cov}(X_i,X_j).\]正相关可能显著降低平均的降噪效果，故不能把独立样本的标准误机械用于相关数据。\end{solve}

\subsection*{训练题~\ref{q:B8.2} 样本方差的无偏性}\label{a:B8.2}{\small 题目见第~\pageref{q:B8.2} 页；来源：编订补充。}
\begin{solve}先作恒等变形：\[\sum_i(X_i-\overline X)^2=\sum_i(X_i-\mu)^2-n(\overline X-\mu)^2.\]右侧两项的期望分别为 $n\sigma^2$ 和 $n\operatorname{Var}(\overline X)=\sigma^2$，所以平方和的期望为 $(n-1)\sigma^2$，除以 $n-1$ 即得无偏性。若分母为 $n$，估计量的期望为 $(n-1)\sigma^2/n$，偏差为 $-\sigma^2/n$。自由度损失来自用同一组数据估计了均值；上述无偏性不要求总体正态。\end{solve}

\subsection*{训练题~\ref{q:B8.3} 正态样本均值和残差的独立性}\label{a:B8.3}{\small 题目见第~\pageref{q:B8.3} 页；来源：编订补充。}
\begin{solve}样本向量为联合正态向量，任何线性变换后的向量仍联合正态。对每个 $i$，\[\operatorname{Cov}(\overline X,X_i-\overline X)=\sigma^2/n-\sigma^2/n=0.\]因此样本均值与残差向量之间的交叉协方差矩阵为零。联合正态向量中，零交叉协方差意味着两个子向量独立；这个结论也允许残差协方差矩阵奇异。$S^2$ 是残差向量的可测函数，因此也与 $\overline X$ 独立。一般总体下，不相关并不推出独立。\end{solve}

\subsection*{训练题~\ref{q:B8.4} 卡方分布及自由度}\label{a:B8.4}{\small 题目见第~\pageref{q:B8.4} 页；来源：编订补充。}
\begin{solve}标准化后记 $Z_i=(X_i-\mu)/\sigma$。取正交矩阵，其第一行为 $n^{-1/2}(1,\ldots,1)$。变换后得到独立标准正态变量 $U_1,\ldots,U_n$，其中 $U_1=\sqrt n(\overline X-\mu)/\sigma$。正交变换保持平方和，故\[\frac{(n-1)S^2}{\sigma^2}=\sum_{i=1}^nZ_i^2-U_1^2=\sum_{j=2}^nU_j^2.\]最后一项是 $n-1$ 个独立标准正态变量平方和，按定义服从 $\chi^2_{n-1}$。估计未知均值去掉了沿全一向量的一个自由方向。\end{solve}

\subsection*{训练题~\ref{q:B8.5} 学生 t 统计量}\label{a:B8.5}{\small 题目见第~\pageref{q:B8.5} 页；来源：编订补充。}
\begin{solve}令 $Z=\sqrt n(\overline X-\mu)/\sigma$，$V=(n-1)S^2/\sigma^2$。正态抽样理论给出 $Z\sim N(0,1)$、$V\sim\chi^2_{n-1}$，且两者独立。于是\[T=\frac{Z}{\sqrt{V/(n-1)}}\sim t_{n-1}.\]分母含有估计方差的不确定性，使 $t$ 分布尾部较厚。$n$ 很大时它趋近标准正态；有限 $n$ 时若使用正态分位数，均值区间会偏窄。非正态总体下，这一有限样本精确结论通常不成立。\end{solve}

\subsection*{训练题~\ref{q:B8.6} 方差比与 F 分布}\label{a:B8.6}{\small 题目见第~\pageref{q:B8.6} 页；来源：编订补充。}
\begin{solve}两个独立的标准化平方和分别服从 $\chi^2_{n-1}$ 和 $\chi^2_{m-1}$。因此\[\frac{S_1^2/\sigma_1^2}{S_2^2/\sigma_2^2}\sim F_{n-1,m-1}.\]若 $W\sim F_{a,b}$，则按定义交换分子分母可得 $1/W\sim F_{b,a}$。记 $F_{a,b;q}$ 为下侧累积概率 $q$ 的分位数，则连续性给出\[F_{a,b;q}=\frac{1}{F_{b,a;1-q}}.\]使用表格时必须核对分子和分母自由度顺序，也要核对采用的是下侧还是上侧分位数。\end{solve}

\subsection*{训练题~\ref{q:B8.7} 顺序统计量的分布}\label{a:B8.7}{\small 题目见第~\pageref{q:B8.7} 页；来源：编订补充。}
\begin{solve}由独立性，所有观测不超过 $x$ 的概率为 $F(x)^n$，所以 $F_{X_{(n)}}(x)=F(x)^n$；类似地 $P(X_{(1)}>x)=[1-F(x)]^n$。第 $k$ 个顺序统计量落在 $x$ 附近，要求恰有 $k-1$ 个观测小于 $x$，一个落在微小区间，其余大于 $x$。求极限或对二项尾和求导，得到\[f_{(k)}(x)=\frac{n!}{(k-1)!(n-k)!}F(x)^{k-1}[1-F(x)]^{n-k}f(x).\]连续性使相等的并列观测概率为零。离散总体中不能直接使用这个密度公式。\end{solve}

\subsection*{训练题~\ref{q:B8.8} 均匀样本最大值}\label{a:B8.8}{\small 题目见第~\pageref{q:B8.8} 页；来源：编订补充。}
\begin{solve}$M$ 的密度为 $nx^{n-1}/\theta^n$，$0<x<\theta$。对 $r>0$ 积分得到\[E(M^r)=\frac{n}{n+r}\theta^r.\]故 $EM=n\theta/(n+1)$，而\[\operatorname{Var}(M)=\left[\frac{n}{n+2}-\frac{n^2}{(n+1)^2}\right]\theta^2=\frac{n\theta^2}{(n+1)^2(n+2)}.\]将最大值乘以 $(n+1)/n$ 得到无偏估计量 $\widetilde\theta=(n+1)M/n$。注意它可能大于实际观测最大值，这不违反参数作为总体上界的含义。\end{solve}

\subsection*{训练题~\ref{q:B8.9} 指数样本的最小值与和}\label{a:B8.9}{\small 题目见第~\pageref{q:B8.9} 页；来源：编订补充。}
\begin{solve}对 $t\ge0$，$P(M>t)=\prod_i e^{-\lambda t}=e^{-n\lambda t}$，所以 $M\sim\operatorname{Exp}(n\lambda)$。拉普拉斯变换相乘给出 $T\sim\operatorname{Gamma}(n,\lambda)$，这里采用率参数。因此\[E(nM)=E\overline X=1/\lambda,\quad\operatorname{Var}(nM)=1/\lambda^2,\quad\operatorname{Var}(\overline X)=1/(n\lambda^2).\]两者都无偏估计总体均值，但 $nM$ 的方差不随样本容量缩小，$\overline X$ 则具有一致性。这说明无偏性本身不能保证估计质量。\end{solve}

\subsection*{训练题~\ref{q:M-F-C13} 中心极限定理的正确标准化}\label{a:M-F-C13}{\small 题目见第~\pageref{q:M-F-C13} 页；来源：MIT 期末试卷，概念题 13（补足有限正方差条件）。}
\begin{solve}依次为错、对、错、对。中心极限定理讨论的是加和后的标准化量，并不改变每个原始观测的分布。其精确渐近表述为\[\frac{\sqrt n(\overline X-\mu)}{\sigma}\xrightarrow{d}N(0,1).\]因此 $\overline X$ 可近似视为 $N(\mu,\sigma^2/n)$，而没有乘 $\sqrt n$ 的量方差是 $1/n$。这里的“近似”没有自动给出某个固定样本容量的误差保证，重尾或强偏态会使收敛较慢。\end{solve}

\clearpage\section*{矩估计、偏差与一致性：详解}\addcontentsline{toc}{section}{矩估计、偏差与一致性：详解}

\subsection*{训练题~\ref{q:B9.1} 伯努利矩估计}\label{a:B9.1}{\small 题目见第~\pageref{q:B9.1} 页；来源：编订补充。}
\begin{solve}总体一阶矩为 $EX=p$，以样本一阶矩替代得到 $\widehat p=\overline X$。于是\[E\widehat p=p,\qquad\operatorname{MSE}(\widehat p)=\operatorname{Var}(\widehat p)=\frac{p(1-p)}n.\]切比雪夫不等式给出对任意 $\varepsilon>0$，\[P(|\widehat p-p|>\varepsilon)\le\frac{p(1-p)}{n\varepsilon^2}\longrightarrow0.\]因此估计量依概率一致。矩估计是一个构造方法，无偏性和一致性则是需要分别验证的性质。\end{solve}

\subsection*{训练题~\ref{q:B9.2} 均匀分布的矩估计与极值估计}\label{a:B9.2}{\small 题目见第~\pageref{q:B9.2} 页；来源：编订补充。}
\begin{solve}由 $EX=\theta/2$ 得 $\widehat\theta_M=2\overline X$。它无偏且方差为 $4\theta^2/(12n)=\theta^2/(3n)$。修正最大值估计也无偏，其方差为\[\operatorname{Var}(\widetilde\theta)=\frac{\theta^2}{n(n+2)}.\]二者均方误差之比为 $(n+2)/3$，所以 $n>1$ 时修正最大值更好，$n=1$ 时相同。这里充分利用参数限制支撑集的信息，比只匹配均值获得更快的误差下降。\end{solve}

\subsection*{训练题~\ref{q:B9.3} 伽马分布的两个矩参数}\label{a:B9.3}{\small 题目见第~\pageref{q:B9.3} 页；来源：编订补充。}
\begin{solve}总体均值、方差分别为 $\alpha/\beta$ 和 $\alpha/\beta^2$。将它们等同于 $m_1,v_n$，解方程得\[\widehat\alpha=\frac{m_1^2}{v_n},\qquad\widehat\beta=\frac{m_1}{v_n}.\]有限样本需 $v_n>0$；当 $n\ge2$ 且总体连续时，样本全相等的概率为零。由大数定律，样本一阶、二阶原点矩收敛，进而 $m_1\to\alpha/\beta$、$v_n\to\alpha/\beta^2$；连续映射定理给出两个估计的一致性。比值形式一般并非无偏。\end{solve}

\subsection*{训练题~\ref{q:B9.4} 正态分布的矩估计}\label{a:B9.4}{\small 题目见第~\pageref{q:B9.4} 页；来源：编订补充。}
\begin{solve}理论矩为 $EX=\mu$、$EX^2=\mu^2+\sigma^2$。配对样本矩可得\[\widehat\mu=\overline X,\qquad\widehat{\sigma^2}=\frac1n\sum_iX_i^2-\overline X^2=\frac1n\sum_i(X_i-\overline X)^2.\]后者等于 $(n-1)S^2/n$，期望为 $(n-1)\sigma^2/n$，有向下偏差。矩法并不保证无偏。若改用 $S^2$，则得到无偏方差估计，但它不再是严格按照这两个原点矩方程得到的估计。\end{solve}

\subsection*{训练题~\ref{q:B9.5} 偏差与均方误差分解}\label{a:B9.5}{\small 题目见第~\pageref{q:B9.5} 页；来源：编订补充。}
\begin{solve}在 $T-\theta=(T-ET)+(ET-\theta)$ 中平方后取期望，交叉项因 $E(T-ET)=0$ 消失，得到分解。对 $a\overline X$，偏差为 $(a-1)\mu$，方差为 $a^2\sigma^2/n$，故\[\operatorname{MSE}_\mu(a\overline X)=a^2\sigma^2/n+(a-1)^2\mu^2.\]若 $0<a<1$ 且 $\mu$ 靠近零，缩小均值可能降低均方误差；但对较大的 $|\mu|$，偏差平方会占主导。不能只比较方差便断言一个有偏估计统一优于无偏估计。\end{solve}

\subsection*{训练题~\ref{q:B9.6} 样本方差的最小均方误差缩放}\label{a:B9.6}{\small 题目见第~\pageref{q:B9.6} 页；来源：编订补充。}
\begin{solve}因为 $(n-1)S^2/\sigma^2\sim\chi^2_{n-1}$，且卡方分布方差为自由度的两倍，故 $\operatorname{Var}(S^2)=2\sigma^4/(n-1)$。于是\[\operatorname{MSE}(aS^2)=\sigma^4\left[\frac{2a^2}{n-1}+(a-1)^2\right].\]对 $a$ 求导并令其为零，得到 $a=(n-1)/(n+1)$；二阶导数为正。最小均方误差为 $2\sigma^4/(n+1)$。这个最优性只针对所指定的常数倍估计量类，且依赖正态假设。\end{solve}

\subsection*{训练题~\ref{q:B9.7} 一致性与渐近正态性}\label{a:B9.7}{\small 题目见第~\pageref{q:B9.7} 页；来源：编订补充。}
\begin{solve}有限均值保证大数定律适用，故 $\overline X\xrightarrow{P}\lambda$。又由中心极限定理，$\sqrt n(\overline X-\lambda)\xrightarrow{d}N(0,\lambda)$。函数 $g(x)=\log x$ 在 $\lambda>0$ 处可微，导数为 $1/\lambda$，德尔塔方法给出\[\sqrt n(\log\overline X-\log\lambda)\xrightarrow{d}N(0,1/\lambda).\]全零样本概率为 $e^{-n\lambda}$，趋于零；可在该事件上将估计另定义为任意有限数，从而得到处处有定义的渐近结论。不能对有限样本直接计算 $\log0$。\end{solve}

\subsection*{训练题~\ref{q:B9.8} 样本分位数的直观比较}\label{a:B9.8}{\small 题目见第~\pageref{q:B9.8} 页；来源：编订补充。}
\begin{solve}两者期望均为 $\theta$，故都无偏；方差分别为 $\theta^2/(3n)$ 与 $\theta^2/3$。前者依概率趋于 $\theta$。后者的分布不随 $n$ 改变；例如对 $0<\varepsilon<\theta$，\[P(|2X_1-\theta|>\varepsilon)=1-\varepsilon/\theta>0,\]所以不一致。拥有越来越多观测，并不使一个始终只用第一个观测的估计量自动变好。此例还说明无偏与一致是不同性质。\end{solve}

\clearpage\section*{极大似然估计：详解}\addcontentsline{toc}{section}{极大似然估计：详解}

\subsection*{训练题~\ref{q:M-C10-1} 离散参数空间中的硬币似然}\label{a:M-C10-1}{\small 题目见第~\pageref{q:M-C10-1} 页；来源：MIT 第 10 讲，题 1。}
\begin{solve}似然为 $L(p)=\binom{80}{49}p^{49}(1-p)^{31}$。在三个候选点的值约为 $2.08\times10^{-7},0.0118,0.0545$，故离散参数空间中的极大似然估计为 $2/3$。连续参数空间内，\[\ell^{\prime}(p)=49/p-31/(1-p)=0\]给出 $\widehat p=49/80=0.6125$；二阶导数严格为负，且两个边界似然为零，所以这是全局最大值。最大值取决于允许的参数集合，而不只是由数据决定。\end{solve}

\subsection*{训练题~\ref{q:M-C10-2} 指数寿命的极大似然估计}\label{a:M-C10-2}{\small 题目见第~\pageref{q:M-C10-2} 页；来源：MIT 第 10 讲，题 2。}
\begin{solve}五个密度相乘得到 $L(\lambda)=\lambda^5e^{-13\lambda}$，所以 $\ell(\lambda)=5\log\lambda-13\lambda$。导数为 $5/\lambda-13$，二阶导数为 $-5/\lambda^2<0$，因而 $\widehat\lambda=5/13$。一般正观测下，\[\widehat\lambda=\frac{n}{\sum_i x_i}=\frac1{\overline x}.\]这里估计的是单位时间事件率；平均寿命的极大似然估计为 $\overline x$。率和平均寿命是互为倒数的两个参数，须避免混用。\end{solve}

\subsection*{训练题~\ref{q:M-E2-II1} 已知形状的伽马率估计}\label{a:M-E2-II1}{\small 题目见第~\pageref{q:M-E2-II1} 页；来源：MIT 第二次考试，第二部分题 1。}
\begin{solve}不含 $\beta$ 的因子记为 $C>0$，则 $L(\beta)=C\beta^{15}e^{-12\beta}$。对数似然为 $15\log\beta-12\beta+\log C$，故驻点满足 $15/\beta-12=0$，得到 $\widehat\beta=5/4$。严格负的二阶导数保证最大值。一般地，$\operatorname{Gamma}(\alpha,\beta)$ 样本的率参数估计为\[\widehat\beta=\frac{n\alpha}{\sum_iX_i}.\]这也等于已知形状时的矩估计；若形状未知，两个参数的似然方程不再具有同样的简单闭式。\end{solve}

\subsection*{训练题~\ref{q:B10.1} 正态模型的联合极大似然}\label{a:B10.1}{\small 题目见第~\pageref{q:B10.1} 页；来源：编订补充。}
\begin{solve}令 $v=\sigma^2>0$，对数似然为\[\ell(\mu,v)=-\frac n2\log(2\pi v)-\frac{1}{2v}\sum_i(x_i-\mu)^2.\]固定 $v$，平方和在 $\mu=\overline x$ 时最小。代回后对 $v$ 求导，得到\[\widehat\mu=\overline x,\qquad\widehat v=\frac1n\sum_i(x_i-\overline x)^2.\]当平方和为正时，该驻点是全局最大值。若所有观测相同，则沿 $\mu$ 等于共同值、$v\downarrow0$，似然无界，在 $v>0$ 的参数空间中没有有限极大值；连续正态样本容量至少二时此事件概率为零。\end{solve}

\subsection*{训练题~\ref{q:B10.2} 均匀端点与支撑集}\label{a:B10.2}{\small 题目见第~\pageref{q:B10.2} 页；来源：编订补充。}
\begin{solve}联合似然为\[L(\theta)=\theta^{-n}\mathbf1_{\{\theta\ge M\}},\qquad M=\max_i x_i.\]在可行域 $\theta\ge M$ 内似然随 $\theta$ 递减，所以最大值在边界 $\widehat\theta=M$。若忽略指示函数，就会错误地把参数推向零，即使这样会使观测落在支撑集之外。密度在端点的取值不影响概率分布，但可影响极大值是否达到；这里明确采用包含端点的密度版本。\end{solve}

\subsection*{训练题~\ref{q:B10.3} 几何分布的估计与边界}\label{a:B10.3}{\small 题目见第~\pageref{q:B10.3} 页；来源：编订补充。}
\begin{solve}第一种记法的似然为 $p^n(1-p)^{T-n}$，$T=\sum_i x_i\ge n$。若 $T>n$，对数导数 $n/p-(T-n)/(1-p)=0$，得到 $\widehat p=n/T=1/\overline x$。若 $T=n$，所有观测都是 $1$，最大值在 $p=1$。若记失败次数 $Y=X-1$，则 $\widehat p=1/(1+\overline Y)$。几何分布存在两种通行支撑约定，估计公式必须随之改变。\end{solve}

\subsection*{训练题~\ref{q:B10.4} 带删失的指数寿命}\label{a:B10.4}{\small 题目见第~\pageref{q:B10.4} 页；来源：编订补充。}
\begin{solve}已失效个体的贡献为 $\lambda e^{-\lambda y_i}$，未失效个体贡献为生存概率 $e^{-\lambda c}$。合并得\[L(\lambda)=\lambda^d\exp\{-\lambda\sum_i y_i\},\qquad d=\sum_i\delta_i.\]当 $d>0$ 时，对数似然求导得到 $\widehat\lambda=d/\sum_i y_i$，即失效次数除以累计观察时间。当 $d=0$ 时，似然随 $\lambda$ 递减，在严格正参数空间中只有 $\lambda\downarrow0$ 的上确界；若扩展模型允许零率，可记边界估计为零。不能把所有未失效寿命误记成恰好在 $c$ 失效。\end{solve}

\subsection*{训练题~\ref{q:B10.5} 极大似然的不变性}\label{a:B10.5}{\small 题目见第~\pageref{q:B10.5} 页；来源：编订补充。}
\begin{solve}记新参数 $\eta=g(\theta)$，新似然只是重新参数化：$\widetilde L(\eta)=L(g^{-1}(\eta))$。一一对应保持候选分布不变，因此最大值在 $\eta=g(\widehat\theta)$ 达到。指数分布中位数 $m$ 满足 $1-e^{-\lambda m}=1/2$，故 $m=\log2/\lambda$，极大似然估计为\[\widehat m=\frac{\log2}{\widehat\lambda}=\overline X\log2.\]对非一一映射也有适当的轮廓似然版本，但不能把“估计量变换”与“变换后仍无偏”混为一谈。\end{solve}

\subsection*{训练题~\ref{q:B10.6} 指数率估计的有限样本偏差}\label{a:B10.6}{\small 题目见第~\pageref{q:B10.6} 页；来源：编订补充。}
\begin{solve}将伽马密度积分可得，对 $n>1$，\[E(T^{-1})=\frac{\lambda^n}{\Gamma(n)}\int_0^\infty t^{n-2}e^{-\lambda t}\,dt=\frac{\lambda}{n-1}.\]故极大似然估计 $n/T$ 的期望为 $n\lambda/(n-1)$，向上有偏；无偏修正为 $(n-1)/T$。当 $n=1$ 时倒数期望发散，不能沿用此结论。若 $n>2$，还可得 $E(T^{-2})=\lambda^2/[(n-1)(n-2)]$，从而无偏修正的方差为 $\lambda^2/(n-2)$。\end{solve}

\clearpage\section*{充分统计量与无偏估计：详解}\addcontentsline{toc}{section}{充分统计量与无偏估计：详解}

\subsection*{训练题~\ref{q:B11.1} 伯努利样本的充分统计量}\label{a:B11.1}{\small 题目见第~\pageref{q:B11.1} 页；来源：编订补充。}
\begin{solve}联合概率为 $p^t(1-p)^{n-t}$，对所有含恰好 $t$ 个 $1$ 的向量相同。因为 $P(T=t)=\binom nt p^t(1-p)^{n-t}$，条件概率为\[P((X_1,\ldots,X_n)=x\mid T=t)=\binom nt^{-1}\]（当 $\sum x_i=t$），其他向量为零。它不含 $p$，这正说明给定 $T$ 后样本排列不再提供关于 $p$ 的信息，也可用因子分解定理证明充分性。\end{solve}

\subsection*{训练题~\ref{q:B11.2} 泊松样本的充分性与条件二项分布}\label{a:B11.2}{\small 题目见第~\pageref{q:B11.2} 页；来源：编订补充。}
\begin{solve}联合质量函数为 $e^{-n\lambda}\lambda^{\sum x_i}/\prod_i x_i!$，按因子分解定理 $T$ 充分。又 $X_1\sim\operatorname{Pois}(\lambda)$ 与 $T-X_1\sim\operatorname{Pois}((n-1)\lambda)$ 独立，故\[P(X_1=k\mid T=t)=\binom tk(1/n)^k(1-1/n)^{t-k}.\]因此条件分布为 $\operatorname{Bin}(t,1/n)$，不含 $\lambda$。这个条件分布还能用于将只使用一个观测的估计改进为使用全样本的信息。\end{solve}

\subsection*{训练题~\ref{q:B11.3} 拉奥—布莱克改进}\label{a:B11.3}{\small 题目见第~\pageref{q:B11.3} 页；来源：编订补充。}
\begin{solve}由条件二项分布，\[\widetilde U=E(U\mid T)=P(X_1=0\mid T)=(1-1/n)^T.\]当 $n=1,T=0$ 时约定 $0^0=1$。迭代期望保证 $E\widetilde U=EU=e^{-\lambda}$。全方差公式给出\[\operatorname{Var}(U)=E\{\operatorname{Var}(U\mid T)\}+\operatorname{Var}(\widetilde U),\]所以方差不会增大。由于二者均无偏，方差比较也就是均方误差比较。这个过程利用条件期望去掉在充分信息给定后仍存在的无关随机波动。\end{solve}

\subsection*{训练题~\ref{q:B11.4} 二项族的完全性}\label{a:B11.4}{\small 题目见第~\pageref{q:B11.4} 页；来源：编订补充。}
\begin{solve}假设等价于\[\sum_{t=0}^ng(t)\binom nt p^t(1-p)^{n-t}=0.\]除以正数 $(1-p)^n$，令 $z=p/(1-p)\in(0,\infty)$，得到一个在正半轴恒为零的多项式 $\sum_tg(t)\binom ntz^t$。多项式各系数必须为零，所以 $g(t)=0$。因此 $T$ 完全；加上充分性，莱曼—谢费定理说明任何作为 $T$ 的函数的无偏估计，都是对应参数函数的唯一几乎处处的均匀最小方差无偏估计。\end{solve}

\subsection*{训练题~\ref{q:B11.5} 估计平方概率}\label{a:B11.5}{\small 题目见第~\pageref{q:B11.5} 页；来源：编订补充。}
\begin{solve}总成功次数 $T$ 满足阶乘矩 $E[T(T-1)]=n(n-1)p^2$，故\[U=\frac{T(T-1)}{n(n-1)}\]无偏。它是完全充分统计量 $T$ 的函数，因此也是 $p^2$ 的均匀最小方差无偏估计。代入估计的期望为\[E(T/n)^2=p^2+\frac{p(1-p)}n,\]所以向上有偏。无偏估计可以理解为对所有不同观测对 $X_iX_j$ 求平均；使用不同观测避免将 $X_i^2=X_i$ 混入平方概率的估计。\end{solve}

\subsection*{训练题~\ref{q:B11.6} 伯努利模型的信息量与下界}\label{a:B11.6}{\small 题目见第~\pageref{q:B11.6} 页；来源：编订补充。}
\begin{solve}单个观测的对数似然导数为 $U_p=X/p-(1-X)/(1-p)=(X-p)/[p(1-p)]$。其期望为零，方差为 $1/[p(1-p)]$，所以 $n$ 个独立观测的信息量为\[I_n(p)=\frac{n}{p(1-p)}.\]在可交换微分与求和等正则条件下，无偏估计 $p$ 的方差至少为 $1/I_n(p)=p(1-p)/n$。样本比例恰好具有这个方差，因而在此正则模型中达到下界。边界 $p=0,1$ 不属于这里的正则参数区间。\end{solve}

\subsection*{训练题~\ref{q:B11.7} 信息下界并非总能达到}\label{a:B11.7}{\small 题目见第~\pageref{q:B11.7} 页；来源：编订补充。}
\begin{solve}单个对数似然为 $\log\lambda-\lambda X$，二阶导数为 $-1/\lambda^2$，故 $I_n(\lambda)=n/\lambda^2$。无偏估计率参数的方差下界为 $\lambda^2/n$。而 $(n-1)/T$ 的方差为 $\lambda^2/(n-2)$，严格大于下界。下界是所有符合条件无偏估计的共同限制，并不保证存在达到它的估计量。两者之比为 $n/(n-2)\to1$，显示渐近效率与有限样本恰好达到下界是不同概念。\end{solve}

\subsection*{训练题~\ref{q:B11.8} 非正则模型中的错误下界}\label{a:B11.8}{\small 题目见第~\pageref{q:B11.8} 页；来源：编订补充。}
\begin{solve}均匀模型的支撑集随 $\theta$ 改变。只在支撑内部求导会遗漏边界贡献，得分的期望也不是零，所以通常的正则克拉美—拉奥推导不适用。修正最大值 $U=(n+1)X_{(n)}/n$ 是无偏的，方差为\[\operatorname{Var}(U)=\frac{\theta^2}{n(n+2)}<\frac{\theta^2}{n}.\]这不违反任何有效定理，而是说明所声称的下界缺少正则假设。应用信息不等式之前，必须检查微分与积分交换以及参数相关支撑问题。\end{solve}

\clearpage\section*{贝叶斯估计与预测：详解}\addcontentsline{toc}{section}{贝叶斯估计与预测：详解}

\subsection*{训练题~\ref{q:M-E2-II2} 三类别观察的贝叶斯更新}\label{a:M-E2-II2}{\small 题目见第~\pageref{q:M-E2-II2} 页；来源：MIT 第二次考试，第二部分题 2。}
\begin{solve}该有序序列的似然为 $(\theta/2)^4(1-\theta)$。与先验相乘，关于 $\theta$ 的部分为 $\theta^6(1-\theta)$，故后验属于 $\operatorname{Beta}(7,2)$。因为\[\int_0^1\theta^6(1-\theta)\,d\theta=\frac17-\frac18=\frac1{56},\]后验密度为 $56\theta^6(1-\theta)$，后验均值为 $7/9$。若只记录三个类别的计数，似然多一个与参数无关的多项式系数，后验不变。\end{solve}

\subsection*{训练题~\ref{q:M-E2-II3} 先验赔率与后验赔率}\label{a:M-E2-II3}{\small 题目见第~\pageref{q:M-E2-II3} 页；来源：MIT 第二次考试，第二部分题 3（抽象分类表述）。}
\begin{solve}赔率 $1:100$ 表示 $P(H)/P(H^c)=1/100$，并不是 $P(H)=1/100$。阳性的似然比为 $1/0.1=10$，所以\[\frac{P(H\mid+)}{P(H^c\mid+)}=\frac1{100}\times10=\frac1{10}.\]因此后验赔率为 $1:10$，后验概率为 $1/11$。高灵敏度并不使一次阳性自动意味着很高的后验概率，先验比例仍参与计算。本题只是给定概率模型下的推断，不把筛选结果解释成确定标签。\end{solve}

\subsection*{训练题~\ref{q:M-F-C14} 贝塔与二项共轭}\label{a:M-F-C14}{\small 题目见第~\pageref{q:M-F-C14} 页；来源：MIT 期末试卷，概念题 14（扩展计算）。}
\begin{solve}先验核为 $\theta^4(1-\theta)^4$，似然核为 $\theta^{20}(1-\theta)^{10}$，所以后验为 $\operatorname{Beta}(25,15)$。于是\[E(\theta\mid D)=25/40=5/8,\qquad\operatorname{mode}(\theta\mid D)=24/38=12/19.\]下一次成功的后验预测概率是对未知参数积分：$P(X_{31}=1\mid D)=E(\theta\mid D)=5/8$。后验均值、众数和极大似然估计 $2/3$ 是不同的概念，有限样本下通常不相同。\end{solve}

\subsection*{训练题~\ref{q:B12.1} 平方损失下的贝叶斯估计}\label{a:B12.1}{\small 题目见第~\pageref{q:B12.1} 页；来源：编订补充。}
\begin{solve}在后验分布下作平方分解，\[E[(a-\Theta)^2\mid D]=\operatorname{Var}(\Theta\mid D)+[a-E(\Theta\mid D)]^2.\]第一项不随 $a$ 改变，第二项在后验均值处唯一最小。因此贝叶斯估计依赖损失函数。对绝对损失，在连续后验下求导得到 $2F_{\Theta\mid D}(a)-1$，其过零点为后验中位数；含原子的情形使用左右导数，任何后验中位数都是最优解。\end{solve}

\subsection*{训练题~\ref{q:B12.2} 伽马与泊松共轭的多次观察}\label{a:B12.2}{\small 题目见第~\pageref{q:B12.2} 页；来源：编订补充。}
\begin{solve}似然中依赖 $\lambda$ 的部分是 $\lambda^{\sum_i x_i}\exp[-\lambda\sum_i t_i]$。因此\[\lambda\mid D\sim\operatorname{Gamma}\left(a+\sum_i x_i,\ b+\sum_i t_i\right),\]后验均值为 $(a+\sum_i x_i)/(b+\sum_i t_i)$。若令 $T=\sum_i t_i$，可将其写成 $\frac{b}{b+T}(a/b)+\frac{T}{b+T}(\sum_i x_i/T)$，即先验均值与数据率的加权平均。计数增加形状，累计观察时间增加率；把次数与暴露时间都简单记成样本容量会出错。\end{solve}

\subsection*{训练题~\ref{q:B12.3} 正态均值的共轭更新}\label{a:B12.3}{\small 题目见第~\pageref{q:B12.3} 页；来源：编订补充。}
\begin{solve}将先验和似然指数中的二次项配方，后验精度等于精度之和：\[v_n=\left(\tau_0^{-2}+n\sigma^{-2}\right)^{-1},\quad m_n=v_n\left(\frac{m_0}{\tau_0^2}+\frac{n\overline x}{\sigma^2}\right).\]故 $\mu\mid D\sim N(m_n,v_n)$。下一次观测可写成 $X_{n+1}=\mu+\varepsilon$，其中给定数据后噪声独立且方差为 $\sigma^2$，所以预测分布为 $N(m_n,\sigma^2+v_n)$。预测方差比参数后验方差更大，因为还包含新观测本身的随机性。\end{solve}

\subsection*{训练题~\ref{q:M-P7-1} 主持人规则改变三门问题}\label{a:M-P7-1}{\small 题目见第~\pageref{q:M-P7-1} 页；来源：MIT 习题 7，题 1。}
\begin{solve}令 $D$ 为“打开 $B$ 且无奖品”。第一种规则下，$D$ 在奖品分别位于 $A,B,C$ 时的似然为 $(1/2,0,1)$。乘以相等先验并归一化，后验为 $(1/3,0,2/3)$，应换到 $C$。第二种规则的似然为 $(1/2,0,1/2)$，后验为 $(1/2,0,1/2)$，换与不换相同。混合规则下似然为 $(1/2,0,0.7+0.3/2)=(1/2,0,17/20)$，归一化得\[(P(A\mid D),P(B\mid D),P(C\mid D))=(10/27,0,17/27).\]仍应换门。观察表面相同，但观察是如何产生的决定了似然。\end{solve}

\subsection*{训练题~\ref{q:M-P7-2} 未知面数骰子的预测}\label{a:M-P7-2}{\small 题目见第~\pageref{q:M-P7-2} 页；来源：MIT 习题 7，题 2（选编）。}
\begin{solve}面数为 $4,6$ 时似然为零。对 $d\in\{8,12,20\}$，似然为 $d^{-n}$，所以\[P(D=d\mid7,\ldots,7)=\frac{d^{-n}}{8^{-n}+12^{-n}+20^{-n}}.\]下一次为 $7$ 的概率是条件概率的后验平均，等于\[\frac{8^{-(n+1)}+12^{-(n+1)}+20^{-(n+1)}}{8^{-n}+12^{-n}+20^{-n}}.\]将分子分母同乘 $8^n$ 可见后验集中到八面骰，预测概率趋于 $1/8$。连续看到同一结果不会使下一次结果的概率趋于一，因为给定骰子之后仍有投掷随机性。\end{solve}

\subsection*{训练题~\ref{q:M-P8-1fg} 显著性与后验概率的区别}\label{a:M-P8-1fg}{\small 题目见第~\pageref{q:M-P8-1fg} 页；来源：MIT 习题 8，题 1(f,g)（选编）。}
\begin{solve}两个点假设的似然比为 $r=0.55^{140}0.45^{110}/0.5^{250}$，所以后验为 $r/(1+r)\approx0.8524$。连续均匀先验给出 $p\mid D\sim\operatorname{Beta}(141,111)$，于是\[P(p>0.5\mid D)=1-F_{\operatorname{Beta}(141,111)}(0.5)\approx0.97.\]这些是在指定先验后对参数的概率陈述。频率学派 $p$ 值则是在零假设成立时，出现至少同样极端数据的概率；它没有把概率分配给零假设本身。选择不同先验或参数空间，后验可以不同。\end{solve}

\clearpage\section*{区间估计与样本量：详解}\addcontentsline{toc}{section}{区间估计与样本量：详解}

\subsection*{训练题~\ref{q:M-C22-1} 小样本正态总体的三类区间}\label{a:M-C22-1}{\small 题目见第~\pageref{q:M-C22-1} 页；来源：MIT 第 22 讲，题 1(a—d)。}
\begin{solve}样本容量 $n=4$，均值 $\overline x=2.5$，离差平方和为 $5$，故 $s^2=5/3$。已知标准差时，均值区间为\[2.5\pm z_{0.95}\frac2{\sqrt4}=(0.8551,4.1449).\]未知标准差时使用 $t_3$ 分位数 $2.3534$，得到\[2.5\pm2.3534\sqrt{5/12}=(0.9809,4.0191).\]方差区间由卡方枢轴得到\[\left(\frac5{\chi^2_{3;0.95}},\frac5{\chi^2_{3;0.05}}\right)\approx(0.6398,14.2107).\]对两个端点取平方根，标准差区间约为 $(0.7999,3.7697)$。方差与标准差区间不可混用；原解中的一个标准差下端点数值已校订。\end{solve}

\subsection*{训练题~\ref{q:M-P10-1} 已知标准差的置信水平比较}\label{a:M-P10-1}{\small 题目见第~\pageref{q:M-P10-1} 页；来源：MIT 习题 10，题 1。}
\begin{solve}标准误为 $9/\sqrt{30}$。使用下侧分位数，两个区间分别为\[15\pm z_{0.975}\frac9{\sqrt{30}}=15\pm3.2205,\]\[15\pm z_{0.85}\frac9{\sqrt{30}}=15\pm1.7030.\]更高覆盖要求更宽区间。样本容量从 $30$ 变为 $90$ 时，宽度按 $1/\sqrt n$ 缩放，成为原来的 $1/\sqrt3$。样本均值只影响区间中心，不影响已知标准差模型中的宽度。\end{solve}

\subsection*{训练题~\ref{q:M-P10-2} 灌装量的正态均值区间}\label{a:M-P10-2}{\small 题目见第~\pageref{q:M-P10-2} 页；来源：MIT 习题 10，题 2（校订舍入）。}
\begin{solve}计算得 $\overline x=356$，离差平方和 $124$，所以 $s^2=15.5$。已知标准差时标准误为 $1$，故两个区间为 $356\pm1.9600$、$356\pm2.3263$。未知标准差时标准误为 $\sqrt{15.5}/3$，自由度为 $8$，故\[I_{0.95}=356\pm t_{8;0.975}\frac{\sqrt{15.5}}3,\]\[I_{0.98}=356\pm t_{8;0.99}\frac{\sqrt{15.5}}3.\]分别约为 $356\pm3.0263$ 与 $356\pm3.8011$。这里未知标准差版本较宽，同时受样本标准差大于 $3$ 和 $t$ 分位数大于正态分位数影响。\end{solve}

\subsection*{训练题~\ref{q:M-P10-3} 正态方差的区间}\label{a:M-P10-3}{\small 题目见第~\pageref{q:M-P10-3} 页；来源：MIT 习题 10，题 3。}
\begin{solve}样本均值约为 $6.1917$，无偏方差 $s^2\approx0.1535606$。因为 $11S^2/\sigma^2\sim\chi^2_{11}$，故\[P\left(\chi^2_{11;0.025}\le\frac{11S^2}{\sigma^2}\le\chi^2_{11;0.975}\right)=0.95.\]对正量求倒数并反向排列，得到方差区间\[\left(\frac{11s^2}{\chi^2_{11;0.975}},\frac{11s^2}{\chi^2_{11;0.025}}\right)\approx(0.07706,0.44268).\]卡方分布偏斜，且反演涉及倒数，所以区间通常不以 $s^2$ 为对称中心。\end{solve}

\subsection*{训练题~\ref{q:M-P11-1} 身高区间与宽度要求}\label{a:M-P11-1}{\small 题目见第~\pageref{q:M-P11-1} 页；来源：MIT 习题 11，题 1。}
\begin{solve}未知方差时使用 $t_{19;0.95}=1.7291$，得到\[69.55\pm1.7291\sqrt{14.26/20}\approx(68.09,71.01).\]已知标准差时使用 $z_{0.95}=1.6449$，得到\[69.55\pm1.6449(3.77/\sqrt{20})\approx(68.16,70.94).\]总宽度为 $2z_{0.95}3.77/\sqrt n$，令其不超过 $1$，得到 $n\ge(2z_{0.95}3.77)^2\approx153.81$，所以至少需要 $154$ 个观测。这里“宽度”是两个端点之差，不是误差半径。\end{solve}

\subsection*{训练题~\ref{q:M-P11-2} 顺序统计量构造端点区间}\label{a:M-P11-2}{\small 题目见第~\pageref{q:M-P11-2} 页；来源：MIT 习题 11，题 2。}
\begin{solve}令 $b_L,b_U$ 分别为 $\operatorname{Beta}(2,n-1)$ 的 $0.025,0.975$ 下侧分位数。则\[P\left(b_L\le\frac{Y_{(2)}}a\le b_U\right)=0.95.\]因所有量为正，反演得\[a\in\left[\frac{Y_{(2)}}{b_U},\frac{Y_{(2)}}{b_L}\right].\]当 $n=9$ 时，$b_L\approx0.02814,b_U\approx0.48250$，故区间约为 $[3.11,53.30]$。区间很宽，说明只使用第二小值时，上端点信息有限。也可以与必然约束 $a\ge X_{(n)}$ 相交而保持覆盖率，因为真实参数以概率一满足该约束。\end{solve}

\subsection*{训练题~\ref{q:M-C22-3} 比例误差与样本量}\label{a:M-C22-3}{\small 题目见第~\pageref{q:M-C22-3} 页；来源：MIT 第 22 讲，题 3（采用精确正态分位数）。}
\begin{solve}正态近似误差半径上界为 $z_{1-\alpha/2}/(2\sqrt n)$。因此 $95\%$ 水平要求\[n\ge\frac{z_{0.975}^2}{4(0.01)^2},\]向上取整得 $9604$；$80\%$ 水平使用 $z_{0.9}=1.28155$，得到 $4106$。若 $n=900$，两个半径分别约为 $0.03267$、$0.02136$。用 $z\approx2$ 的简便规则会得到 $10000$，这与精确正态分位数的规划值不同。这里仍是近似区间规划，尤其接近边界比例时应另查有限样本覆盖。\end{solve}

\subsection*{训练题~\ref{q:B13.1} 指数率的精确区间}\label{a:B13.1}{\small 题目见第~\pageref{q:B13.1} 页；来源：编订补充。}
\begin{solve}枢轴分布不含未知参数，因此\[P\left(\chi^2_{2n;\alpha/2}\le2\lambda T\le\chi^2_{2n;1-\alpha/2}\right)=1-\alpha.\]在 $T>0$ 时除以 $2T$，率区间为\[\left[\frac{\chi^2_{2n;\alpha/2}}{2T},\frac{\chi^2_{2n;1-\alpha/2}}{2T}\right].\]对端点取倒数并交换顺序，平均寿命区间为\[\left[\frac{2T}{\chi^2_{2n;1-\alpha/2}},\frac{2T}{\chi^2_{2n;\alpha/2}}\right].\]这里是指数抽样假设下的有限样本精确覆盖，不需要中心极限定理。\end{solve}

\subsection*{训练题~\ref{q:B13.2} 均匀端点的单侧区间}\label{a:B13.2}{\small 题目见第~\pageref{q:B13.2} 页；来源：编订补充。}
\begin{solve}对 $0<u<1$，$P(M/\theta\le u)=u^n$。由于 $M\le\theta$ 以概率一成立，区间包含参数只需满足 $\theta\le cM$，即 $M/\theta\ge1/c$。其概率为 $1-c^{-n}$。令它等于 $1-\alpha$，得到 $c=\alpha^{-1/n}$，所以\[[M,M/\alpha^{1/n}].\]此区间不以某个点估计为中心，而是利用总体支撑给出的确定性下界。构造置信区间不必总是“估计值加减标准误”。\end{solve}

\subsection*{训练题~\ref{q:B13.3} 两个均值差的合并方差区间}\label{a:B13.3}{\small 题目见第~\pageref{q:B13.3} 页；来源：编订补充。}
\begin{solve}定义\[S_p^2=\frac{(n-1)S_1^2+(m-1)S_2^2}{n+m-2}.\]由独立卡方变量相加，$(n+m-2)S_p^2/\sigma^2\sim\chi^2_{n+m-2}$；它与均值差独立。因此标准化统计量\[\frac{(\overline X-\overline Y)-(\mu_1-\mu_2)}{S_p\sqrt{1/n+1/m}}\sim t_{n+m-2}.\]反演得到均值差区间为 $(\overline X-\overline Y)\pm t_{n+m-2;1-\alpha/2}S_p\sqrt{1/n+1/m}$。合并平方和共用同一个真实方差，正是共同方差假设的用途。\end{solve}

\subsection*{训练题~\ref{q:B13.4} 韦尔奇近似区间}\label{a:B13.4}{\small 题目见第~\pageref{q:B13.4} 页；来源：编订补充。}
\begin{solve}差的估计标准误为 $SE=\sqrt{s_1^2/n+s_2^2/m}$，萨特思韦特近似自由度为\[\nu=\frac{(s_1^2/n+s_2^2/m)^2}{(s_1^2/n)^2/(n-1)+(s_2^2/m)^2/(m-1)}.\]本题 $SE=1$，$\nu=1/(0.16/9+0.36/14)\approx22.991$，区间为 $2\pm t_{\nu;0.975}$，约 $(-0.069,4.069)$。它通常是近似区间，不应标成自由度整数的精确 $t$ 结论。相比先做方差检验再决定方法，事先合理选定允许异方差的方法更容易保持清晰的推断规则。\end{solve}

\subsection*{训练题~\ref{q:B13.5} 配对差的区间}\label{a:B13.5}{\small 题目见第~\pageref{q:B13.5} 页；来源：编订补充。}
\begin{solve}差值均值为 $2$，离差平方和为 $10$，所以 $s_D^2=10/4=2.5$。标准误为 $\sqrt{2.5/5}=\sqrt{0.5}$，自由度为 $4$，区间为\[2\pm t_{4;0.975}\sqrt{0.5}\approx2\pm1.9632.\]同一受试者前后观测通常相关。差值方差包含 $\operatorname{Var}(X)+\operatorname{Var}(Y)-2\operatorname{Cov}(X,Y)$；把前后当作独立会遗漏协方差。配对分析所需正态假设针对差值，而不是仅分别针对前后边缘分布。\end{solve}

\subsection*{训练题~\ref{q:M-P10-5} 置信区间与后验区间不能互换}\label{a:M-P10-5}{\small 题目见第~\pageref{q:M-P10-5} 页；来源：MIT 习题 10，题 5（选编）。}
\begin{solve}标准误为 $5/8=0.625$，区间为 $1.6\pm1.959964(0.625)$，即约 $(0.3750,2.8250)$。它包含候选值 $1,2$，所以先验概率和为 $0.04$。后验权重为\[w_j=\pi_j\exp\left[-\frac{(1.6-j)^2}{2(0.625)^2}\right],\quad j=0,1,2,3.\]归一化后约为 $(0.537,0.191,0.247,0.025)$，所以区间的后验概率约为 $0.438$，不是 $0.95$。频率覆盖是对固定真实参数反复抽样时的随机区间而言；后验概率是在给定数据和先验以后对参数的不确定性而言。\end{solve}

\subsection*{训练题~\ref{q:B13.6} 自助法百分位区间}\label{a:B13.6}{\small 题目见第~\pageref{q:B13.6} 页；来源：编订补充。}
\begin{solve}把经验分布视为总体近似。每次从已有 $n$ 个观测中有放回抽取 $n$ 个值，计算自助样本中位数，重复 $B$ 次得到 $T_1^*,\ldots,T_B^*$。以其经验 $\alpha/2$ 与 $1-\alpha/2$ 分位数作为百分位区间端点。原始抽样随机性产生数据，自助重抽样随机性近似统计量的抽样波动，有限 $B$ 还引入蒙特卡洛误差。若总体在中位数附近有正且连续的密度，通常具有良好渐近基础；若分位数不唯一、存在大原子或样本极小，普通百分位法可能覆盖不佳。增大 $B$ 不能弥补原始样本代表性不足。\end{solve}

\clearpage\section*{假设检验与功效：详解}\addcontentsline{toc}{section}{假设检验与功效：详解}

\subsection*{训练题~\ref{q:M-P8-1ab} 硬币的精确双侧检验}\label{a:M-P8-1ab}{\small 题目见第~\pageref{q:M-P8-1ab} 页；来源：MIT 习题 8，题 1(a,b)。}
\begin{solve}在零假设下 $X\sim\operatorname{Bin}(250,0.5)$，分布关于 $125$ 对称。按偏离 $125$ 的绝对距离度量极端程度，\[p_{\rm exact}=P(X\le110)+P(X\ge140)=2P(X\ge140)\approx0.06642.\]因此在 $0.10$ 水平拒绝，在 $0.05$ 水平不拒绝。未修正正态近似使用 $z=(140-125)/\sqrt{62.5}$，给出约 $0.05778$；若右尾阈值改用 $139.5$，通常更接近精确概率。离散分布下不同双侧定义可能不同，本题因对称性定义很自然。\end{solve}

\subsection*{训练题~\ref{q:M-P8-1cd} 离散拒绝域与功效}\label{a:M-P8-1cd}{\small 题目见第~\pageref{q:M-P8-1cd} 页；来源：MIT 习题 8，题 1(c,d)。}
\begin{solve}$0.01$ 水平的拒绝域取 $X\le104$ 或 $X\ge146$，每个尾的零假设概率不超过 $0.005$。$0.05$ 水平对应 $X\le109$ 或 $X\ge141$。因此在任意给定 $p$ 下，功效为\[\pi(p)=\sum_{k=0}^{109}\binom{250}{k}p^k(1-p)^{250-k}+\sum_{k=141}^{250}\binom{250}{k}p^k(1-p)^{250-k}.\]计算得 $\pi(0.55)\approx0.35237$，$\pi(0.60)\approx0.88963$。效应远离零假设时，分布落入拒绝域的概率增大。离散检验的实际大小通常小于指定水平，不能总写成严格等号。\end{solve}

\subsection*{训练题~\ref{q:M-P8-2} 第一类与第二类错误}\label{a:M-P8-2}{\small 题目见第~\pageref{q:M-P8-2} 页；来源：MIT 习题 8，题 2（抽象分类表述）。}
\begin{solve}第一类错误是对真实零假设作拒绝，因此估计错误率为 $9/140\approx0.06429$。第二类错误是阳性对象未被识别，其估计概率为 $15/140\approx0.10714$，功效为 $125/140\approx0.89286$。阳性结果为假的概率是逆向条件概率，需结合目标人群中的真实阳性比例。此处两类各抽取 $140$ 个是实验设计，不保证目标人群也各占一半。错误率的条件方向不可交换。\end{solve}

\subsection*{训练题~\ref{q:M-P8-3} 复合零假设与功效规划}\label{a:M-P8-3}{\small 题目见第~\pageref{q:M-P8-3} 页；来源：MIT 习题 8，题 3（单位改称测量单位）。}
\begin{solve}对固定右尾拒绝域，误拒概率随 $\mu$ 增大，故零假设中的最不利值为 $40$。阈值为\[c=40+z_{0.96}\frac5{\sqrt3}\approx45.0538.\]在 $\mu=45$ 时，功效为 $1-\Phi((c-45)/(5/\sqrt3))\approx0.4926$。一般样本量下，功效为 $\Phi(\sqrt n-z_{0.96})$；要求其至少 $0.9$，即\[\sqrt n\ge z_{0.96}+z_{0.9}\approx3.0322.\]所以最小整数样本量为 $10$。复合零假设要求控制其中所有参数值下的错误率，不能任取某个内部点代替最不利点。\end{solve}

\subsection*{训练题~\ref{q:M-P8-4} 直接计算两类错误}\label{a:M-P8-4}{\small 题目见第~\pageref{q:M-P8-4} 页；来源：MIT 习题 8，题 4。}
\begin{solve}零假设下拒绝域在支撑内总长度为 $0.2$，所以第一类错误概率为 $0.2/2=0.1$。当 $\theta=2.5$ 时，不拒绝区间长度为 $1.8$，第二类错误概率为 $1.8/2.5=0.72$。一般功效为\[\pi(\theta)=\begin{cases}1,&0<\theta\le0.1,\\0.1/\theta,&0.1<\theta\le1.9,\\1-1.8/\theta,&\theta>1.9.\end{cases}\]这是一条关于真实参数的函数，而不是单一数字。功效在某些接近零假设的备择处也可能偏低。\end{solve}

\subsection*{训练题~\ref{q:M-P9-1} 单侧与双侧 t 检验}\label{a:M-P9-1}{\small 题目见第~\pageref{q:M-P9-1} 页；来源：MIT 习题 9，题 1。}
\begin{solve}零假设下\[T=\frac{\overline X-10}{S/\sqrt{16}}\sim t_{15},\qquad t_{\rm obs}=2.\]双侧 $p$ 值为 $2[1-F_{t_{15}}(2)]\approx0.06395$，故不拒绝；右侧 $p$ 值为 $1-F_{t_{15}}(2)\approx0.03197$，故拒绝。两种检验回答不同问题、采用不同拒绝域。不能看到样本均值偏高后才临时把双侧检验改成右侧检验并仍声称维持原先的检验水平。\end{solve}

\subsection*{训练题~\ref{q:M-P9-2} 停止规则改变抽样分布}\label{a:M-P9-2}{\small 题目见第~\pageref{q:M-P9-2} 页；来源：MIT 习题 9，题 2。}
\begin{solve}固定次数方案以正面数 $H\sim\operatorname{Bin}(12,0.5)$ 为统计量，$p$ 值为 $P(H\le3)\approx0.07300$。停止方案记录第三个正面前反面数 $N$，其负二项概率为 $P(N=k)=\binom{k+2}{2}2^{-(k+3)}$，故\[P(N\ge9)\approx0.03271.\]两个抽样空间不同，频率检验结果可以不同。对于预先规定且不另含参数依赖因素的上述停止规则，观测数据似然都与 $p^3(1-p)^9$ 成比例；$\operatorname{Beta}(a,b)$ 先验均更新为 $\operatorname{Beta}(a+3,b+9)$。这不意味着任意自适应筛选和未报告数据都可忽略。\end{solve}

\subsection*{训练题~\ref{q:M-E2-II4} 正态右尾检验的分位数}\label{a:M-E2-II4}{\small 题目见第~\pageref{q:M-E2-II4} 页；来源：MIT 第二次考试，第二部分题 4。}
\begin{solve}临界值满足 $P(Z>c)=0.1$，所以 $c=\Phi^{-1}(0.9)\approx1.28155$。观测 $p$ 值为\[p=P(Z\ge2)=1-\Phi(2)\approx0.02275.\]因 $2>c$ 或等价地 $p<0.1$，拒绝零假设。分位数函数由给定尾概率反求横轴阈值；分布函数由给定横轴值计算累积概率。把 $\Phi(0.9)$ 当作临界值，是将两种函数方向弄反。\end{solve}

\subsection*{训练题~\ref{q:M-F-C10} 检验与置信区间对偶}\label{a:M-F-C10}{\small 题目见第~\pageref{q:M-F-C10} 页；来源：MIT 期末试卷，概念题 10。}
\begin{solve}依次为错、错、对、不能作这种频率学派概率解释。$p=0.08>0.05$，故在 $0.05$ 水平不拒绝，零位于对应 $95\%$ 区间；但 $p<0.10$，所以在 $0.10$ 水平拒绝，零不在 $90\%$ 区间。区间由接受域反演而来时，这种对偶精确成立。最后一句把随机区间的重复抽样覆盖与给定数据后的参数后验概率混在一起，需要额外先验才能作后者陈述。\end{solve}

\subsection*{训练题~\ref{q:B14.1} 奈曼—皮尔逊检验}\label{a:B14.1}{\small 题目见第~\pageref{q:B14.1} 页；来源：编订补充。}
\begin{solve}备择与零假设的对数似然比为\[\log\frac{L(\mu_1)}{L(\mu_0)}=\frac{n(\mu_1-\mu_0)}{\sigma^2}\overline X-\frac{n(\mu_1^2-\mu_0^2)}{2\sigma^2}.\]它随 $\overline X$ 严格增加。奈曼—皮尔逊引理说明应在大均值时拒绝，阈值由零假设大小确定为 $c=\mu_0+z_{1-\alpha}\sigma/\sqrt n$。在 $\mu_1$ 下功效为\[1-\Phi\left(z_{1-\alpha}-\frac{(\mu_1-\mu_0)\sqrt n}{\sigma}\right).\]最强性针对指定的简单备择以及所有同水平检验。似然比排序确定拒绝方向，零假设分布确定阈值。\end{solve}

\subsection*{训练题~\ref{q:B14.2} 正态均值的似然比检验}\label{a:B14.2}{\small 题目见第~\pageref{q:B14.2} 页；来源：编订补充。}
\begin{solve}全参数空间中的最大似然在 $\widehat\mu=\overline X$ 达到。利用平方和分解，\[\sum_i(X_i-\mu_0)^2-\sum_i(X_i-\overline X)^2=n(\overline X-\mu_0)^2.\]所以\[-2\log\Lambda=\frac{n(\overline X-\mu_0)^2}{\sigma^2}=Z^2.\]零假设下它精确服从 $\chi^2_1$，等价拒绝域为 $|Z|>z_{1-\alpha/2}$。本例无需调用渐近的威尔克斯定理，因为统计量的有限样本分布已经精确已知。\end{solve}

\subsection*{训练题~\ref{q:B14.3} 多重检验与邦费罗尼控制}\label{a:B14.3}{\small 题目见第~\pageref{q:B14.3} 页；来源：编订补充；主题参照 MIT 习题 8，题 5(b)。}
\begin{solve}独立时无一次误拒的概率为 $0.95^{20}$，所以至少一次误拒的概率为\[1-0.95^{20}\approx0.64151.\]若不保证独立，可令每个检验的显著性水平不超过 $0.05/20=0.0025$。由并事件上界，至少一次误拒的概率不超过各真实零假设误拒概率之和，因而不超过 $0.05$。这个邦费罗尼控制不需要独立性，并且在只有部分零假设真实时仍成立。代价是单个检验的功效可能降低。\end{solve}

\subsection*{训练题~\ref{q:B14.4} 不拒绝不等于等效}\label{a:B14.4}{\small 题目见第~\pageref{q:B14.4} 页；来源：编订补充。}
\begin{solve}未拒绝可能由差确实小，也可能由数据精度不足造成。此例的 $95\%$ 近似区间为 $0.2\pm1.96$，相当宽，包含许多实际重要的差异。等效检验把“不等效”作为零假设，即差不大于 $-0.5$ 或不小于 $0.5$。在双单侧检验中，各单侧使用 $\alpha=0.05$，等价于检查 $90\%$ 双侧区间是否完全包含在 $(-0.5,0.5)$ 内。本例 $90\%$ 区间为 $0.2\pm1.6449$，并未落入等效界，不能认定等效。界限应在观察结果前由问题背景确定。\end{solve}

\clearpage\section*{拟合优度与列联表：详解}\addcontentsline{toc}{section}{拟合优度与列联表：详解}

\subsection*{训练题~\ref{q:M-P9-3} 单个正态总体的方差检验}\label{a:M-P9-3}{\small 题目见第~\pageref{q:M-P9-3} 页；来源：MIT 习题 9，题 3。}
\begin{solve}样本均值为 $0.05$，离差平方和为 $25.7372$，故 $s^2\approx2.33975$。在零假设下\[Q=\frac{(n-1)S^2}{1}\sim\chi^2_{11}.\]观察值 $q=25.7372$，右尾 $p$ 值约为 $0.0071$，小于 $0.05$，所以拒绝零假设。均值是未知并由样本估计的，因此自由度为 $11$。如果均值已知而用相对已知均值的平方和，统计量和自由度都不同。\end{solve}

\subsection*{训练题~\ref{q:M-E2-II5} 两个总体的类别比例比较}\label{a:M-E2-II5}{\small 题目见第~\pageref{q:M-E2-II5} 页；来源：MIT 第二次考试，第二部分题 5（教学数据）。}
\begin{solve}合并后的类别比例估计为 $20/100=0.2$ 与 $0.8$，所以每组期望计数均为 $(10,40)$。皮尔逊统计量为\[Q=\frac{(15-10)^2}{10}+\frac{(35-40)^2}{40}+\frac{(5-10)^2}{10}+\frac{(45-40)^2}{40}=6.25.\]自由度为 $(2-1)(2-1)=1$，渐近 $p$ 值约为 $0.01242$，故拒绝比例相同的零假设。结论是模型和抽样条件下的比例差异证据；不能仅凭观察列联表推出一个类别对另一个变量有因果作用。\end{solve}

\subsection*{训练题~\ref{q:B15.1} 多项分类的拟合优度}\label{a:B15.1}{\small 题目见第~\pageref{q:B15.1} 页；来源：编订补充。}
\begin{solve}零假设指定每类概率均为 $1/6$，所以期望计数都是 $20$。统计量为\[Q=\sum_{j=1}^6\frac{(O_j-20)^2}{20}=\frac{4+4+1+1+9+9}{20}=1.4.\]六个计数和受总数约束，未从数据估计额外概率参数，所以自由度为 $6-1=5$。$\chi^2_{5;0.95}\approx11.0705$，观察值远小于临界值，不拒绝公平性。这个结果不证明骰子绝对公平，只说明当前样本未给出足够的偏离证据。\end{solve}

\subsection*{训练题~\ref{q:B15.2} 估计参数后的自由度}\label{a:B15.2}{\small 题目见第~\pageref{q:B15.2} 页；来源：编订补充。}
\begin{solve}若零假设中一个自由参数通过适当的同一分组多项似然方法估计，且满足常规正则性、各类别概率为正并有足够期望计数，则皮尔逊统计量渐近自由度为 $k-1-1=k-2$。其中一个自由度来自总概率为一，另一个来自估计参数。若用未分组信息估计参数而又对分组计数作检验，精确渐近分布未必简单地等于该卡方分布，不能无条件扣一。分组应覆盖全部支撑，尾部低期望类别可合并；根据观察偏差反复调整分组会影响原定检验。\end{solve}

\subsection*{训练题~\ref{q:B15.3} 列联表独立性检验}\label{a:B15.3}{\small 题目见第~\pageref{q:B15.3} 页；来源：编订补充。}
\begin{solve}总数为 $200$，行和均为 $100$，列和为 $50,60,90$。在独立模型下，期望计数为行和乘列和除以总数，两行均为 $(25,30,45)$。故\[Q=2\left(\frac{25}{25}+0+\frac{25}{45}\right)=\frac{28}{9}\approx3.1111.\]自由度为 $(2-1)(3-1)=2$，右尾概率为 $e^{-Q/2}\approx0.2111$。在 $0.05$ 水平不拒绝独立性。独立模型含有行比例与列比例，约束扣除后才得到上述自由度。\end{solve}

\subsection*{训练题~\ref{q:B15.4} 小样本的费希尔精确检验}\label{a:B15.4}{\small 题目见第~\pageref{q:B15.4} 页；来源：编订补充。}
\begin{solve}两个行和与列和均为 $5$，总数为 $10$。在零假设下，条件分布为超几何分布：\[P(A=a)=\frac{\binom5a\binom5{5-a}}{\binom{10}5},\quad a=0,\ldots,5.\]观测 $A=4$，右侧 $p$ 值为 $(25+1)/252=13/126\approx0.10317$。双侧按题设概率排序纳入 $a=0,1,4,5$，概率总和为 $52/252=13/63\approx0.20635$。小期望计数时，不能把大样本卡方近似当作精确分布；双侧精确检验还须说明采用的极端程度定义。\end{solve}

\subsection*{训练题~\ref{q:M-P9-5} 两个正态方差的 F 检验}\label{a:M-P9-5}{\small 题目见第~\pageref{q:M-P9-5} 页；来源：MIT 习题 9，题 5（给出原数据的汇总量）。}
\begin{solve}统计量取固定顺序的比值 $F=S_1^2/S_2^2$，零假设下服从 $F_{19,19}$；观测值约为 $0.97038$。按双倍较小尾概率定义，\[p=2\min\{F_{19,19}^{\rm cdf}(f),1-F_{19,19}^{\rm cdf}(f)\}\approx0.9484.\]不拒绝等方差。方差比 $r=\sigma_1^2/\sigma_2^2$ 的枢轴为 $F/r$，故区间为\[\left[\frac{f}{F_{19,19;0.975}},\frac{f}{F_{19,19;0.025}}\right]\approx[0.3841,2.4516].\]区间较宽，说明“不显著”也可能伴随相当大的不确定性；此检验对偏离正态较敏感。\end{solve}

\clearpage\section*{回归与方差分析：详解}\addcontentsline{toc}{section}{回归与方差分析：详解}

\subsection*{训练题~\ref{q:M-C26-1a} 三个点的最小二乘直线}\label{a:M-C26-1a}{\small 题目见第~\pageref{q:M-C26-1a} 页；来源：MIT 第 26 讲，题 1(a) 与题 2(a—c)。}
\begin{solve}模型为 $y_i=a x_i+b+e_i$。最小化\[Q(a,b)=(3-a-b)^2+(1-2a-b)^2+(4-4a-b)^2.\]对两个参数求偏导得到正规方程\[21a+7b=21,\qquad7a+3b=8.\]解得 $a=1/2,b=3/2$，所以 $\widehat y=x/2+3/2$。$\overline x=7/3,\overline y=8/3$，代入直线确有 $\widehat y(7/3)=8/3$。设计矩阵两列线性无关，平方和的 Hessian 正定，因此该驻点是唯一全局最小值。\end{solve}

\subsection*{训练题~\ref{q:M-C26-1b} 二次多项式仍是线性回归}\label{a:M-C26-1b}{\small 题目见第~\pageref{q:M-C26-1b} 页；来源：MIT 第 26 讲，题 1(b,d)（选编）。}
\begin{solve}设计矩阵的三行为 $(1,1,1),(4,2,1),(16,4,1)$，响应向量为 $(3,1,4)^\mathsf T$。三个横坐标不同，矩阵可逆；解得\[a=7/6,\qquad b=-11/2,\qquad c=22/3.\]曲线经过所有三个点，残差平方和为零。所谓线性回归是指对未知参数线性，而不要求对自变量 $x$ 线性；三次模型只需增加 $x^3$ 这一列。但三个观测估计三个参数时残差自由度为零，不能再用残差估计误差方差，也不能据零残差断言泛化预测完美。\end{solve}

\subsection*{训练题~\ref{q:M-C26-1c} 对数尺度的指数拟合}\label{a:M-C26-1c}{\small 题目见第~\pageref{q:M-C26-1c} 页；来源：MIT 第 26 讲，题 1(c)（补充模型区别）。}
\begin{solve}令 $z_i=\log y_i$。对数尺度的目标是\[Q(a,b)=\sum_i(z_i-ax_i-b)^2,\]解为 $a=\sum_i(x_i-\overline x)(z_i-\overline z)/\sum_i(x_i-\overline x)^2$，$b=\overline z-a\overline x$。代入本题得 $a\approx0.18122,b\approx0.40547$。两种目标通常给出不同拟合，因为对数尺度侧重相对误差，原尺度侧重绝对误差。若模型为 $\log Y=ax+b+\varepsilon$、$\varepsilon\sim N(0,\sigma^2)$，则 $e^{ax+b}$ 是条件中位数，条件均值还需乘 $e^{\sigma^2/2}$。\end{solve}

\subsection*{训练题~\ref{q:M-C26-2d} 正态误差下的最小二乘与似然}\label{a:M-C26-2d}{\small 题目见第~\pageref{q:M-C26-2d} 页；来源：MIT 第 26 讲，题 2(d)。}
\begin{solve}对数似然为\[\ell(a,b,\sigma^2)=-\frac n2\log(2\pi\sigma^2)-\frac{1}{2\sigma^2}\sum_i(y_i-ax_i-b)^2.\]固定正方差时，最大化似然等价于最小化残差平方和 $SSE$。代回最小二乘参数，再对方差求导，得到 $\widehat\sigma^2_{\rm ML}=SSE/n$。当设计秩为二且 $n>2$ 时，无偏估计则为 $SSE/(n-2)$。若 $SSE=0$，在严格正方差空间可能不存在有限极大值。只说“独立同方差”不足以推出似然等价，还需要正态误差或另行指定的误差分布。\end{solve}

\subsection*{训练题~\ref{q:B16.1} 正规方程与残差正交}\label{a:B16.1}{\small 题目见第~\pageref{q:B16.1} 页；来源：编订补充。}
\begin{solve}平方和为 $(Y-X\beta)^\mathsf T(Y-X\beta)$，梯度为 $-2X^\mathsf T(Y-X\beta)$。令其为零得到正规方程 $X^\mathsf TX\widehat\beta=X^\mathsf TY$。满列秩保证 $X^\mathsf TX$ 正定可逆，故\[\widehat\beta=(X^\mathsf TX)^{-1}X^\mathsf TY.\]残差 $e=Y-X\widehat\beta$ 满足 $X^\mathsf Te=0$。若全一向量是 $X$ 的一列，则 $\sum_ie_i=0$，于是样本响应均值等于拟合值均值。没有截距项时，残差和为零和通常的中心化平方和分解都不一定成立。\end{solve}

\subsection*{训练题~\ref{q:B16.2} 高斯—马尔可夫最优性}\label{a:B16.2}{\small 题目见第~\pageref{q:B16.2} 页；来源：编订补充。}
\begin{solve}任意线性无偏估计可写为 $AY$，其中 $AX=I_p$。记 $A_0=(X^\mathsf TX)^{-1}X^\mathsf T$，令 $D=A-A_0$，则 $DX=0$。因为 $A_0D^\mathsf T=(X^\mathsf TX)^{-1}X^\mathsf TD^\mathsf T=0$，\[\operatorname{Cov}(AY)=\sigma^2A_0A_0^\mathsf T+\sigma^2DD^\mathsf T.\]第二项半正定，所以任何线性组合的方差都不小于最小二乘对应方差。这是在线性无偏估计类中的最优性，不是对所有可能的非线性或有偏估计的统一最优性。证明只用一二阶矩条件。\end{solve}

\subsection*{训练题~\ref{q:B16.3} 斜率估计的方差与区间}\label{a:B16.3}{\small 题目见第~\pageref{q:B16.3} 页；来源：编订补充。}
\begin{solve}斜率估计为 $\widehat a=\sum_i(x_i-\overline x)Y_i/S_{xx}$。代入模型，利用中心化和为零，得 $E\widehat a=a$；独立同方差误差给出\[\operatorname{Var}(\widehat a)=\frac{\sigma^2}{S_{xx}}.\]正态假设使它为正态变量，且与 $s^2=SSE/(n-2)$ 独立，因此\[\frac{\widehat a-a}{s/\sqrt{S_{xx}}}\sim t_{n-2}.\]区间为 $\widehat a\pm t_{n-2;1-\alpha/2}s/\sqrt{S_{xx}}$。在模型可靠和误差方差不变时，增大 $S_{xx}$ 可提高斜率精度，但超出已观察范围的外推还面临模型形式风险。\end{solve}

\subsection*{训练题~\ref{q:B16.4} 均值响应区间与预测区间}\label{a:B16.4}{\small 题目见第~\pageref{q:B16.4} 页；来源：编订补充。}
\begin{solve}令 $\widehat y_0=\widehat a x_0+\widehat b$，其方差为\[\sigma^2h_0,\qquad h_0=\frac1n+\frac{(x_0-\overline x)^2}{S_{xx}}.\]均值响应区间为 $\widehat y_0\pm t_{n-2;1-\alpha/2}s\sqrt{h_0}$。独立新响应还含新误差 $\varepsilon_0$，所以预测误差方差为 $\sigma^2(1+h_0)$，预测区间为\[\widehat y_0\pm t_{n-2;1-\alpha/2}s\sqrt{1+h_0}.\]额外的 $1$ 对应个体噪声。即使样本无限多，个体预测区间通常也不会收缩为零，而均值响应估计的不确定性可随数据增加而缩小。\end{solve}

\subsection*{训练题~\ref{q:B16.5} 单因素方差分析}\label{a:B16.5}{\small 题目见第~\pageref{q:B16.5} 页；来源：编订补充。}
\begin{solve}三组均值为 $3,5,7$，总均值为 $5$。组间平方和\[SS_B=3[(3-5)^2+0+(7-5)^2]=24,\]自由度为 $2$；各组内平方和都是 $2$，故 $SS_W=6$，自由度为 $9-3=6$。统计量为\[F=\frac{SS_B/2}{SS_W/6}=12.\]零假设下服从 $F_{2,6}$，右尾概率为 $(6/(6+2\times12))^3=0.008$，因此拒绝均值全相等。该结论说明至少有一对均值不同；具体哪些组不同需另作有误差控制的比较。\end{solve}

\subsection*{训练题~\ref{q:B16.6} 相关系数与回归解释度}\label{a:B16.6}{\small 题目见第~\pageref{q:B16.6} 页；来源：编订补充。}
\begin{solve}记 $y_i-\overline y=(\widehat y_i-\overline y)+e_i$。由残差与常数列、自变量列正交，交叉项和为零，所以 $SST=SSR+SSE$。又 $\widehat a=S_{xy}/S_{xx}$，因此 $SSR=\widehat a^2S_{xx}=S_{xy}^2/S_{xx}$，从而\[R^2=\frac{SSR}{S_{yy}}=\frac{S_{xy}^2}{S_{xx}S_{yy}}=r^2.\]需要 $S_{xx}>0,S_{yy}>0$，否则斜率或相关系数可能未定义。该等式针对含截距的一元模型；多元回归中不能用某一个自变量与响应的简单相关系数平方代替整个模型的 $R^2$。\end{solve}

\clearpage\section*{综合训练与解题辨析：详解}\addcontentsline{toc}{section}{综合训练与解题辨析：详解}

\subsection*{训练题~\ref{q:M-F-5} 帕累托模型从估计到区间}\label{a:M-F-5}{\small 题目见第~\pageref{q:M-F-5} 页；来源：MIT 期末试卷，题 5；区间为编订延伸。}
\begin{solve}设 $T=\sum_i\log x_i$。似然与 $\theta^n\exp(-\theta T)$ 成比例，故对数导数为 $n/\theta-T$，当 $T>0$ 时极大似然估计为 $n/T$。在 $y\ge0$ 上，\[P(Y>y)=P(X>e^y)=e^{-\theta y},\]所以 $Y\sim\operatorname{Exp}(\theta)$，总和 $T\sim\operatorname{Gamma}(n,\theta)$。因此\[\left[\frac{\chi^2_{2n;\alpha/2}}{2T},\frac{\chi^2_{2n;1-\alpha/2}}{2T}\right]\]是精确 $1-\alpha$ 区间。最后，$EX=\int_1^\infty\theta x^{-\theta}\,dx$ 仅在 $\theta>1$ 时有限，值为 $\theta/(\theta-1)$。率参数始终可估计，并不要求原始变量存在均值；不能对所有 $\theta>0$ 都使用基于 $\overline X$ 的有限均值渐近理论。\end{solve}

\subsection*{训练题~\ref{q:M-PF-26} 瑞利参数与变换统计量}\label{a:M-PF-26}{\small 题目见第~\pageref{q:M-PF-26} 页；来源：MIT 模拟期末，题 26；后两问为编订延伸。}
\begin{solve}对数似然中参数相关项为 $n\log\tau-(\tau/2)\sum_i x_i^2$，故\[\widehat\tau=\frac{2n}{\sum_i x_i^2}.\]对 $z>0$，$P(X^2/2>z)=e^{-\tau z}$，因此 $Z$ 为率 $\tau$ 的指数变量，$T=\sum_iX_i^2/2\sim\operatorname{Gamma}(n,\tau)$。当 $n>1$，无偏率估计为 $(n-1)/T$；又 $2\tau T=\tau\sum_iX_i^2\sim\chi^2_{2n}$，所以区间为\[\left[\frac{\chi^2_{2n;\alpha/2}}{\sum_iX_i^2},\frac{\chi^2_{2n;1-\alpha/2}}{\sum_iX_i^2}\right].\]同一模型可通过变量变换复用指数抽样理论，但必须明确瑞利分布这里采用的参数化。\end{solve}

\subsection*{训练题~\ref{q:M-F-7} 定位中的连续更新与预测}\label{a:M-F-7}{\small 题目见第~\pageref{q:M-F-7} 页；来源：MIT 期末试卷，题 7。}
\begin{solve}初始先验为 $(1/4,1/4,1/4,1/4)$。橙色读数的似然为 $(0.7,0.7,0.3,0.3)$，归一化后得到\[(0.35,0.35,0.15,0.15).\]运动后概率按位置循环右移，成为 $(0.15,0.35,0.35,0.15)$。下一次蓝色读数的似然为 $(0.3,0.3,0.7,0.7)$，全概率公式给出预测概率\[0.15(0.3)+0.35(0.3)+0.35(0.7)+0.15(0.7)=0.5.\]若实际又读到蓝色，未归一化后验为 $(0.045,0.105,0.245,0.105)$，归一化得 $(0.09,0.21,0.49,0.21)$。这是“运动预测—观测更新”的两步，不能把移动前的门标签直接沿用。\end{solve}

\subsection*{训练题~\ref{q:M-F-8} 六类别概率的综合检验}\label{a:M-F-8}{\small 题目见第~\pageref{q:M-F-8} 页；来源：MIT 期末试卷，题 8（删除品牌及时效性陈述）。}
\begin{solve}期望计数为 $(200,200,200,150,150,100)$。皮尔逊统计量为\[Q=2+0.5+1.125+\frac23+1.5+4=9.791666\ldots.\]零假设完全指定概率，故自由度为 $5$。右尾概率约为 $0.08136$，在 $0.05$ 水平不拒绝。若同一包装内的对象颜色相关，计数未必服从简单独立多项模型；单把总数当成独立样本容量会误估波动。题设独立性是采用该参考分布的重要依据，而不是由对象数量大自动保证的条件。\end{solve}

\subsection*{训练题~\ref{q:M-F-9} 两样本检验的完整流程}\label{a:M-F-9}{\small 题目见第~\pageref{q:M-F-9} 页；来源：MIT 期末试卷，题 9（改称甲乙两总体）。}
\begin{solve}先说明抽样独立和共同方差是假设，不能由两个样本标准差接近自动证明。合并方差为\[s_p^2=\frac{10(64)+10(36)}{20}=50.\]均值差标准误为 $\sqrt{50(2/11)}=\sqrt{100/11}$，统计量为 $t=5/\sqrt{100/11}\approx1.65831$，零假设下自由度为 $20$。双侧 $p$ 值约 $0.1129$，故不拒绝。甲总体均值区间独立地使用甲样本的 $t_{10}$ 分布：\[115\pm t_{10;0.975}\frac8{\sqrt{11}}\approx(109.63,120.37).\]两组比较的自由度与单组区间的自由度不同，不能照搬。\end{solve}

\subsection*{训练题~\ref{q:M-F-11} 双侧检验在指定备择下的功效}\label{a:M-F-11}{\small 题目见第~\pageref{q:M-F-11} 页；来源：MIT 期末试卷，题 11。}
\begin{solve}零假设下标准误为 $0.1$，不拒绝区间为 $[-0.1960,0.1960]$（端点已舍入），其补集才是拒绝域。在真实均值为 $0.4$ 时，将端点相对 $0.4$ 标准化，功效为\[\begin{aligned}\pi(0.4)&=P_{0.4}(|\overline X|>0.195996)\\&=\Phi(-5.959964)+1-\Phi(-2.040036)\\&\approx0.97933.\end{aligned}\]若误把中央区间叫作拒绝域，会计算成第二类错误概率，约为 $0.02067$。核对区间方向可用直观检验：当真实均值与零假设相差四个标准误时，功效理应较高。\end{solve}

\subsection*{训练题~\ref{q:M-F-12} 折线基函数的最小二乘}\label{a:M-F-12}{\small 题目见第~\pageref{q:M-F-12} 页；来源：MIT 期末试卷，题 12（以文字替代作图）。}
\begin{solve}记 $z_i=|x_i-2|$，本题为 $(2,4,2)$，模型对参数 $b$ 线性且无截距。平方和导数给出\[\widehat b=\frac{\sum_i z_iY_i}{\sum_i z_i^2}=\frac{2(1)+4(1.5)+2(0.5)}{4+16+4}=\frac38.\]独立同方差误差保证 $E\widehat b=b$、$\operatorname{Var}(\widehat b)=\sigma^2/24$。拟合曲线为 $\widehat y=(3/8)|x-2|$，顶点 $(2,0)$，左段斜率 $-3/8$，右段斜率 $3/8$。这里自变量变换后的列向量只有一列，不应额外加入未在题设中的截距。\end{solve}

\subsection*{训练题~\ref{q:M-F-13} 自助分位数的读法}\label{a:M-F-13}{\small 题目见第~\pageref{q:M-F-13} 页；来源：MIT 期末试卷，题 13（保留所需次序统计量）。}
\begin{solve}经验分布逆函数将概率 $q$ 的分位数定义为排序后的第 $\lceil Bq\rceil$ 个值。对 $B=500$，$0.05$ 与 $0.95$ 分位数分别对应第 $25$ 和第 $475$ 个值，所以区间为 $[6.740,7.224]$。不同软件可能采用插值分位数，因此末位可略有差异。增加 $B$ 主要减少自助分位数的蒙特卡洛误差，并不增加原始样本包含的总体信息。百分位区间不是自动精确的 $90\%$ 区间，其覆盖依赖统计量的正则性、样本代表性和近似质量。\end{solve}

\subsection*{训练题~\ref{q:M-PF-18} 抽球机制与逆向条件概率}\label{a:M-PF-18}{\small 题目见第~\pageref{q:M-PF-18} 页；来源：MIT 模拟期末，题 18。}
\begin{solve}令 $O_i,B_i$ 分别表示第 $i$ 次为橙、蓝。第一种路径概率为\[P(O_1O_2)=\frac35\frac24=\frac3{10}.\]若第一次为蓝，第二次袋中有四橙二蓝，故 $P(B_1O_2)=(2/5)(4/6)=4/15$。于是 $P(O_2)=3/10+4/15=17/30$，贝叶斯公式给出\[P(B_1\mid O_2)=\frac{4/15}{17/30}=\frac8{17}.\]两次取球不独立；把第二次橙球概率固定成最初的 $3/5$ 会忽略袋中构成随路径变化。\end{solve}

\subsection*{训练题~\ref{q:M-PF-19} 停时分布的边界计算}\label{a:M-PF-19}{\small 题目见第~\pageref{q:M-PF-19} 页；来源：MIT 模拟期末，题 19；均值为编订延伸。}
\begin{solve}一次点数最多为 $6$，所以第一项为零。两次总点数大于 $6$ 的有序组合有 $21$ 个，因此第二项为 $21/36=7/12$。七次即使都为一点，总和也为 $7$，所以第三项为一。对 $1\le k\le6$，事件 $N>k$ 等价于 $k$ 个正整数点数之和不超过 $6$；每个自动不超过六。正整数分拆的有序计数为 $\binom6k$，故\[EN=\sum_{k=0}^6P(N>k)=\sum_{k=0}^6\binom6k6^{-k}=(7/6)^6.\]其中 $k=0$ 项为一。这个停时有界，所以尾和和期望交换没有收敛困难。\end{solve}

\subsection*{训练题~\ref{q:M-PF-20} 分段密度与分位数}\label{a:M-PF-20}{\small 题目见第~\pageref{q:M-PF-20} 页；来源：MIT 模拟期末，题 20。}
\begin{solve}两个三角形面积各为 $1/2$，密度归一化。对 $0\le x\le1/2$，$F(x)=2x^2$；对 $1/2<x\le1$，$F(x)=1-2(1-x)^2$。在区间左、右侧分别延拓为 $0,1$。所以通过概率为\[P(X\ge0.55)=2(1-0.55)^2=0.405.\]分位数位于右半区，解 $1-2(1-q)^2=0.875$ 得 $q=0.75$。另一个代数根 $1.25$ 不在支撑内，必须舍去。求分位数时先判断所在的分段，不能对整段使用同一公式。\end{solve}

\subsection*{训练题~\ref{q:M-PF-21} 从分布函数回到密度与期望}\label{a:M-PF-21}{\small 题目见第~\pageref{q:M-PF-21} 页；来源：MIT 模拟期末，题 21；二阶矩为编订延伸。}
\begin{solve}分布函数在两个端点连续，因此没有端点原子。在内部求导得 $f(x)=2-2x$。直接积分可得\[EX=\int_0^1x(2-2x)\,dx=\frac13,\]\[EX^2=\int_0^1x^2(2-2x)\,dx=\frac16.\]因此方差为 $1/6-1/9=1/18$。因无点质量，$P(X<0.4)=F(0.4)=0.64$。一般地，分布函数某点的跳跃大小才是该点概率，导数密度只捕捉绝对连续部分。\end{solve}

\subsection*{训练题~\ref{q:M-PF-29} 参数自助法与基本区间}\label{a:M-PF-29}{\small 题目见第~\pageref{q:M-PF-29} 页；来源：MIT 模拟期末，题 29。}
\begin{solve}均值为 $1/(1-p)$，矩估计也是极大似然估计：$\widehat p=1-1/\overline X$。以该拟合参数为真，从上述几何分布模拟独立自助样本，容量仍为 $n$；计算每次 $\widehat p^*=1-1/\overline X^*$，重复足够多次。令估计误差 $\Delta^*=\widehat p^*-\widehat p$ 的经验 $0.025,0.975$ 分位数为 $d_L,d_U$，基本区间为\[[\widehat p-d_U,\widehat p-d_L].\]它是对估计误差区间作反演，与直接取 $\widehat p^*$ 分位数的百分位区间不同。可与参数空间 $[0,1]$ 相交；若估计落在边界导致拟合分布退化，普通自助近似需特别谨慎。\end{solve}

\subsection*{训练题~\ref{q:M-PF-30} 极小尾概率与数值稳定性}\label{a:M-PF-30}{\small 题目见第~\pageref{q:M-PF-30} 页；来源：MIT 模拟期末，题 30（校订非零尾概率）。}
\begin{solve}标准误为 $2/10=0.2$，拒绝域为 $|\overline X-3|>1.959964(0.2)$。当均值为 $5$，观测 $z=10$，\[p=2[1-\Phi(10)]\approx1.52397\times10^{-23}.\]数值计算宜使用生存函数直接计算右尾，避免浮点数中 $1-\Phi(10)$ 的消减误差。在真实均值为 $4$ 时，接受区间相对该均值标准化的两端为 $-6.959964,-3.040036$，所以\[\pi(4)=1-\Phi(-3.040036)+\Phi(-6.959964)\approx0.99882.\]有限实数处正态尾概率严格为正；数值打印为零只表示精度不足。\end{solve}

\subsection*{训练题~\ref{q:M-PF-31} 正态测量的后验精度规划}\label{a:M-PF-31}{\small 题目见第~\pageref{q:M-PF-31} 页；来源：MIT 模拟期末，题 31。}
\begin{solve}先验精度为 $1/4$，单次观测精度为 $2$。因此后验方差为 $v=(1/4+2)^{-1}=4/9$，均值为\[m=\frac{80/4+85/0.5}{1/4+2}=\frac{760}{9}.\]故后验为 $N(760/9,4/9)$。在 $n$ 次测量后，后验方差为 $v_n=(1/4+2n)^{-1}=4/(1+8n)$。要求 $v_n<0.01$，即 $n>399/8=49.875$，所以最少总测量次数为 $50$。方差只依赖次数和误差方差，后验均值则还依赖实际测量值。\end{solve}

\subsection*{训练题~\ref{q:B17.1} 泊松分裂与参数推断}\label{a:B17.1}{\small 题目见第~\pageref{q:B17.1} 页；来源：编订补充。}
\begin{solve}给定 $N=n$，$X\sim\operatorname{Bin}(n,p)$，且 $Y=n-X$。对非负整数 $x,y$，\[P(X=x,Y=y)=e^{-\lambda}\frac{\lambda^{x+y}}{x!y!}p^x(1-p)^y.\]它可分解为率 $\lambda p$ 与 $\lambda(1-p)$ 的两个泊松质量函数，故 $X,Y$ 独立。多天数据的参数相关似然为\[e^{-m\lambda}\lambda^{T_X+T_Y}p^{T_X}(1-p)^{T_Y}.\]当总数 $T=T_X+T_Y>0$，极大似然为 $\widehat\lambda=T/m$、$\widehat p=T_X/T$。若 $T=0$，关于 $p$ 的似然恒定，数据不能识别分类比例；率参数的最大值位于零边界（若参数空间允许）。\end{solve}

\subsection*{训练题~\ref{q:B17.2} 有限总体抽样与修正因子}\label{a:B17.2}{\small 题目见第~\pageref{q:B17.2} 页；来源：编订补充。}
\begin{solve}每个大小为 $n$ 的子集等可能，所以\[P(X=k)=\frac{\binom Mk\binom{N-M}{n-k}}{\binom Nn}.\]令 $p=M/N$，单个抽样位置的指示变量均值为 $p$；不同位置的联合成功概率为 $M(M-1)/[N(N-1)]$，故协方差为 $-p(1-p)/(N-1)$。于是\[E\widehat p=p,\qquad\operatorname{Var}(\widehat p)=\frac{p(1-p)}n\frac{N-n}{N-1}.\]最后的因子是有限总体修正。抽样比例很小时近似有放回模型；当 $n=N$ 时普查没有抽样误差，方差恰好为零。\end{solve}

\subsection*{训练题~\ref{q:B17.3} 随机化回答的估计误差}\label{a:B17.3}{\small 题目见第~\pageref{q:B17.3} 页；来源：编订补充。}
\begin{solve}单人回答“是”的概率为 $q=p/2+1/2$，所以 $p=2q-1$。估计量 $\widehat p=2\widehat q-1$ 无偏，且\[\operatorname{Var}(\widehat p)=4q(1-q)/n=(1-p^2)/n.\]因为有限样本可能出现 $\widehat q<1/2$，该无偏估计可能为负。将它截断到 $[0,1]$ 可满足参数约束并减小针对该区间真值的平方误差，但通常引入偏差。该机制牺牲一部分统计效率以保护个体回答隐私；独立硬币机制与真实状态的关系必须明确。\end{solve}

\subsection*{训练题~\ref{q:B17.4} 尾和公式与收集全套}\label{a:B17.4}{\small 题目见第~\pageref{q:B17.4} 页；来源：编订补充。}
\begin{solve}已经有 $k$ 种类型时，下一次得到新类型的概率为 $(m-k)/m$。记此阶段等待次数为 $W_k$，包括成功一次。独立抽样与到达阶段时的重启性质保证各阶段等待时间相互独立，且 $W_k$ 为支撑 $1,2,\ldots$ 的几何变量。因此\[ET=\sum_{k=0}^{m-1}\frac{m}{m-k}=m\sum_{j=1}^m\frac1j.\]几何方差为 $(1-p)/p^2$，故\[\operatorname{Var}(T)=\sum_{j=1}^m\left(\frac{m^2}{j^2}-\frac mj\right).\]最后几种类型的成功概率很小，导致收集末期等待较久。只用“每种平均等 $m$ 次”相加会忽略状态变化。\end{solve}

\subsection*{训练题~\ref{q:B17.5} 随机和的复合泊松结构}\label{a:B17.5}{\small 题目见第~\pageref{q:B17.5} 页；来源：编订补充。}
\begin{solve}给定 $N$，条件均值与方差分别为 $N\mu,N\sigma^2$。全期望、全方差给出\[ES=\lambda\mu,\qquad\operatorname{Var}(S)=\lambda(\sigma^2+\mu^2).\]若 $\varphi_Z(t)=Ee^{itZ_1}$，则\[\varphi_S(t)=E[\varphi_Z(t)^N]=\exp\{\lambda[\varphi_Z(t)-1]\}.\]只知道 $\mu,\sigma^2$ 不能确定整个特征函数，也就不能精确确定一般尾概率。可以给出切比雪夫界或在额外大强度条件下作正态近似，但必须标明近似性质。\end{solve}

\subsection*{训练题~\ref{q:B17.6} 正态变量平方和与比值}\label{a:B17.6}{\small 题目见第~\pageref{q:B17.6} 页；来源：编订补充。}
\begin{solve}$U\sim\chi^2_2$、$V\sim\chi^2_1$ 且独立，所以 $(U/2)/(V/1)\sim F_{2,1}$，支撑为正半轴。将卡方视为同率 $1/2$ 的伽马变量，变换 $R=U/(U+V),T=U+V$ 的逆式为 $U=RT,V=(1-R)T$，雅可比为 $T$。联合密度可分解为贝塔核和伽马核，得到\[R\sim\operatorname{Beta}(1,1/2),\qquad T\sim\chi^2_3,\]并且独立。$R$ 的支撑是 $(0,1)$。二者虽都由平方和比值形成，但缩放比与份额比的分布不同。\end{solve}

\subsection*{训练题~\ref{q:B17.7} 指数竞争与获胜标记}\label{a:B17.7}{\small 题目见第~\pageref{q:B17.7} 页；来源：编订补充。}
\begin{solve}生存函数相乘得到 $P(T>t)=e^{-(\lambda+\mu)t}$，故 $T\sim\operatorname{Exp}(\lambda+\mu)$。对小区间 $dt$，联合密度标记为\[P(T\in dt,J=1)=\lambda e^{-\lambda t}e^{-\mu t}\,dt.\]将其写成 $[\lambda/(\lambda+\mu)]\,[(\lambda+\mu)e^{-(\lambda+\mu)t}]dt$，可见 $P(J=1)=\lambda/(\lambda+\mu)$ 且条件时间分布与无条件相同；$J=0$ 同理。因此 $T$ 与先发生哪一类事件独立。这一性质依赖指数分布的恒定风险率，一般寿命分布不一定成立。\end{solve}

\subsection*{训练题~\ref{q:B17.8} 条件期望作为最小平方预测}\label{a:B17.8}{\small 题目见第~\pageref{q:B17.8} 页；来源：编订补充。}
\begin{solve}记 $m(X)=E(Y\mid X)$。对二阶可积的 $g(X)$，展开平方并使用\[E[(Y-m(X))(m(X)-g(X))]=0,\]其理由是先条件于 $X$，残差的条件期望为零。因而\[E(Y-g(X))^2=E(Y-m(X))^2+E(m(X)-g(X))^2.\]最小值在 $g=m$ 几乎处处达到。联合正态且 $\operatorname{Var}(X)>0$ 时，\[m(X)=EY+\frac{\operatorname{Cov}(X,Y)}{\operatorname{Var}(X)}(X-EX).\]对一般联合分布，这个线性式只是最优线性预测，未必等于全部可测函数中的最优预测。\end{solve}

\subsection*{训练题~\ref{q:B17.9} 混合正态的方差分解}\label{a:B17.9}{\small 题目见第~\pageref{q:B17.9} 页；来源：编订补充。}
\begin{solve}令组标记为 $G$。全期望给出 $EX=w\mu_1+(1-w)\mu_2$。全方差分为组内波动和组间波动：\[\begin{aligned}\operatorname{Var}(X)&=w\sigma_1^2+(1-w)\sigma_2^2\\&\quad+w(1-w)(\mu_1-\mu_2)^2.\end{aligned}\]即使均值相同而方差不同，混合分布一般也非正态。例如共同均值为零时四阶矩为 $3[w\sigma_1^4+(1-w)\sigma_2^4]$，而相同总体方差的正态四阶矩是 $3[w\sigma_1^2+(1-w)\sigma_2^2]^2$；非退化权重及不同方差使前者严格更大。\end{solve}

\subsection*{训练题~\ref{q:B17.10} 德尔塔方法估计风险函数}\label{a:B17.10}{\small 题目见第~\pageref{q:B17.10} 页；来源：编订补充。}
\begin{solve}中心极限定理给出 $\sqrt n(\widehat p-p)\xrightarrow dN(0,p(1-p))$。导数为 $g^{\prime}(p)=1/[p(1-p)]$，故\[\sqrt n[g(\widehat p)-g(p)]\xrightarrow dN\left(0,\frac1{p(1-p)}\right).\]代入估计标准误为 $1/\sqrt{n\widehat p(1-\widehat p)}$，可在对数赔率尺度构造正态区间，再用逻辑函数 $e^u/(1+e^u)$ 映回比例区间。当样本全成功或全失败时公式未定义；固定内部真值下这类事件概率趋于零，但有限样本仍需使用适合边界的区间或明确的平滑方案。\end{solve}

\subsection*{训练题~\ref{q:B17.11} 威尔逊比例区间的反演}\label{a:B17.11}{\small 题目见第~\pageref{q:B17.11} 页；来源：编订补充。}
\begin{solve}两侧平方并移项得到\[(n+z^2)p^2-(2n\widehat p+z^2)p+n\widehat p^2\le0.\]求二次方程的两个根，区间中心为\[c=\frac{\widehat p+z^2/(2n)}{1+z^2/n},\]半径为\[h=\frac{z}{1+z^2/n}\sqrt{\frac{\widehat p(1-\widehat p)}n+\frac{z^2}{4n^2}}.\]因此区间为 $[c-h,c+h]$。它保留零假设方差中的未知 $p$ 再反演，因此即使 $\widehat p=0$ 或 $1$ 也不会退化成单点；沃尔德区间在这些情形半径为零。该推导仍基于正态近似，并非一般有限样本精确覆盖。\end{solve}

\subsection*{训练题~\ref{q:B17.12} 正态模型的充分统计量}\label{a:B17.12}{\small 题目见第~\pageref{q:B17.12} 页；来源：编订补充。}
\begin{solve}展开平方和后，联合密度可写为\[(2\pi\sigma^2)^{-n/2}\exp\left[-\frac{\sum_i x_i^2-2\mu\sum_i x_i+n\mu^2}{2\sigma^2}\right].\]样本仅通过两个给定总和进入参数相关部分，因子分解定理给出充分性。容量固定且 $n\ge2$ 时，由\[\sum_iX_i=n\overline X,\quad\sum_iX_i^2=(n-1)S^2+n\overline X^2\]可互相恢复两个表示。只保留均值会丢失离差信息，一般不能同时充分反映均值和方差两个未知参数。充分性说的是模型内参数信息，不意味着保留了所有可视化或模型诊断信息。\end{solve}

\subsection*{训练题~\ref{q:B17.13} 最小值模型的非对称估计}\label{a:B17.13}{\small 题目见第~\pageref{q:B17.13} 页；来源：编订补充。}
\begin{solve}似然中参数相关部分为 $e^{n\lambda\theta}\mathbf1_{\{\theta\le M\}}$，$M=\min_iX_i$，所以最大值在 $\widehat\theta=M$ 达到。因为 $M-\theta\sim\operatorname{Exp}(n\lambda)$，有 $EM=\theta+1/(n\lambda)$，无偏修正为 $M-1/(n\lambda)$。又\[P\left(0\le M-\theta\le\frac{-\log\alpha}{n\lambda}\right)=1-\alpha,\]所以区间为 $[M+\log\alpha/(n\lambda),M]$。它以最小值为上界，体现模型只允许观测位于参数右侧；对称正态区间在这里并非自然构造。\end{solve}

\subsection*{训练题~\ref{q:B17.14} 贝塔—二项预测的额外变异}\label{a:B17.14}{\small 题目见第~\pageref{q:B17.14} 页；来源：编订补充。}
\begin{solve}对参数积分得到\[P(K=k\mid D)=\binom mk\frac{B(a+k,b+m-k)}{B(a,b)}.\]令 $s=a+b$，则 $Ep=a/s$，$\operatorname{Var}(p)=ab/[s^2(s+1)]$。全期望给出 $EK=ma/s$；全方差给出\[\begin{aligned}\operatorname{Var}(K)&=mE[p(1-p)]+m^2\operatorname{Var}(p)\\&=\frac{mab(s+m)}{s^2(s+1)}.\end{aligned}\]固定均值参数的二项方差是 $mab/s^2$，预测方差多了因子 $(s+m)/(s+1)$。当 $m>1$ 时它大于一，因为共同未知参数使未来试验在边际预测下相关。\end{solve}

\subsection*{训练题~\ref{q:B17.15} 正态性、独立性与零相关}\label{a:B17.15}{\small 题目见第~\pageref{q:B17.15} 页；来源：编订补充。}
\begin{solve}正态分布对称，所以随机乘以独立符号不改变边缘分布，$X,Y$ 都为标准正态。又 $EXY=E(SZ^2)=ES\,EZ^2=0$，两个均值均为零，故协方差为零。另一方面 $Y^2=X^2$ 以概率一成立；若二者独立，则 $X^2,Y^2$ 也独立，但同一个非退化随机量不可能与自身独立。比如事件 $\{X^2\le1\}$ 与 $\{Y^2\le1\}$ 实际相同，其概率严格介于零与一。该向量集中在两条直线上，并非联合正态；边缘分别正态不等于联合正态。\end{solve}

\subsection*{训练题~\ref{q:B17.16} 支撑集破坏独立性}\label{a:B17.16}{\small 题目见第~\pageref{q:B17.16} 页；来源：编订补充。}
\begin{solve}积分支撑得 $f_X(x)=2(1-x)$、$f_Y(y)=2y$，各自在 $(0,1)$ 上。联合密度不等于边缘乘积，并且三角形外仍有两个边缘密度都正而联合密度为零的区域。计算得 $EX=1/3,EY=2/3$，\[EXY=\int_0^1\int_0^y2xy\,dx\,dy=\frac14.\]所以协方差为 $1/4-2/9=1/36>0$。检查独立性需要把支撑指示函数一起考虑，只看内部代数表达式能否拆成两个因子是不充分的。\end{solve}

\subsection*{训练题~\ref{q:B17.17} 蒙特卡洛积分与标准误}\label{a:B17.17}{\small 题目见第~\pageref{q:B17.17} 页；来源：编订补充。}
\begin{solve}$EU_i^2=1/3$，$EU_i^4=1/5$，所以 $\operatorname{Var}(U_i^2)=1/5-1/9=4/45$。因此 $\widehat I$ 无偏且方差为 $4/(45n)$。中心极限定理给出近似区间\[\widehat I\pm z_{1-\alpha/2}\sqrt{4/(45n)}.\]切比雪夫不等式给出有限样本保证 $P(|\widehat I-I|\ge\varepsilon)\le4/(45n\varepsilon^2)$。标准误按 $n^{-1/2}$ 下降，将标准误减半通常要将独立样本数增至四倍。数值积分真值在此已知是为了核查估计方法，实际复杂积分中通常需用样本方差估计未知方差。\end{solve}

\subsection*{训练题~\ref{q:B17.18} 反向变量的方差缩减}\label{a:B17.18}{\small 题目见第~\pageref{q:B17.18} 页；来源：编订补充。}
\begin{solve}由于 $U$ 与 $1-U$ 的边缘都均匀，$EA=1/3$。化简为 $A=U^2-U+1/2$，平方积分得\[EA^2=\int_0^1(U^2-U+1/2)^2\,dU=\frac7{60}.\]所以 $\operatorname{Var}(A)=7/60-1/9=1/180$。两个独立样本平方的平均方差为 $(4/45)/2=2/45$，是配对法的八倍。原因是两项负相关，反向取样抵消波动。反向变量法并非对所有函数都有改善，必须检查变换后函数的协方差方向。\end{solve}

\subsection*{训练题~\ref{q:B17.19} 切比雪夫界与正态近似的角色}\label{a:B17.19}{\small 题目见第~\pageref{q:B17.19} 页；来源：编订补充。}
\begin{solve}样本均值方差为 $0.04$。切比雪夫不等式给出\[P(|\overline X-\mu|\ge0.5)\le0.04/0.25=0.16,\]所以所求概率至少为 $0.84$。标准化阈值为 $0.5/0.2=2.5$，正态近似给出 $2\Phi(2.5)-1\approx0.98758$。如果总体本来正态，该数值就是精确概率；一般总体下只能视为大样本近似，容量 $100$ 不能单独保证误差很小。前者保守但不依赖分布形状，后者通常更精细但依赖额外近似依据。\end{solve}

\subsection*{训练题~\ref{q:B17.20} 学生化与斯卢茨基定理}\label{a:B17.20}{\small 题目见第~\pageref{q:B17.20} 页；来源：编订补充。}
\begin{solve}中心极限定理给出 $Z_n=\sqrt n(\overline X-\mu)/\sigma\xrightarrow dN(0,1)$。有限二阶矩保证样本一阶、二阶矩的强大数定律适用，故 $S^2\to\sigma^2$，于是 $\sigma/S\to1$ 依概率成立。斯卢茨基定理得到\[Z_n\frac\sigma S\xrightarrow dN(0,1).\]这里不要求 $Z_n$ 与 $S$ 独立。若总体正态，则有任意有限 $n$ 下的精确 $t_{n-1}$ 分布；一般总体中没有这一精确结论。小样本强偏态时，形式上使用 $t$ 分位数并不能自动保证所标注的置信水平。\end{solve}

\subsection*{训练题~\ref{q:B17.21} 三种收敛之间的辨析}\label{a:B17.21}{\small 题目见第~\pageref{q:B17.21} 页；来源：编订补充。}
\begin{solve}若 $E|X_n-X|^2\to0$，马尔可夫不等式给出\[P(|X_n-X|>\varepsilon)\le\frac{E|X_n-X|^2}{\varepsilon^2}\to0.\]取 $U\sim U(0,1)$，令 $X_n=n\mathbf1_{\{U\le1/n\}}$。对任意固定 $\varepsilon>0$ 和足够大的 $n$，$P(X_n>\varepsilon)=1/n\to0$，但 $EX_n=1$ 始终不变。罕见但越来越大的取值保留了期望。依概率收敛本身不足以交换极限和期望；常见充分条件包括支配收敛条件或一致可积性。\end{solve}

\subsection*{训练题~\ref{q:B17.22} 参数空间约束下的正态估计}\label{a:B17.22}{\small 题目见第~\pageref{q:B17.22} 页；来源：编订补充。}
\begin{solve}最大化似然等价于最小化 $n(\mu-\overline X)^2$，同时满足 $\mu\ge0$，因此\[\widehat\mu=\max(0,\overline X).\]对每个真实 $\mu\ge0$ 和每个观测值，将负数投影到零都不会增大到 $\mu$ 的距离，所以逐点有 $(\widehat\mu-\mu)^2\le(\overline X-\mu)^2$，进而均方误差不大于 $1/n$。但在 $\mu=0$ 时，$E\widehat\mu=1/\sqrt{2\pi n}>0$，故有偏。这个例子说明参数约束可以使有偏估计在允许的参数空间内具有更小的平方风险。\end{solve}

\subsection*{训练题~\ref{q:B17.23} 遗漏变量导致斜率变化}\label{a:B17.23}{\small 题目见第~\pageref{q:B17.23} 页；来源：编订补充。}
\begin{solve}总体线性预测的斜率为 $\operatorname{Cov}(X,Y)/\operatorname{Var}(X)$。利用协方差线性性，\[a=\beta_1+\beta_2\frac{\operatorname{Cov}(X,Z)}{\operatorname{Var}(X)}.\]当 $\beta_2=0$ 或 $\operatorname{Cov}(X,Z)=0$ 时额外项为零；也可能由其他特定退化关系使其为零。回归系数因此依赖模型中包含哪些变量，不能只根据两变量相关就赋予因果解释。这里推导的是总体最小平方投影系数，不是有限样本斜率的逐个数值等式。\end{solve}

\subsection*{训练题~\ref{q:B17.24} 加权最小二乘}\label{a:B17.24}{\small 题目见第~\pageref{q:B17.24} 页；来源：编订补充。}
\begin{solve}用白化矩阵 $\Sigma^{-1/2}$ 变换数据，得到协方差为单位阵的误差。最小化\[(Y-X\beta)^\mathsf T\Sigma^{-1}(Y-X\beta),\]正规方程给出\[\widehat\beta=(X^\mathsf T\Sigma^{-1}X)^{-1}X^\mathsf T\Sigma^{-1}Y.\]其均值为 $\beta$，协方差为 $(X^\mathsf T\Sigma^{-1}X)^{-1}$。独立异方差时，$\Sigma$ 为对角阵，目标函数是 $\sum_i(Y_i-x_i^\mathsf T\beta)^2/\sigma_i^2$，精度高的观测权重大。若这些方差需从数据估计，则以上公式的有限样本推断不能直接视为已知协方差情形。\end{solve}

\clearpage




